% Les types de tourisme et leurs problèmes — Allemand 1ère A — Cours signé OnBuch+
% Programme officiel MINESEC. Compiler avec : tectonic ALA05-les-types-de-tourisme-et-leurs-problemes.tex
\newcommand{\DOCMATIERE}{Allemand}
\newcommand{\DOCNIVEAU}{1ère A}
\newcommand{\DOCMODULE}{Chapitre 2 : Vie économique}
\newcommand{\DOCLECON}{ALA05}
\newcommand{\DOCTITRE}{Les types de tourisme et leurs problèmes}
% OnBuch+ — préambule « cours signés » ENRICHI (compilé avec tectonic / XeLaTeX)
% Système visuel « Pop » d'OnBuch+ : fond crème (jamais blanc), bordures encre
% épaisses, ombres « dures » décalées, titres Archivo Black en capitales.
% Version MATHÉMATIQUES : graphes (pgfplots/tikz), boîtes pédagogiques
% (définition, propriété, méthode, attention, le savais-tu, exercice résolu,
% exercice, TP, bilan), page de garde et signature OnBuch+ ancrée en bas.
%
% Macros attendues dans main.tex AVANT \input{preamble} :
%   \DOCMATIERE  \DOCNIVEAU  \DOCMODULE  \DOCLECON (numéro)  \DOCTITRE
\documentclass[11pt]{article}

\usepackage{fontspec}
\usepackage[ngerman,french]{babel} % français = langue principale ; l'allemand via \all{..} / textebox
\usepackage[a4paper,margin=2cm,top=3.4cm,bottom=2.8cm,headheight=1.4cm,footskip=1.3cm]{geometry}
\usepackage{amsmath,amssymb,amsfonts}
\usepackage{mathtools} % \xmapsto, \coloneqq... souvent utilisés par les modèles
\usepackage{cancel} % simplifications barrées (bilans de forces)
% \abs{z} (module, valeur absolue) et \norm{u} (norme), utilisés par les modèles
\providecommand{\abs}{}\let\abs\relax\DeclarePairedDelimiter{\abs}{\lvert}{\rvert}
\providecommand{\norm}{}\let\norm\relax\DeclarePairedDelimiter{\norm}{\lVert}{\rVert}
\usepackage[table]{xcolor} % [table] : fournit aussi \rowcolors (alternance de lignes)
\usepackage{pagecolor}
\usepackage{tikz}
\usetikzlibrary{arrows.meta,positioning,calc,shapes.geometric,shapes.misc,shapes.arrows,decorations.pathmorphing,decorations.pathreplacing,decorations.markings,patterns,shadows,mindmap,trees,fit,backgrounds,matrix,intersections,angles,quotes}
\usepackage{pgfplots}
\usepackage[version=4]{mhchem} % équations nucléaires \ce{^{235}_{92}U}
\usepackage[european,siunitx]{circuitikz} % schémas électriques
\pgfplotsset{compat=1.18}
\usepgfplotslibrary{fillbetween,groupplots} % aires entre courbes (calcul intégral)
\usepackage{enumitem}
\usepackage{fancyhdr}
% [most] charge toutes les bibliothèques tcolorbox (listings, minted, raster,
% documentation...) et gonfle inutilement la mémoire TeX sur un document long
% (~300 boîtes) : on ne charge que ce qui est réellement utilisé.
\usepackage[skins,breakable]{tcolorbox}
\usepackage{graphicx}
\usepackage{colortbl}
\usepackage{booktabs}
\usepackage{tabularx}
\usepackage{diagbox} % case d'en-tête diagonale des tableaux à double entrée
% Colonnes extensibles centrée / gauche / droite (C, L, R), souvent utilisées par les modèles
\newcolumntype{C}{>{\centering\arraybackslash}X}
\newcolumntype{L}{>{\raggedright\arraybackslash}X}
\newcolumntype{R}{>{\raggedleft\arraybackslash}X}
\usepackage{array}
\usepackage{multicol}
\usepackage{multirow}
\usepackage{siunitx}
% parse-numbers=false : accepte aussi \SI{2,5 \times 10^{-3}}{...}
\sisetup{locale=FR,output-decimal-marker={,},group-minimum-digits=4,parse-numbers=false}
\DeclareSIUnit\ohm{\ensuremath{\symup{\Omega}}} % U+2126 absent de la police
\DeclareSIUnit\division{div} % réglages d'oscilloscope (ms/div, V/div)
\DeclareSIUnit\tr{tr} % tours (vitesse de rotation, ex. \SI{1500}{\tr\per\minute})
\DeclareSIUnit\molar{M} % concentration molaire (ex. \SI{3}{\molar}, \SI{1}{\micro\molar})
\DeclareSIUnit\an{an} % année (le modèle confond parfois avec \year, un compteur TeX) — ex. \SI{5}{\kilo\meter\per\an}
\DeclareSIUnit\FCFA{FCFA} % franc CFA (ex. \SI{500}{\FCFA\per\kilo\gram})
\DeclareSIUnit\degreeBrix{\textdegree Brix} % teneur en sucre d'un sirop/jus (réfractométrie)
\DeclareSIUnit\month{mois} % le modèle utilise parfois \month, un compteur TeX, comme unité de durée
\usepackage{qrcode}
% \degree : commande fréquemment utilisée par les modèles pour un angle ou une
% température en dehors de \SI{}{} (non fournie par les paquets chargés).
\providecommand{\degree}{^{\circ}}
% \qed (fin de démonstration), utilisé par les modèles sans amsthm
\providecommand{\qed}{\hfill$\square$}
% \barre : double barre des valeurs interdites dans un tableau de variations (tkz-tab)
\providecommand{\barre}{\ensuremath{\|}}
% environnement proof (amsthm non chargé) utilisé par les modèles
\ifdefined\proof\else\newenvironment{proof}[1][Démonstration]{\par\noindent\textit{#1.}\ }{\qed\par}\fi
\usepackage{lastpage}
\usepackage{hyperref}
\hypersetup{hidelinks}

% ============================================================
% Palette « Pop » OnBuch+ — mêmes hex que la classe P (plus_ui.dart)
% ============================================================
\definecolor{poporange}{HTML}{F59321}
\definecolor{popdark}{HTML}{E07A0C}
\definecolor{popcream}{HTML}{FFF6E8}
\definecolor{popink}{HTML}{1C1714}
\definecolor{poppurple}{HTML}{7A5AE0}
\definecolor{popgreen}{HTML}{2E9E6A}
\definecolor{poppink}{HTML}{E0457B}
\definecolor{popblue}{HTML}{2F7DE1}
\definecolor{popgold}{HTML}{C98A12}
\definecolor{popmuted}{HTML}{8A7F76}
\definecolor{popline}{HTML}{EADFD2}
\definecolor{popgray}{HTML}{8A7F76}
\colorlet{popgrey}{popgray}
% Alias tolérants (le modèle nomme parfois les couleurs différemment)
\colorlet{popwhite}{white}
\colorlet{popblack}{popink}
\colorlet{popred}{poppink}
\colorlet{popyellow}{popgold}
% Teintes claires (fonds de tableaux/encadrés) — le modèle nomme parfois
% les couleurs avec un suffixe L (ex. popblueL) sans les avoir définies.
\colorlet{poporangeL}{poporange!20}
\colorlet{popdarkL}{popdark!20}
\colorlet{poppurpleL}{poppurple!20}
\colorlet{popgreenL}{popgreen!20}
\colorlet{poppinkL}{poppink!20}
\colorlet{popblueL}{popblue!20}
\colorlet{popgoldL}{popgold!20}
\colorlet{popredL}{popred!20}
\colorlet{popgrayL}{popgray!20}
% Teintes claires pour fonds de tableaux / zones
\colorlet{popblueL}{popblue!10!white}
\colorlet{poporangeL}{poporange!14!white}
\colorlet{popgreenL}{popgreen!12!white}
\colorlet{poppurpleL}{poppurple!12!white}
\colorlet{poppinkL}{poppink!12!white}
\colorlet{popgoldL}{popgold!14!white}

\pagecolor{popcream}

% ============================================================
% Typographie
% ============================================================
\defaultfontfeatures{Ligatures=TeX}
\setmainfont{PlusJakartaSans-Medium.ttf}[Path=../../../fonts/,BoldFont=PlusJakartaSans-Bold.ttf]
% Maths en sans-serif (Fira Math) avec chiffres et lettres droites de Plus
% Jakarta, pour que les formules se fondent dans le texte.
\usepackage[math-style=ISO,bold-style=ISO]{unicode-math}
\setmathfont{FiraMath-Regular.otf}[Path=../../../fonts/]
\setmathfont{PlusJakartaSans-Medium.ttf}[Path=../../../fonts/,range={up/num,up/{Latin,latin},\mathup}]
\usepackage{icomma} % virgule décimale sans espace en mode math
% Caractères mathématiques Unicode absents des polices de texte (Plus Jakarta,
% Archivo Black) : rendus par leur équivalent mathématique, en texte comme en
% formule — sinon ils s'affichent en carré vide (ex. « Arithmétique dans ℤ »).
\usepackage{newunicodechar}
\newunicodechar{ℝ}{\ensuremath{\mathbb{R}}}
\newunicodechar{ℤ}{\ensuremath{\mathbb{Z}}}
\newunicodechar{ℂ}{\ensuremath{\mathbb{C}}}
\newunicodechar{ℕ}{\ensuremath{\mathbb{N}}}
\newunicodechar{ℚ}{\ensuremath{\mathbb{Q}}}
\newunicodechar{𝔻}{\ensuremath{\mathbb{D}}}
\newunicodechar{α}{\ensuremath{\alpha}}
\newunicodechar{β}{\ensuremath{\beta}}
\newunicodechar{γ}{\ensuremath{\gamma}}
\newunicodechar{δ}{\ensuremath{\delta}}
\newunicodechar{ε}{\ensuremath{\varepsilon}}
\newunicodechar{θ}{\ensuremath{\theta}}
\newunicodechar{λ}{\ensuremath{\lambda}}
\newunicodechar{μ}{\ensuremath{\mu}}
\newunicodechar{σ}{\ensuremath{\sigma}}
\newunicodechar{τ}{\ensuremath{\tau}}
\newunicodechar{φ}{\ensuremath{\varphi}}
\newunicodechar{ψ}{\ensuremath{\psi}}
\newunicodechar{ω}{\ensuremath{\omega}}
\newunicodechar{Σ}{\ensuremath{\Sigma}}
% majuscules produites par \MakeUppercase dans les titres (α-aminés -> Α)
\newunicodechar{Α}{\ensuremath{\alpha}}
\newunicodechar{Β}{\ensuremath{\beta}}
\newunicodechar{Γ}{\ensuremath{\Gamma}}
\newunicodechar{Δ}{\ensuremath{\Delta}}
\newunicodechar{Ε}{\ensuremath{\varepsilon}}
\newunicodechar{Θ}{\ensuremath{\Theta}}
\newunicodechar{Λ}{\ensuremath{\Lambda}}
\newunicodechar{Μ}{\ensuremath{\mu}}
\newunicodechar{Τ}{\ensuremath{\tau}}
\newunicodechar{Φ}{\ensuremath{\Phi}}
\newunicodechar{Ψ}{\ensuremath{\Psi}}
\newunicodechar{Ω}{\ensuremath{\Omega}}
\newunicodechar{↦}{\ensuremath{\mapsto}}
\newunicodechar{∈}{\ensuremath{\in}}
\newunicodechar{∉}{\ensuremath{\notin}}
\newunicodechar{∧}{\ensuremath{\wedge}}
\newunicodechar{⇔}{\ensuremath{\Leftrightarrow}}
\newunicodechar{⟺}{\ensuremath{\Longleftrightarrow}}
\newunicodechar{⇒}{\ensuremath{\Rightarrow}}
\newunicodechar{⟹}{\ensuremath{\Longrightarrow}}
\newunicodechar{∘}{\ensuremath{\circ}}
\newunicodechar{∪}{\ensuremath{\cup}}
\newunicodechar{∩}{\ensuremath{\cap}}
\newunicodechar{≡}{\ensuremath{\equiv}}
\newunicodechar{⊂}{\ensuremath{\subset}}
\newunicodechar{⊥}{\ensuremath{\perp}}
\newunicodechar{∃}{\ensuremath{\exists}}
\newunicodechar{∀}{\ensuremath{\forall}}
\newunicodechar{∛}{\ensuremath{\sqrt[3]{\,}}}
\newunicodechar{∜}{\ensuremath{\sqrt[4]{\,}}}
\newunicodechar{⃗}{\ensuremath{{}^{\to}}}
\newunicodechar{ⁿ}{\ifmmode^{n}\else\textsuperscript{\ensuremath{n}}\fi}
\newunicodechar{ˣ}{\ifmmode^{x}\else\textsuperscript{\ensuremath{x}}\fi}
\newunicodechar{ᵇ}{\ifmmode^{b}\else\textsuperscript{\ensuremath{b}}\fi}
\newunicodechar{ʳ}{\ifmmode^{r}\else\textsuperscript{\ensuremath{r}}\fi}
\newunicodechar{ᵘ}{\ifmmode^{u}\else\textsuperscript{\ensuremath{u}}\fi}
\newunicodechar{ʸ}{\ifmmode^{y}\else\textsuperscript{\ensuremath{y}}\fi}
\newunicodechar{ᵃ}{\ifmmode^{a}\else\textsuperscript{\ensuremath{a}}\fi}
\newunicodechar{ᵉ}{\ifmmode^{e}\else\textsuperscript{\ensuremath{e}}\fi}
\newunicodechar{ᵏ}{\ifmmode^{k}\else\textsuperscript{\ensuremath{k}}\fi}
\newunicodechar{ᵖ}{\ifmmode^{p}\else\textsuperscript{\ensuremath{p}}\fi}
\newunicodechar{ᴺ}{\ifmmode^{N}\else\textsuperscript{\ensuremath{N}}\fi}
\newunicodechar{ᶜ}{\ifmmode^{c}\else\textsuperscript{\ensuremath{c}}\fi}
\newunicodechar{ʰ}{\ifmmode^{h}\else\textsuperscript{\ensuremath{h}}\fi}
\newunicodechar{ᵠ}{\ifmmode^{q}\else\textsuperscript{\ensuremath{q}}\fi}
\newunicodechar{ₙ}{\ifmmode_{n}\else\textsubscript{\ensuremath{n}}\fi}
\newunicodechar{ₐ}{\ifmmode_{a}\else\textsubscript{\ensuremath{a}}\fi}
\newunicodechar{ᵢ}{\ifmmode_{i}\else\textsubscript{\ensuremath{i}}\fi}
\newunicodechar{ᵣ}{\ifmmode_{r}\else\textsubscript{\ensuremath{r}}\fi}
\newunicodechar{ₖ}{\ifmmode_{k}\else\textsubscript{\ensuremath{k}}\fi}
\newunicodechar{ᵦ}{\ifmmode_{\beta}\else\textsubscript{\ensuremath{\beta}}\fi}
\newunicodechar{ₓ}{\ifmmode_{x}\else\textsubscript{\ensuremath{x}}\fi}
\newfontfamily\popdisplay{ArchivoBlack-Regular.ttf}[Path=../../../fonts/]
\newfontfamily\poplogo{SpaceGrotesk-Bold.ttf}[Path=../../../fonts/]
\newfontfamily\popsemi{PlusJakartaSans-SemiBold.ttf}[Path=../../../fonts/]
\newfontfamily\popxbold{PlusJakartaSans-ExtraBold.ttf}[Path=../../../fonts/]
\newfontfamily\popxboldls{PlusJakartaSans-ExtraBold.ttf}[Path=../../../fonts/,LetterSpace=3.0]

\setlist[enumerate]{leftmargin=1.7em,itemsep=5pt,topsep=4pt}
\setlist[itemize]{leftmargin=1.7em,itemsep=4pt,topsep=4pt,label={\color{poporange}\textbullet}}
\setlength{\parindent}{0pt}
\setlength{\parskip}{5pt}
\renewcommand{\arraystretch}{1.25}

% pgfplots : style Pop par défaut
\pgfplotsset{
  every axis/.append style={axis line style={popink,line width=1pt},tick style={popink},
    grid style={popline,line width=0.5pt},label style={font=\small},tick label style={font=\footnotesize}},
  popcurve/.style={poporange,line width=1.8pt,no markers,smooth},
}
% Même style disponible dans un \draw TikZ simple (hors pgfplots)
\tikzset{popcurve/.style={poporange,line width=1.8pt}}
\tikzset{popgrid/.style={popline,very thin}}
\colorlet{popcurve}{poporange} % le nom du style est parfois utilisé comme couleur

% ============================================================
% Pill / squiggle
% ============================================================
\newcommand{\poppill}[3]{%
  \tikz[baseline=-0.55ex]\node[draw=popink,line width=1.2pt,rounded corners=2.6pt,%
    fill=#2,inner xsep=6.5pt,inner ysep=3pt]{%
    \popxboldls\fontsize{7}{7}\selectfont\color{#3}\MakeUppercase{#1}};%
}
\newcommand{\popsquiggle}[1]{%
  \tikz[baseline=-0.4ex]{
    \draw[#1,line width=2.6pt,line cap=round]
      plot [smooth] coordinates {(0,0.09) (0.34,0.01) (0.68,0.21) (1.02,0.01) (1.36,0.10)};
  }%
}
% Mot-clé surligné (marqueur orange sous le texte)
\newcommand{\cle}[1]{{\popxbold\color{popdark}#1}}

% ============================================================
% En-tête / pied
% ============================================================
\pagestyle{fancy}
\fancyhf{}
\renewcommand{\headrulewidth}{0pt}
\renewcommand{\footrulewidth}{0pt}
\lhead{\raisebox{-0.15ex}{\poplogo\fontsize{13}{13}\selectfont\color{popink}OnBuch\color{poporange}+}}
\rhead{\poppill{\DOCMATIERE\ \textbullet\ \DOCNIVEAU\ \textbullet\ Leçon \DOCLECON}{white}{popink}}
\lfoot{\popxbold\fontsize{7.6}{7.6}\selectfont\color{popmuted}\MakeUppercase{Cours signé OnBuch+}\ \textbullet\ {\popsemi\DOCTITRE}}
\rfoot{\tikz[baseline=-0.6ex]\node[rounded corners=8.5pt,fill=popink,minimum height=17pt,minimum width=17pt,inner xsep=5pt,inner sep=0]{\popxbold\fontsize{8}{8}\selectfont\color{popcream}\thepage\,/\,\pageref*{LastPage}};}

% ============================================================
% Titres
% ============================================================
\newcommand{\chaptitre}[2]{%
  \begin{tcolorbox}[enhanced,colback=poporange,colframe=popink,boxrule=2.6pt,arc=11pt,
      drop shadow=popink,left=15pt,right=15pt,top=12pt,bottom=11pt,before skip=3pt,after skip=18pt]
    \poppill{#2}{popink}{popcream}\par\vspace{9pt}
    {\raggedright\hyphenpenalty=10000\exhyphenpenalty=10000\popdisplay\fontsize{22}{24}\selectfont\color{popcream}\MakeUppercase{#1}\par}
  \end{tcolorbox}
}
% Section numérotée : pastille orange avec le numéro + titre Archivo
\newcounter{popsec}
\newcommand{\coursec}[1]{%
  \stepcounter{popsec}\setcounter{popsub}{0}%
  \par\vspace{16pt}\noindent
  \begin{minipage}{\linewidth}
  \tikz[baseline=-0.7ex]\node[draw=popink,line width=1.6pt,fill=poporange,rounded corners=4pt,minimum size=22pt,inner sep=0]{\popdisplay\fontsize{12}{12}\selectfont\color{popcream}\thepopsec};%
  \hspace{8pt}{\popdisplay\fontsize{15}{17}\selectfont\color{popink}\MakeUppercase{#1}}\par
  \vspace{4pt}\noindent{\color{poporange}\rule{36pt}{3.6pt}}
  \end{minipage}\par\nopagebreak\vspace{8pt}
}
\newcounter{popsub}
\newcommand{\courssub}[1]{%
  \stepcounter{popsub}%
  \par\vspace{9pt}\noindent{\popxbold\fontsize{11.5}{13}\selectfont\color{popink}{\color{poporange}\thepopsec.\thepopsub}\hspace{6pt}#1}\par\nopagebreak\vspace{4pt}
}

% ============================================================
% Boîtes pédagogiques — même recette : fond blanc, bordure épaisse colorée,
% ombre dure encre, badge plein. Titre optionnel [#1].
% ============================================================
\tcbset{popbase/.style={breakable,enhanced,colback=white,boxrule=2.3pt,arc=9pt,
  shadow={3pt}{-3pt}{0pt}{popink},left=14pt,right=14pt,top=11pt,bottom=11pt,before skip=11pt,after skip=15pt,
  fontupper=\popsemi\fontsize{10.6}{14.5}\selectfont\color{popink},
  fontlower=\popsemi\fontsize{10.6}{14.5}\selectfont\color{popink}}}
% Coche dessinée (le glyphe ✓ n'existe pas dans les polices du document)
\newcommand{\popcheck}[1]{\tikz[baseline=-0.5ex]\draw[#1,line width=1.6pt,line cap=round,line join=round](-0.13,0.0)--(-0.04,-0.09)--(0.13,0.1);}
\providecommand{\coche}{\popcheck{popgreen}}
\providecommand{\resultat}[1]{\par\smallskip\noindent\fcolorbox{popgreen}{popgreenL}{\parbox{\dimexpr\linewidth-2\fboxsep-2\fboxrule\relax}{\textbf{Résultat.} #1}}\par\smallskip}
\newcommand{\popboxhead}[3]{\poppill{#1}{#2}{white}\ifx\relax#3\relax\else\hspace{6pt}{\popxbold\fontsize{10.6}{12}\selectfont\color{popink}#3}\fi\par\vspace{7pt}}

\newtcolorbox{aretenir}[1][]{popbase,colframe=popgreen,before upper={\popboxhead{À retenir}{popgreen}{#1}}}
% Texte en allemand (document de lecture, dialogue, extrait) : cadre
% d'étude. Un titre optionnel (source, auteur) s'affiche dans l'en-tête.
\newtcolorbox{textebox}[1][]{popbase,colframe=popblue,before upper={\popboxhead{Text}{popblue}{#1}\begin{otherlanguage}{ngerman}},after upper={\end{otherlanguage}}}
% Marques ✓ / ✗ (forme correcte / fautive) souvent utilisées par les modèles
\providecommand{\cross}{\textcolor{poppink}{\ensuremath{\times}}}
\providecommand{\xmark}{\cross}\providecommand{\crossmark}{\cross}\providecommand{\cmark}{\ensuremath{\checkmark}}
% Ligne de réponse / justification à compléter (inventée par les modèles)
\providecommand{\justification}[1]{\par\noindent\textit{Justification~:} \dotfill\par}
% \textipa (package tipa non chargé) : la police ne couvre pas l'API -> simple passage
\providecommand{\textipa}[1]{#1}
% Soulignement (ulem non chargé) : \uline, \ul
\providecommand{\uline}[1]{\underline{#1}}\providecommand{\ul}[1]{\underline{#1}}
% variantes ASCII d'ordinaux écrites en macro
\providecommand{\iere}{\textsuperscript{re}}\providecommand{\ieme}{\textsuperscript{e}}\providecommand{\ere}{\textsuperscript{re}}\providecommand{\eme}{\textsuperscript{e}}\providecommand{\er}{\textsuperscript{er}}\providecommand{\ie}{\textsuperscript{e}}
% Ordinaux français écrits en macro par les modèles : 1\ière, 3\ième
\newcommand{\ière}{\textsuperscript{re}}\newcommand{\ième}{\textsuperscript{e}}
% \exemple (inventée par les modèles) : étiquette « Exemple : »
\providecommand{\exemple}{\textbf{Exemple~:}\ }
% \square utilisable aussi hors mode math (cases à cocher)
\AtBeginDocument{\let\squareMath\square\renewcommand{\square}{\ensuremath{\squareMath}}}
% Mot ou phrase en allemand à l'intérieur d'un paragraphe français
\newcommand{\all}[1]{\ifmmode\text{\textit{#1}}\else\begingroup\selectlanguage{ngerman}\itshape #1\endgroup\fi}
\let\esp\all % alias toléré
\newtcolorbox{exemplebox}[1][]{popbase,colframe=poporange,before upper={\popboxhead{Exemple}{poporange}{#1}}}
\newtcolorbox{definition}[1][]{popbase,colframe=popblue,before upper={\popboxhead{Définition}{popblue}{#1}}}
\newtcolorbox{propriete}[1][]{popbase,colframe=poppurple,before upper={\popboxhead{Propriété}{poppurple}{#1}}}
\newtcolorbox{methode}[1][]{popbase,colframe=popgold,before upper={\popboxhead{Méthode}{popgold}{#1}}}
\newtcolorbox{attention}[1][]{popbase,colframe=poppink,before upper={\popboxhead{Attention}{poppink}{#1}}}
\newtcolorbox{experience}[1][]{popbase,colframe=popblue,colback=popblueL,before upper={\popboxhead{Expérience}{popblue}{#1}}}
% « Le savais-tu ? » : carte violette pleine (contexte, culture, Cameroun)
\newtcolorbox{savaistu}[1][]{popbase,colback=poppurple,colframe=popink,
  fontupper=\popsemi\fontsize{10.4}{14.2}\selectfont\color{white},
  before upper={\poppill{Le savais-tu\,?}{white}{poppurple}\ifx\relax#1\relax\else\hspace{6pt}{\popxbold\color{white}#1}\fi\par\vspace{7pt}}}
% Exercice résolu : énoncé en haut, \tcblower puis la solution sur fond crème
\newcounter{popexo}\newcounter{popexr}
\newtcolorbox{exoresolu}[1][]{popbase,colframe=poporange,colbacklower=popcream,
  segmentation style={popink,line width=1.2pt,dashed},
  before upper={\stepcounter{popexr}\popboxhead{Exercice résolu \thepopexr}{poporange}{#1}},
  before lower={\poppill{Solution}{popink}{popcream}\par\vspace{6pt}}}
% Exercice (non résolu en place) — la correction est dans la partie Corrigés
\newtcolorbox{exercice}[1][]{popbase,colframe=popink,
  before upper={\stepcounter{popexo}\popboxhead{Exercice \thepopexo}{popink}{#1}}}
% Corrigé
\newtcolorbox{corrige}[1][]{popbase,colframe=popgreen,colback=popgreenL,
  before upper={\popboxhead{Corrigé}{popgreen}{#1}}}
% Objectifs / prérequis / bilan
\newtcolorbox{objectifs}{popbase,colframe=popink,colback=poporangeL,
  before upper={\popboxhead{Objectifs de la leçon}{popink}{}}}
\newtcolorbox{prerequis}{popbase,colframe=popink,colback=popblueL,
  before upper={\popboxhead{Prérequis}{popblue}{}}}
\newtcolorbox{bilan}{popbase,colframe=popink,colback=white,boxrule=2.8pt,
  before upper={\popboxhead{Fiche bilan}{poporange}{L'essentiel à connaître pour l'examen}}}
% Figure encadrée avec légende
\newtcolorbox{popfigure}[1][]{enhanced,breakable=false,colback=white,colframe=popline,boxrule=1.4pt,arc=8pt,
  left=10pt,right=10pt,top=10pt,bottom=8pt,before skip=10pt,after skip=12pt,halign=center,
  lower separated=false,halign lower=center,
  fontlower=\popsemi\fontsize{9}{11.5}\selectfont\color{popmuted},
  #1}
\newcommand{\legende}[1]{\par\vspace{6pt}{\popsemi\fontsize{9}{11.5}\selectfont\color{popmuted}#1}}

% ============================================================
% Signature OnBuch+
% ============================================================
\newcommand{\poplogoinline}[1]{{\poplogo\fontsize{#1}{#1}\selectfont\color{popink}OnBuch\color{poporange}+}}
\newcommand{\signatureonbuch}{%
  \begin{tcolorbox}[enhanced,colback=popink,colframe=popink,boxrule=2.2pt,arc=10pt,
      drop shadow=poporange,left=14pt,right=14pt,top=10pt,bottom=10pt,before skip=0pt,after skip=0pt]
    \tikz[baseline=-0.7ex]\node[circle,fill=popgreen,draw=popcream,line width=1.3pt,minimum size=22pt,inner sep=0]{\popcheck{white}};%
    \hspace{9pt}\begin{minipage}[c]{0.62\linewidth}
      {\popxbold\fontsize{12}{13}\selectfont\color{popcream}Signé\ \textbullet\ {\poplogo OnBuch\color{poporange}+}}\par\vspace{2pt}
      {\raggedright\popsemi\fontsize{8}{10.5}\selectfont\color{popline}\MakeUppercase{Équipe pédagogique OnBuch+}\\\MakeUppercase{Conforme au programme officiel MINESEC}\par}
    \end{minipage}\hfill
    {\poplogo\fontsize{22}{22}\selectfont\color{popcream}OnBuch\color{poporange}+}
  \end{tcolorbox}%
}

% ============================================================
% Page de garde
% #1 titre · #2 matière · #3 niveau · #4 module · #5 n° leçon · #6 durée ·
% #7 description / objectifs (texte ou itemize)
% ============================================================
\newcommand{\pagegarde}[7]{%
  \thispagestyle{empty}
  \newgeometry{a4paper,margin=1.7cm,top=1.6cm,bottom=1.4cm}%
  \begin{tikzpicture}[remember picture,overlay]
    \fill[poporange!18] ([xshift=-3.2cm,yshift=-3.6cm]current page.north east) circle (4.6cm);
    \fill[poppurple!14] ([xshift=2.4cm,yshift=7.6cm]current page.south west) circle (3.8cm);
  \end{tikzpicture}%
  {\poplogo\fontsize{30}{30}\selectfont\color{popink}OnBuch\color{poporange}+}\hfill
  \raisebox{0.6ex}{\poppill{Cours signé \textbullet\ Premium}{popink}{popcream}}\par
  \vspace{1pt}\popsquiggle{poppurple}\par
  \vspace{8pt}
  \begin{tcolorbox}[enhanced,colback=poporange,colframe=popink,boxrule=2.8pt,arc=13pt,
      drop shadow=popink,left=18pt,right=18pt,top=12pt,bottom=13pt,after skip=10pt]
    \poppill{#2 \textbullet\ #3}{popink}{popcream}\hspace{5pt}\poppill{#4}{white}{popink}\par\vspace{8pt}
    {\popdisplay\fontsize{14}{14}\selectfont\color{popink}LEÇON #5}\par\vspace{4pt}
    {\raggedright\hyphenpenalty=10000\exhyphenpenalty=10000\popdisplay\fontsize{25}{27}\selectfont\color{popcream}\MakeUppercase{#1}\par}
    \vspace{8pt}
    {\popxbold\fontsize{9.5}{9.5}\selectfont\color{popink}\MakeUppercase{Durée conseillée : #6}}
  \end{tcolorbox}
  \begin{tcolorbox}[enhanced,colback=white,colframe=popink,boxrule=2.2pt,arc=10pt,
      drop shadow=popink,left=16pt,right=16pt,top=10pt,bottom=10pt,after skip=8pt]
    \poppill{Dans cette leçon}{poppurple}{white}\par\vspace{5pt}
    \popsemi\fontsize{9.3}{12.6}\selectfont\color{popink}#7
  \end{tcolorbox}
  \noindent\begin{minipage}[t]{0.487\linewidth}
    \begin{tcolorbox}[enhanced,colback=popgreen,colframe=popink,boxrule=2.2pt,arc=10pt,
        drop shadow=popink,left=11pt,right=11pt,top=10pt,bottom=10pt,height=3.1cm,valign=center]
      \tcbox[colback=white,colframe=popink,boxrule=1.3pt,arc=5pt,boxsep=0pt,nobeforeafter,
        left=3pt,right=3pt,top=3pt,bottom=3pt,valign=center]{\qrcode[height=1.9cm]{https://play.google.com/store/apps/details?id=com.luvvix.onbuch&hl=fr}}%
      \hspace{8pt}\begin{minipage}[c]{0.5\linewidth}
        {\popdisplay\fontsize{11}{12.5}\selectfont\color{white}\MakeUppercase{Télécharge OnBuch}}\par\vspace{4pt}
        {\raggedright\popsemi\fontsize{8.4}{11}\selectfont\color{white}Gratuit \textbullet\ Google Play\\Tuteur IA, annales, concours, résultats\par}
      \end{minipage}
    \end{tcolorbox}
  \end{minipage}\hfill
  \begin{minipage}[t]{0.487\linewidth}
    \begin{tcolorbox}[enhanced,colback=poppurple,colframe=popink,boxrule=2.2pt,arc=10pt,
        drop shadow=popink,left=11pt,right=11pt,top=10pt,bottom=10pt,height=3.1cm,valign=center]
      \tikz[baseline=-0.7ex]\node[circle,fill=white,draw=popink,line width=1.3pt,minimum size=1.6cm,inner sep=0]{\popdisplay\fontsize{24}{24}\selectfont\color{poppurple}+};%
      \hspace{8pt}\begin{minipage}[c]{0.6\linewidth}
        {\popdisplay\fontsize{11}{12.5}\selectfont\color{white}\MakeUppercase{Passe à OnBuch+}}\par\vspace{4pt}
        {\raggedright\popsemi\fontsize{8.4}{11}\selectfont\color{white}Tous les cours signés, Tuteur IA illimité, fiches PDF — onglet OnBuch+ de l'app\par}
      \end{minipage}
    \end{tcolorbox}
  \end{minipage}
  \vspace{8pt}
  \popcontenu
  \vfill
  \signatureonbuch
  \restoregeometry
  \newpage
}

% Rangée « Ce cours contient » (page de garde)
\newcommand{\poptile}[3]{% #1 couleur #2 gros texte #3 légende
  \begin{tcolorbox}[enhanced,colback=#1,colframe=popink,boxrule=2pt,arc=9pt,shadow={2.5pt}{-2.5pt}{0pt}{popink},
      left=6pt,right=6pt,top=9pt,bottom=9pt,halign=center,valign=center,height=1.75cm,width=\linewidth,nobeforeafter]
    {\popdisplay\fontsize{12.5}{13}\selectfont\color{white}\MakeUppercase{#2}}\par\vspace{3pt}
    {\centering\popsemi\fontsize{7.8}{9.5}\selectfont\color{white}#3\par}
  \end{tcolorbox}}
\newcommand{\popcontenu}{%
  \par\noindent{\popxboldls\fontsize{7.5}{7.5}\selectfont\color{popmuted}CE COURS SIGNÉ CONTIENT}\par\vspace{7pt}
  \noindent\begin{minipage}[t]{0.235\linewidth}\poptile{popblue}{Cours}{complet, illustré, pas à pas}\end{minipage}\hfill
  \begin{minipage}[t]{0.235\linewidth}\poptile{poporange}{Méthodes}{et exercices résolus}\end{minipage}\hfill
  \begin{minipage}[t]{0.235\linewidth}\poptile{poppink}{Exercices}{type examen + corrigés}\end{minipage}\hfill
  \begin{minipage}[t]{0.235\linewidth}\poptile{popgreen}{Fiche bilan}{l'essentiel à retenir}\end{minipage}\par}

% Sommaire
\newtcolorbox{sommaire}{enhanced,colback=white,colframe=popink,boxrule=2.2pt,arc=10pt,
  drop shadow=popink,left=18pt,right=18pt,top=13pt,bottom=13pt,before skip=6pt,after skip=20pt,
  before upper={\poppill{Sommaire}{poppurple}{white}\par\vspace{9pt}\popsemi\fontsize{10.6}{14.5}\selectfont\color{popink}}}

% Signature de fin de cours — poussée tout en bas de la dernière page
\newcommand{\findecours}{%
  \par\vfill\nopagebreak
  \begin{center}\popsquiggle{poporange}\end{center}\vspace{4pt}
  \signatureonbuch
}
\AtBeginDocument{\def\llbracket{\mathopen{[\mkern-3.5mu[}}\def\rrbracket{\mathclose{]\mkern-3.5mu]}}} % intervalles d'entiers (glyphe ⟦ absent de la police)
% Glyphes absents de Fira Math : redéfinis à partir de glyphes disponibles
\AtBeginDocument{%
  \def\setminus{\mathbin{\backslash}}\let\smallsetminus\setminus
  \def\vdots{\mathord{\vbox{\baselineskip3.2pt\lineskiplimit0pt\kern4pt\hbox{.}\hbox{.}\hbox{.}}}}%
  \def\ddots{\mathinner{\mkern1mu\raise6.5pt\hbox{.}\mkern2mu\raise3.6pt\hbox{.}\mkern2mu\raise0.7pt\hbox{.}\mkern1mu}}%
  \def\triangle{\mathord{\scalebox{0.9}{$\symup{\Delta}$}}}%
  \def\nabla{\mathord{\raisebox{\depth}{\scalebox{1}[-1]{$\symup{\Delta}$}}}}%
  \def\checkmark{\popcheck{popdark}}%
  \def\star{\mathbin{\ast}}\def\bigstar{\mathord{\ast}}%
  \def\diamond{\mathbin{\scalebox{0.6}{$\Diamond$}}}%
  \def\bigcup{\mathop{\mathchoice{\raisebox{-0.3em}{\scalebox{1.7}{$\cup$}}}{\scalebox{1.25}{$\cup$}}{\cup}{\cup}}\displaylimits}%
  \def\bigcap{\mathop{\mathchoice{\raisebox{-0.3em}{\scalebox{1.7}{$\cap$}}}{\scalebox{1.25}{$\cap$}}{\cap}{\cap}}\displaylimits}%
  \def\lesseqqgtr{\mathrel{\lessgtr}}\def\bigoplus{\mathop{\oplus}\displaylimits}%
}
\providecommand{\ssi}{si et seulement si} % abréviation fréquente
\AtBeginDocument{\providecommand{\wideparen}{\overparen}} % arc de cercle

%% FIGFIX-BEGIN v4 — correctifs globaux des figures (docs/onbuch-generation/tools/figfix)
\usepackage{adjustbox}\usepackage{etoolbox}
% toute figure TikZ/pgfplots plus large que la ligne est réduite à la largeur disponible
\BeforeBeginEnvironment{tikzpicture}{\begin{adjustbox}{max width=\linewidth}}
\AfterEndEnvironment{tikzpicture}{\end{adjustbox}}
% légendes pgfplots lisibles par-dessus les courbes
\pgfplotsset{every axis/.append style={legend style={fill=white,fill opacity=0.92,text opacity=1,draw=black!15,inner sep=3pt}}}
% étiquettes d'arêtes (syntaxe edge["..."]) avec fond blanc
\tikzset{every edge quotes/.append style={fill=white,inner sep=1.2pt}}
%% FIGFIX-END
\begin{document}

\pagegarde{Les types de tourisme et leurs problèmes}{Allemand}{1ère A}{Chapitre 2 : Vie économique}{ALA05}{2 h}{Leçon de lecture-commentaire sur les types de tourisme (balnéaire, culturel, écotourisme, tourisme de masse) et leurs problèmes économiques, sociaux et écologiques. Elle développe le lexique du voyage, les prépositions de lieu, les subordonnées en weil/dass et le parfait avec sein.
\vspace{4pt}
\begin{multicols}{2}\small
\begin{itemize}
\item À la fin de cette leçon, je sais comprendre globalement puis en détail un texte allemand sur les types de tourisme et leurs problèmes
\item À la fin de cette leçon, je sais employer le vocabulaire thématique du tourisme (der Tourismus, der Tourist, das Hotel, die Reise, der Strand, die Sehenswürdigkeit) avec articles et pluriels corrects
\item À la fin de cette leçon, je sais utiliser les prépositions de lieu à deux cas (accusatif = mouvement, datif = position) pour parler de voyages
\item À la fin de cette leçon, je sais construire des subordonnées de cause avec weil et d'opinion avec dass, verbe en fin de phrase
\item À la fin de cette leçon, je sais conjuguer au parfait les verbes de déplacement avec sein (reisen, fahren, fliegen, kommen)
\item À la fin de cette leçon, je sais répondre par écrit à des questions de compréhension en phrases complètes
\end{itemize}\end{multicols}}

\begin{sommaire}
\begin{multicols}{2}
\begin{enumerate}
\item Texte support : Tourismus – Segen oder Problem ?
\item Lexique thématique : der Tourismus
\item Grammaire 1 : les prépositions de lieu (accusatif/datif)
\item Grammaire 2 : weil et dass + parfait avec sein
\item Méthode du commentaire et commentaire modèle
\item Production guidée : Mein Traumurlaub
\item Activité d'intégration
\item Méthodes et exercices résolus
\item Exercices
\item Corrigés des exercices
\item Fiche bilan
\end{enumerate}
\end{multicols}
\end{sommaire}

\begin{objectifs}
\begin{itemize}
\item À la fin de cette leçon, je sais comprendre globalement puis en détail un texte allemand sur les types de tourisme et leurs problèmes
\item À la fin de cette leçon, je sais employer le vocabulaire thématique du tourisme (der Tourismus, der Tourist, das Hotel, die Reise, der Strand, die Sehenswürdigkeit) avec articles et pluriels corrects
\item À la fin de cette leçon, je sais utiliser les prépositions de lieu à deux cas (accusatif = mouvement, datif = position) pour parler de voyages
\item À la fin de cette leçon, je sais construire des subordonnées de cause avec weil et d'opinion avec dass, verbe en fin de phrase
\item À la fin de cette leçon, je sais conjuguer au parfait les verbes de déplacement avec sein (reisen, fahren, fliegen, kommen)
\item À la fin de cette leçon, je sais répondre par écrit à des questions de compréhension en phrases complètes
\item À la fin de cette leçon, je sais rédiger un court paragraphe guidé sur le tourisme au Cameroun
\item À la fin de cette leçon, je sais dégager l'idée principale d'un texte et présenter un mini-commentaire oral
\end{itemize}
\end{objectifs}

\begin{prerequis}
\begin{itemize}
\item Le parfait avec haben (Seconde) : ich habe gemacht
\item Les articles définis et indéfinis au nominatif, accusatif et datif
\item La place du verbe en position 2 dans la phrase principale
\item Le vocabulaire de base des voyages et des loisirs (Seconde)
\item Les subordonnées simples déjà vues (dass en début d'année)
\end{itemize}
\end{prerequis}

\newpage

\coursec{Texte support : Tourismus – Segen oder Problem ?}

\courssub{Situation d'accroche et problématique}

Chaque été, des milliers de touristes visitent la plage de \cle{Kribi}, les chutes de la Lobé ou le mont Cameroun, tandis que des Camerounais rêvent de découvrir les Alpes suisses ou les musées de Berlin. Mais le tourisme n'est pas qu'une affaire de plaisir : il rapporte de l'argent, crée des emplois, et pose aussi de graves problèmes — pollution des plages, prix qui montent, cultures locales transformées en spectacle. En Allemagne, on parle de \all{Massentourismus} à Majorque et de \all{sanfter Tourismus} (« tourisme doux ») dans les Alpes.

\textbf{Problématique :} Quels types de tourisme existe-t-il, et quels problèmes posent-ils ? C'est ce que nous allons découvrir à travers un texte allemand, puis nous apprendrons à en parler avec le vocabulaire et la grammaire du chapitre.

\courssub{Méthode : comment lire un texte allemand}

\begin{methode}[Lire et comprendre un texte en cinq étapes]
\begin{enumerate}
\item \textbf{Lire le titre} et deviner le thème. Ici : \all{Tourismus – Segen oder Problem ?} Le mot \all{Segen} (« bienfait ») annonce déjà une réflexion en deux parties : les avantages et les inconvénients.
\item \textbf{Première lecture globale}, en silence et \emph{sans dictionnaire} : chercher l'idée générale, pas chaque mot.
\item \textbf{Souligner les mots connus} et les mots de la famille du thème : \all{reisen, der Strand, das Hotel, die Touristen}\dots
\item \textbf{Deuxième lecture détaillée} avec le glossaire : comprendre phrase par phrase.
\item \textbf{Répondre aux questions en phrases complètes}, en reprenant les mots de la question : \all{Warum reisen viele Touristen ans Meer?} → \all{Viele Touristen reisen ans Meer, weil\dots}
\end{enumerate}
\end{methode}

\courssub{Le texte : lecture silencieuse puis lecture à voix haute}

\begin{textebox}[Tourismus – Segen oder Problem ?]
Der Tourismus ist heute sehr wichtig für viele Länder. Viele Touristen reisen im Sommer ans Meer, weil sie die Sonne und den Strand lieben. Man nennt das den Strandtourismus. Andere besuchen Städte wie Berlin oder Wien, weil sie sich für Kultur und Sehenswürdigkeiten interessieren : das ist der Kulturtourismus.

Der Tourismus bringt Geld und Arbeit. Die Hotels, die Restaurants und die Geschäfte verdienen viel, und die Einnahmen sind wichtig für das Land. In Kamerun, zum Beispiel in Kribi, finden viele Jugendliche Arbeit im Tourismus.

Aber der Tourismus hat auch Probleme. Beim Massentourismus reisen zu viele Menschen an denselben Ort : Hotels verschmutzen die Strände, der Müll liegt überall, und die Preise steigen für die Einheimischen. Auch die Natur leidet : die Umwelt ist in Gefahr.

Ich denke, dass der Ökotourismus eine gute Lösung ist, weil er die Natur schützt. Die Touristen schlafen in kleinen Häusern, essen lokale Produkte und respektieren die Kultur der Einheimischen. So ist der Tourismus ein Segen und kein Problem.
\end{textebox}

\textbf{Glossaire (à écrire en marge du texte) :}

\begin{center}
\begin{tabularx}{\linewidth}{|X|X|}\hline
\rowcolor{popblueL} \textbf{Deutsch} & \textbf{Français} \\ \hline
der Segen, - & le bienfait, la bénédiction \\ \hline
das Meer, -e & la mer \\ \hline
die Sehenswürdigkeit, -en & le site touristique, la curiosité \\ \hline
das Geschäft, -e & le magasin, le commerce \\ \hline
verdienen & gagner (de l'argent) \\ \hline
die Einnahmen (Pl.) & les recettes \\ \hline
der Müll & les ordures, les déchets \\ \hline
verschmutzen & polluer, salir \\ \hline
steigen (ist gestiegen) & monter, augmenter \\ \hline
der Einheimische, -n & l'habitant local, l'autochtone \\ \hline
die Umwelt & l'environnement \\ \hline
die Gefahr, -en & le danger \\ \hline
leiden (hat gelitten) & souffrir \\ \hline
die Lösung, -en & la solution \\ \hline
schützen & protéger \\ \hline
respektieren & respecter \\ \hline
\end{tabularx}
\end{center}

\begin{attention}[Lecture à voix haute : prononciation]
Prononce : \all{Segen} « sé-guenn », \all{Sehenswürdigkeit} « sé-henss-vur-dikh-kaït », \all{verschmutzen} « fèr-chmou-tsenn », \all{Ökotourismus} « é-ko-tou-riss-mouss », \all{Wien} « vine ». Rappel : le \all{w} allemand se prononce « v », le \all{v} se prononce « f », et \all{sch} se prononce « ch ».
\end{attention}

\courssub{Repérage global : thème, lieux, idée générale}

Après la première lecture, réponds mentalement à trois questions :

\begin{itemize}
\item \textbf{Thème :} le texte parle du \cle{tourisme}, de ses avantages et de ses problèmes.
\item \textbf{Lieux cités :} \all{Berlin} et \all{Wien} (villes culturelles), \all{Kribi} au Cameroun (tourisme balnéaire), la mer en général. On peut ajouter, pour la discussion, \all{die Alpen} et \all{die Schweiz}, destinations de rêve des touristes européens.
\item \textbf{Idée générale :} le tourisme est à la fois un \all{Segen} (argent, emplois) et un \all{Problem} (pollution, hausse des prix) ; l'écotourisme est présenté comme une solution.
\end{itemize}

\begin{attention}[Ne traduis pas mot à mot !]
Le français et l'allemand ne construisent pas les phrases de la même façon. Observe :
\begin{itemize}
\item \all{ans Meer reisen} = aller à la mer (\textbf{mouvement} → accusatif : \all{an das Meer}).
\item \all{am Meer sein} = être à la mer (\textbf{position} → datif : \all{an dem Meer}).
\end{itemize}
De même, \all{sich interessieren für} = s'intéresser \emph{à} : \all{Sie interessieren sich für Kultur.} « Ils s'intéressent à la culture. » Ne dis jamais \all{für Kultur interessieren} sans le \all{sich} ! Attention aussi au faux ami : \all{das Meer} = la mer, mais \all{die See} = la mer aussi (plus littéraire), alors que \all{der See} = le lac !
\end{attention}

\begin{exemplebox}[Deux phrases clés du texte, traduites]
\all{Viele Touristen reisen ans Meer, weil sie die Sonne lieben.} « Beaucoup de touristes vont à la mer parce qu'ils aiment le soleil. » Remarque : dans la subordonnée en \all{weil}, le verbe conjugué \all{lieben} va \textbf{à la fin}.\par\medskip
\all{Ich denke, dass der Ökotourismus eine gute Lösung ist.} « Je pense que l'écotourisme est une bonne solution. » Même règle avec \all{dass} : le verbe \all{ist} est rejeté en fin de phrase.
\end{exemplebox}

\courssub{Compréhension détaillée : les types de tourisme et leurs problèmes}

Relevons dans le texte les types de tourisme cités, avec un avantage et un problème pour chacun :

\begin{center}
\begin{tabularx}{\linewidth}{|l|X|X|}\hline
\rowcolor{popblueL} \textbf{Type de tourisme} & \textbf{Avantage (der Vorteil)} & \textbf{Problème (das Problem)} \\ \hline
der Strandtourismus & Sonne, Strand, Erholung (détente) & verschmutzte Strände \\ \hline
der Kulturtourismus & Kultur und Sehenswürdigkeiten & zu viele Besucher in den Städten \\ \hline
der Massentourismus & viele Einnahmen, Arbeit & Müll, steigende Preise für die Einheimischen \\ \hline
der Ökotourismus & schützt die Natur und die Kultur & weniger Geld, kleine Gruppen \\ \hline
\end{tabularx}
\end{center}

\begin{popfigure}
\begin{tikzpicture}[mindmap, grow cyclic,
  every node/.style={concept, align=center, font=\small},
  root concept/.append style={concept color=popblue, font=\bfseries},
  level 1/.append style={level distance=3.4cm, sibling angle=90},
  level 2/.append style={level distance=2.4cm, sibling angle=50, font=\scriptsize}]
\node[root concept]{Tourismus}
  child[concept color=poporange]{ node[align=center]{Strand-\\tourismus}
      child { node[concept color=poporange!40, align=center]{+ Sonne\\+ Erholung} }
      child { node[concept color=poporange!40, align=center]{-- schmutzige\\Strände} } }
  child[concept color=poppurple]{ node[align=center]{Kultur-\\tourismus}
      child { node[concept color=poppurple!40, align=center]{+ Kultur\\+ Museen} }
      child { node[concept color=poppurple!40, align=center]{-- zu viele\\Besucher} } }
  child[concept color=poppink]{ node[align=center]{Massen-\\tourismus}
      child { node[concept color=poppink!40, align=center]{+ Geld\\+ Arbeit} }
      child { node[concept color=poppink!40, align=center]{-- Müll\\-- hohe Preise} } }
  child[concept color=popgreen]{ node[align=center]{Öko-\\tourismus}
      child { node[concept color=popgreen!40, align=center]{+ schützt\\die Natur} }
      child { node[concept color=popgreen!40, align=center]{-- kleine\\Gruppen} } };
\end{tikzpicture}
\legende{Carte mentale : les quatre types de tourisme du texte, avec un avantage (+) et un problème (--) pour chacun.}
\end{popfigure}

\courssub{Questions de compréhension}

\begin{exercice}[Compréhension écrite]
\textbf{A. Vrai ou faux ? Justifie avec une phrase du texte.}
\begin{enumerate}
\item Alle Touristen reisen ans Meer.
\item Der Tourismus bringt Geld und Arbeit.
\item Die Preise steigen für die Touristen.
\item Der Ökotourismus schützt die Natur.
\end{enumerate}
\textbf{B. Réponds en allemand, en phrases complètes.}
\begin{enumerate}
\item Warum reisen viele Touristen ans Meer?
\item Welche Städte besuchen die Kulturtouristen?
\item Was ist ein Problem des Massentourismus?
\end{enumerate}
\textbf{C. Question personnelle :} \all{Möchtest du als Tourist nach Deutschland reisen? Warum?}
\end{exercice}

\begin{corrige}[Exercice : compréhension écrite]
\textbf{A.} 1. \textbf{Faux} : \all{Andere besuchen Städte wie Berlin oder Wien.} 2. \textbf{Vrai} : \all{Der Tourismus bringt Geld und Arbeit.} 3. \textbf{Faux} : \all{Die Preise steigen für die Einheimischen} (pour les habitants, pas pour les touristes). 4. \textbf{Vrai} : \all{Der Ökotourismus ist eine gute Lösung, weil er die Natur schützt.}\par
\textbf{B.} 1. \all{Viele Touristen reisen ans Meer, weil sie die Sonne und den Strand lieben.} 2. \all{Die Kulturtouristen besuchen Städte wie Berlin oder Wien.} 3. \all{Die Hotels verschmutzen die Strände, und die Preise steigen für die Einheimischen.}\par
\textbf{C.} Modèles de réponse : \all{Ja, ich möchte nach Deutschland reisen, weil ich Berlin und die Museen sehen möchte.} / \all{Nein, ich möchte lieber in Kamerun bleiben, weil Kribi sehr schön ist.}
\end{corrige}

\courssub{Reformulation : résumer le texte}

\begin{exoresolu}[Résumé en trois phrases avec weil et dass]
\textbf{Énoncé :} Résume le texte en trois phrases : une phrase avec \all{weil}, une phrase avec \all{dass}, une phrase libre.
\tcblower
\textbf{Solution détaillée :} On reprend l'idée générale (avantages + problèmes + solution) sans recopier le texte.\par
1. \all{Viele Menschen reisen gern, weil der Tourismus schöne Ferien bringt.} (subordonnée en \all{weil}, verbe \all{bringt} à la fin)\par
2. \all{Ich denke, dass der Massentourismus die Umwelt verschmutzt.} (subordonnée en \all{dass}, verbe \all{verschmutzt} à la fin)\par
3. \all{Der Ökotourismus ist eine gute Lösung für die Natur und die Einheimischen.}\par
\textbf{Méthode :} une idée par phrase ; vérifier la place du verbe conjugué à la fin après \all{weil} et \all{dass}.
\end{exoresolu}

\begin{experience}[Débat en classe : Segen oder Problem ?]
Par groupes de quatre : deux élèves défendent le tourisme (\all{Der Tourismus ist ein Segen, weil\dots}), deux élèves le critiquent (\all{Der Tourismus ist ein Problem, weil\dots}). Chacun doit produire au moins deux phrases avec \all{weil} ou \all{dass}, en utilisant le vocabulaire du texte : \all{die Einnahmen, die Arbeit, der Müll, die Umwelt, die Einheimischen}. Exemple de départ : \all{Ich denke, dass der Tourismus wichtig für Kamerun ist, weil er Arbeit bringt.}
\end{experience}

\begin{savaistu}[Les Allemands, champions du voyage]
L'Allemagne est l'un des premiers pays émetteurs de touristes en Europe : les Allemands voyagent énormément, surtout vers l'Espagne, l'Italie et l'Autriche. L'île de Majorque est tellement prisée des touristes allemands qu'on la surnomme en plaisantant le « 17. Land » — comme si elle était un dix-septième État fédéré allemand ! C'est là que le mot \all{Massentourismus} prend tout son sens. En réaction, le \all{sanfte Tourismus} (« tourisme doux ») se développe dans les Alpes : randonnée, petits hôtels, produits locaux — un modèle proche de l'écotourisme. Et au Cameroun ? Les chutes de la Lobé, uniques au monde parce qu'elles se jettent directement dans l'océan, sont un trésor pour un écotourisme à développer.
\end{savaistu}

\begin{aretenir}[Ce qu'il faut retenir de cette première étape]
\begin{itemize}
\item Le texte oppose le tourisme comme \all{Segen} (Geld, Arbeit) et comme \all{Problem} (Müll, verschmutzte Strände, steigende Preise).
\item Quatre types de tourisme : \all{der Strandtourismus, der Kulturtourismus, der Massentourismus, der Ökotourismus}.
\item Mouvement → \all{ans Meer reisen} ; position → \all{am Meer sein}.
\item Après \all{weil} et \all{dass}, le verbe conjugué va à la fin de la phrase.
\item Vocabulaire clé : \all{die Sehenswürdigkeit, die Einnahmen, der Müll, verschmutzen, die Umwelt, schützen, der Einheimische}.
\end{itemize}
\end{aretenir}

\coursec{Lexique thématique : der Tourismus}

Dans la section précédente, nous avons lu et commenté le texte \all{Tourismus – Segen oder Problem ?}. Nous y avons rencontré de nombreux mots nouveaux : \all{der Tourist}, \all{die Sehenswürdigkeit}, \all{die Umweltverschmutzung}… Pour bien comprendre un texte et surtout pour pouvoir en parler et écrire à notre tour, il faut maintenant \cle{organiser ce vocabulaire} et l'apprendre méthodiquement. C'est l'objectif de cette section : classer les mots par \cle{champs lexicaux}, apprendre chaque nom avec son \cle{article} et son \cle{pluriel}, découvrir les \cle{familles de mots} et mémoriser les \cle{expressions utiles} du thème.

\begin{methode}[Comment apprendre un nom allemand ?]
Un nom allemand ne s'apprend jamais seul. Pour chaque nom, retiens toujours trois informations :
\begin{enumerate}
\item \textbf{Le mot} : \all{Strand}
\item \textbf{Son article (son genre)} : \all{der Strand} — masculin
\item \textbf{Son pluriel} : \all{die Strände}
\end{enumerate}
On note donc dans son cahier : \all{der Strand (¨-e)}. Le signe ¨ indique que le pluriel prend un Umlaut (\all{a} devient \all{ä}). Apprendre un nom sans son article, c'est apprendre la moitié du mot : sans article, impossible de décliner correctement !
\end{methode}

\courssub{Champ lexical 1 : les acteurs du tourisme}

Qui voyage ? Qui accueille ? Qui guide ? Voici les personnes que l'on rencontre dans le monde du tourisme.

\begin{center}
\begin{tabularx}{\linewidth}{|X|X|X|}
\hline
\rowcolor{popblueL} \textbf{Deutsch} & \textbf{Français} & \textbf{Remarque} \\
\hline
der Tourist (-en) & le touriste & nom faible : voir l'encadré attention \\
\hline
die Touristin (-nen) & la touriste & féminin formé avec -in \\
\hline
der Reisende (-n) & le voyageur & nom faible (participe substantivé) : den Reisenden \\
\hline
der Gast (¨-e) & le client, l'invité & pluriel : die Gäste \\
\hline
der Reiseleiter (-) & le guide (accompagnateur) & mot composé : Reise + Leiter ; pluriel : die Reiseleiter \\
\hline
die Einheimischen (pl.) & les habitants du pays, les autochtones & n'existe qu'au pluriel ; adjectif substantivé \\
\hline
\end{tabularx}
\end{center}

\begin{exemplebox}[Les acteurs en contexte]
\all{Viele Touristen besuchen jedes Jahr die kamerunische Küste.} « Beaucoup de touristes visitent chaque année la côte camerounaise. »

\all{Der Reiseleiter zeigt den Gästen die Stadt Douala.} « Le guide montre la ville de Douala aux clients. »

\all{Die Einheimischen sind sehr freundlich zu den Touristen.} « Les habitants sont très aimables avec les touristes. »

\all{Eine Touristin aus Österreich fotografiert den Hafen.} « Une touriste d'Autriche photographie le port. »
\end{exemplebox}

\begin{attention}[der Tourist : un nom faible !]
\all{Der Tourist} appartient au groupe des \cle{noms faibles} (comme \all{der Student}, \all{der Junge}, \all{der Reisende}). Ces noms masculins prennent la terminaison \textbf{-en} à tous les cas \textbf{sauf au nominatif singulier} :

\begin{center}
\begin{tabular}{|l|l|}
\hline
Nominatif & der Tourist \\
\hline
Accusatif & den Tourist\textbf{en} \\
\hline
Datif & dem Tourist\textbf{en} \\
\hline
Génitif & des Tourist\textbf{en} \\
\hline
\end{tabular}
\end{center}

On écrit donc : \all{Ich helfe dem Touristen.} « J'aide le touriste. » (Datif) et \all{Wir sehen einen Touristen am Strand.} « Nous voyons un touriste à la plage. » (Accusatif).

Erreur typique du francophone : dire \all{mit einem Tourist} au lieu de \all{mit einem Touristen}. Même après \all{einem}, la terminaison \textbf{-en} est obligatoire !
\end{attention}

\courssub{Champ lexical 2 : les lieux et l'hébergement}

Où loge-t-on ? Que visite-t-on ? Voici les lieux essentiels du voyage.

\begin{center}
\begin{tabularx}{\linewidth}{|X|X|X|}
\hline
\rowcolor{popblueL} \textbf{Deutsch} & \textbf{Français} & \textbf{Remarque} \\
\hline
das Hotel (-s) & l'hôtel & pluriel régulier en -s : die Hotels \\
\hline
die Jugendherberge (-n) & l'auberge de jeunesse & Jugend (jeunesse) + Herberge (auberge) \\
\hline
der Campingplatz (¨-e) & le camping & Platz = place, terrain ; pluriel : die Campingplätze \\
\hline
der Strand (¨-e) & la plage & pluriel : die Strände \\
\hline
das Meer (-e) & la mer & attention : féminin en français, neutre en allemand ! Pluriel : die Meere \\
\hline
die Sehenswürdigkeit (-en) & le site touristique, la curiosité & voir la définition ci-dessous \\
\hline
das Museum (Museen) & le musée & pluriel irrégulier : die Museen \\
\hline
\end{tabularx}
\end{center}

\begin{definition}[die Sehenswürdigkeit (-en)]
Une \all{Sehenswürdigkeit} est littéralement « ce qui mérite d'être vu » : un site touristique, un monument, une curiosité à visiter. C'est un mot composé que l'on peut découper ainsi :

\all{sehen} (voir) + \all{würdig} (digne) + \all{-keit} (suffixe qui forme un nom féminin)

Donc : « dignité-d'être-vu » = ce qui vaut la peine d'être vu. Exemple : \all{Die Wasserfälle von Ekom sind eine berühmte Sehenswürdigkeit Kameruns.} « Les chutes d'Ekom sont un site célèbre du Cameroun. »
\end{definition}

\begin{exemplebox}[Les lieux en contexte]
\all{Wir übernachten in einem kleinen Hotel am Meer.} « Nous passons la nuit dans un petit hôtel au bord de la mer. »

\all{Die Jugendherberge ist billiger als das Hotel.} « L'auberge de jeunesse est moins chère que l'hôtel. »

\all{Am Strand von Kribi gibt es viele Palmen.} « Sur la plage de Kribi, il y a beaucoup de palmiers. »

\all{Das Nationalmuseum in Yaoundé ist jeden Tag geöffnet.} « Le musée national de Yaoundé est ouvert tous les jours. »
\end{exemplebox}

\courssub{Champ lexical 3 : les activités}

Que fait un touriste ? Il voyage, il visite, il se repose… Voici les noms et les verbes des activités.

\begin{center}
\begin{tabularx}{\linewidth}{|X|X|X|}
\hline
\rowcolor{popblueL} \textbf{Deutsch} & \textbf{Français} & \textbf{Remarque} \\
\hline
die Reise (-n) & le voyage & du verbe reisen \\
\hline
der Ausflug (¨-e) & l'excursion & pluriel : die Ausflüge \\
\hline
die Rundreise (-n) & le circuit (voyage organisé) & rund = circulaire \\
\hline
besuchen & visiter (des personnes, des lieux) & verbe régulier, auxiliaire haben \\
\hline
besichtigen & visiter (un monument, un site) & plus précis que besuchen, auxiliaire haben \\
\hline
wandern & randonner, marcher & verbe fort : ist gewandert (auxiliaire sein) \\
\hline
baden & se baigner & hat gebadet (auxiliaire haben) \\
\hline
fotografieren & photographier & hat fotografiert \\
\hline
sich erholen & se reposer, se détendre & verbe réfléchi : ich erhole mich, hat sich erholt \\
\hline
\end{tabularx}
\end{center}

\begin{attention}[besuchen ou besichtigen ?]
Les deux verbes se traduisent par « visiter », mais ils ne s'emploient pas dans les mêmes contextes :

\begin{itemize}
\item \all{besuchen} : visiter des \textbf{personnes} ou fréquenter un lieu : \all{Ich besuche meine Großeltern.} « Je rends visite à mes grands-parents. » \all{Er besucht die Schule.} « Il fréquente l'école. »
\item \all{besichtigen} : visiter un \textbf{monument, un site, une ville} en touriste : \all{Wir besichtigen das Schloss.} « Nous visitons le château. »
\end{itemize}

Pour un site touristique, le verbe naturel est donc \all{besichtigen} : \all{Die Touristen besichtigen die Sehenswürdigkeiten der Stadt.} « Les touristes visitent les sites de la ville. »
\end{attention}

\begin{exemplebox}[Expressions utiles à mémoriser]
\all{eine Reise machen} = faire un voyage : \all{Wir machen eine Reise nach Kribi.} « Nous faisons un voyage à Kribi. »

\all{in den Urlaub fahren} = partir en vacances : \all{Im Sommer fahren wir in den Urlaub.} « En été, nous partons en vacances. » (verbe fort : fährt, fuhr, ist gefahren)

\all{am Strand liegen} = être allongé sur la plage : \all{Die Touristen liegen am Strand und baden im Meer.} « Les touristes sont allongés sur la plage et se baignent dans la mer. »

\all{sich gut erholen} = bien se reposer : \all{Nach der Reise haben wir uns gut erholt.} « Après le voyage, nous nous sommes bien reposés. »
\end{exemplebox}

\courssub{Champ lexical 4 : les problèmes du tourisme}

Le texte de la section 1 le montrait : le tourisme apporte de l'argent, mais aussi des problèmes. Voici le vocabulaire pour en parler — avec un aspect positif à la fin.

\begin{center}
\begin{tabularx}{\linewidth}{|X|X|X|}
\hline
\rowcolor{popblueL} \textbf{Deutsch} & \textbf{Français} & \textbf{Remarque} \\
\hline
die Umweltverschmutzung (-en) & la pollution de l'environnement & Umwelt = environnement \\
\hline
die Preise steigen & les prix augmentent & verbe fort : steigt, stieg, ist gestiegen \\
\hline
der Müll & les déchets, les ordures & pas de pluriel courant \\
\hline
die Natur zerstören & détruire la nature & zerstören = détruire \\
\hline
die Arbeitsplätze schaffen & créer des emplois & aspect positif ! Arbeit = travail, Platz = place \\
\hline
\end{tabularx}
\end{center}

\begin{exemplebox}[Les problèmes en contexte]
\all{Der Massentourismus bringt oft Umweltverschmutzung.} « Le tourisme de masse apporte souvent la pollution de l'environnement. »

\all{In der Hochsaison steigen die Preise für Hotels.} « En haute saison, les prix des hôtels augmentent. »

\all{Viele Touristen lassen ihren Müll am Strand.} « Beaucoup de touristes laissent leurs déchets sur la plage. »

\all{Der Tourismus schafft viele Arbeitsplätze für die Einheimischen.} « Le tourisme crée beaucoup d'emplois pour les habitants. »
\end{exemplebox}

\begin{propriete}[Le genre des noms composés]
L'allemand adore les mots composés. Règle d'or : \textbf{un nom composé prend toujours le genre (l'article) de son dernier élément}.

\begin{itemize}
\item \all{das Meer} + \all{die Verschmutzung} $\rightarrow$ \all{die Meeresverschmutzung} (féminin, comme \all{Verschmutzung}, avec Fugen-\all{s})
\item \all{die Jugend} + \all{die Herberge} $\rightarrow$ \all{die Jugendherberge} (féminin)
\item \all{das Camping} + \all{der Platz} $\rightarrow$ \all{der Campingplatz} (masculin)
\item \all{die Reise} + \all{der Leiter} $\rightarrow$ \all{der Reiseleiter} (masculin)
\item \all{rund} (adj.) + \all{die Reise} $\rightarrow$ \all{die Rundreise} (féminin)
\end{itemize}

Quand tu rencontres un mot composé inconnu, découpe-le : le dernier élément te donne le genre et souvent le sens général. C'est une stratégie de compréhension très puissante !
\end{propriete}

\courssub{Les familles de mots : apprendre plus vite}

En allemand comme en français, les mots se regroupent en \cle{familles} autour d'une même racine. Connaître la famille d'un mot, c'est apprendre quatre ou cinq mots pour le prix d'un.

\begin{itemize}
\item \all{reisen} (voyager) $\rightarrow$ \all{die Reise} (le voyage), \all{der Reisende} (le voyageur), \all{das Reisebüro} (l'agence de voyages), \all{der Reiseleiter} (le guide)
\item \all{besuchen} (visiter) $\rightarrow$ \all{der Besuch} (la visite), \all{der Besucher} (le visiteur)
\item \all{sehen} (voir) $\rightarrow$ \all{die Sehenswürdigkeit} (le site à voir), \all{sehenswert} (qui vaut la peine d'être vu)
\end{itemize}

\begin{popfigure}
\begin{center}
\begin{tikzpicture}[
  root/.style={circle, draw=popblue, fill=popblueL, thick, align=center, minimum size=1.6cm, font=\bfseries},
  child1/.style={rectangle, rounded corners, draw=poporange, fill=poporangeL, thick, align=center},
  child2/.style={rectangle, rounded corners, draw=popgreen, fill=popgreenL, thick, align=center},
  child3/.style={rectangle, rounded corners, draw=poppurple, fill=poppurpleL, thick, align=center},
  child4/.style={rectangle, rounded corners, draw=poppink, fill=poppinkL, thick, align=center},
  link/.style={thick, popmuted}
]
\node[root, align=center] (R){reisen\\ voyager};
\node[child1, above left=1.4cm and 2.6cm of R, align=center] (A){die Reise\\ le voyage};
\node[child2, above right=1.4cm and 2.6cm of R, align=center] (B){der Reisende\\ le voyageur};
\node[child3, below left=1.4cm and 2.6cm of R, align=center] (C){das Reisebüro\\ l'agence de voyages};
\node[child4, below right=1.4cm and 2.6cm of R, align=center] (D){der Reiseleiter\\ le guide};
\draw[link] (R) -- (A);
\draw[link] (R) -- (B);
\draw[link] (R) -- (C);
\draw[link] (R) -- (D);
\end{tikzpicture}
\end{center}
\legende{Carte mentale : la famille du verbe \all{reisen}. À partir d'un seul verbe, on construit quatre noms utiles.}
\end{popfigure}

\begin{aretenir}[Tableau récapitulatif du lexique essentiel]
\begin{center}
\begin{tabularx}{\linewidth}{|l|X|X|}
\hline
\rowcolor{popblueL} \textbf{Thème} & \textbf{Deutsch} & \textbf{Français} \\
\hline
Acteurs & der Tourist (-en), der Reisende (-n), der Gast (¨-e), der Reiseleiter (-), die Einheimischen & le touriste, le voyageur, le client, le guide, les habitants \\
\hline
Lieux & das Hotel (-s), die Jugendherberge (-n), der Campingplatz (¨-e), der Strand (¨-e), das Meer (-e), das Museum (Museen) & l'hôtel, l'auberge de jeunesse, le camping, la plage, la mer, le musée \\
\hline
Activités & die Reise (-n), der Ausflug (¨-e), besuchen, besichtigen, wandern, baden, sich erholen & le voyage, l'excursion, visiter, visiter (un site), randonner, se baigner, se reposer \\
\hline
Problèmes & die Umweltverschmutzung, der Müll, die Preise steigen, die Natur zerstören, Arbeitsplätze schaffen & la pollution, les déchets, les prix augmentent, détruire la nature, créer des emplois \\
\hline
\end{tabularx}
\end{center}
\end{aretenir}

\courssub{Texte d'entraînement : Unsere Rundreise durch Kamerun}

Lisons maintenant un court texte qui réemploie tout ce vocabulaire en situation. Lis-le d'abord en entier pour le sens global, puis mot à mot avec le glossaire.

\begin{textebox}[Unsere Rundreise durch Kamerun]
Im Juli haben meine Familie und ich eine Rundreise durch Kamerun gemacht. Zuerst sind wir nach Douala gefahren. Dort haben wir zwei Nächte in einem Hotel am Hafen verbracht. Unser Reiseleiter, Herr Mbarga, hat uns die Stadt gezeigt: den Markt, den Hafen und das alte Viertel Bonanjo.

Dann sind wir an die Küste gefahren. In Kribi haben wir am Strand gelegen und im Meer gebadet. Der Strand war wunderschön, aber wir haben auch ein Problem gesehen: Manche Touristen lassen ihren Müll im Sand. Das ist schlecht für die Natur. Die Einheimischen sagen: „Der Tourismus ist wichtig für uns, er schafft Arbeitsplätze. Aber die Touristen müssen die Umwelt respektieren.“

Am dritten Tag haben wir einen Ausflug zu den Wasserfällen von Ekom gemacht. Diese Sehenswürdigkeit ist fantastisch! Wir haben viele Fotos gemacht und sind im Wald gewandert. Am Abend waren wir müde, aber glücklich.

Zum Schluss sind wir nach Yaoundé gefahren. Dort haben wir das Nationalmuseum besichtigt und Souvenirs gekauft. Die Preise in der Hauptstadt sind höher als in den Dörfern, aber die Stadt hat viel zu bieten.

Diese Reise war ein unvergessliches Erlebnis. Wir haben uns gut erholt und viel über unser Land gelernt. Kamerun ist wirklich einer Reise wert!
\end{textebox}

\textbf{Glossaire :} \all{der Hafen (¨-n)} le port (pl. die Häfen) ; \all{das Viertel (-)} le quartier (pl. die Viertel) ; \all{die Küste (-n)} la côte ; \all{der Wasserfall (¨-e)} la chute d'eau ; \all{das Souvenir (-s)} le souvenir ; \all{das Dorf (¨-er)} le village (pl. die Dörfer) ; \all{das Erlebnis (-se)} l'expérience, l'aventure ; \all{unvergesslich} inoubliable ; \all{etwas zu bieten haben} avoir quelque chose à offrir ; \all{einer Reise wert sein} valoir le voyage (génitif).

\textbf{Questions de compréhension :}
\begin{enumerate}
\item \all{Wo hat die Familie die ersten zwei Nächte verbracht ?}
\item \all{Welches Problem sieht die Familie am Strand von Kribi ?}
\item \all{Was sagen die Einheimischen über den Tourismus ?}
\item \all{Welche Sehenswürdigkeit besucht die Familie am dritten Tag ?}
\item \all{Was hat die Familie in Yaoundé gemacht ?}
\end{enumerate}

\begin{corrige}[Questions de compréhension]
\begin{enumerate}
\item \all{Sie hat die ersten zwei Nächte in einem Hotel am Hafen in Douala verbracht.} « Elle a passé les deux premières nuits dans un hôtel au port de Douala. »
\item \all{Manche Touristen lassen ihren Müll im Sand.} « Certains touristes laissent leurs déchets dans le sable. »
\item \all{Der Tourismus schafft Arbeitsplätze, aber die Touristen müssen die Umwelt respektieren.} « Le tourisme crée des emplois, mais les touristes doivent respecter l'environnement. »
\item \all{Sie besucht die Wasserfälle von Ekom.} « Elle visite les chutes d'Ekom. »
\item \all{Sie hat das Nationalmuseum besichtigt und Souvenirs gekauft.} « Elle a visité le musée national et acheté des souvenirs. »
\end{enumerate}
\end{corrige}

\courssub{Exercice résolu et entraînement}

\begin{exoresolu}[Complète avec le bon mot du lexique]
Complète les phrases avec : \all{Sehenswürdigkeiten}, \all{Reiseleiter}, \all{Jugendherberge}, \all{Müll}, \all{erholen}, \all{Reise}.

\begin{enumerate}
\item Im Sommer machen wir eine \ldots\ nach Deutschland.
\item Der \ldots\ zeigt den Touristen die Stadt.
\item Die Schüler schlafen in einer \ldots , weil sie nicht viel Geld haben.
\item Die Touristen besichtigen die \ldots\ der Hauptstadt.
\item Wirf deinen \ldots\ nicht auf den Strand !
\item Nach den Ferien wollen wir uns gut \ldots .
\end{enumerate}
\tcblower
\textbf{Solution détaillée :}
\begin{enumerate}
\item \all{eine Reise} — expression figée : \all{eine Reise machen} (faire un voyage), accusatif.
\item \all{Der Reiseleiter} — c'est la personne qui montre (\all{zeigt}) la ville aux touristes, nominatif.
\item \all{einer Jugendherberge} — datif après \all{in} (lieu où l'on dort) ; les élèves sans argent logent en auberge de jeunesse.
\item \all{die Sehenswürdigkeiten} — pluriel accusatif : on visite les sites de la capitale.
\item \all{deinen Müll} — accusatif : les déchets qu'on jette.
\item \all{erholen} — verbe réfléchi : \all{wir wollen uns erholen} (nous voulons nous reposer), infinitif sans \all{zu} après \all{wollen}.
\end{enumerate}
\end{exoresolu}

\begin{exercice}[À toi de jouer !]
\begin{enumerate}
\item Traduis en allemand : a) « Nous faisons un voyage à Kribi. » b) « Les touristes visitent les sites de la ville. » c) « Le tourisme crée des emplois. »
\item Trouve le bon article : \ldots Hotel, \ldots Strand, \ldots Meer, \ldots Museum, \ldots Reiseleiter, \ldots Umweltverschmutzung.
\item Construis la famille de mots de \all{besuchen} (deux noms) et de \all{sehen} (un nom, un adjectif).
\end{enumerate}
\end{exercice}

\begin{corrige}[Exercice « À toi de jouer ! »]
\begin{enumerate}
\item a) \all{Wir machen eine Reise nach Kribi.} b) \all{Die Touristen besichtigen die Sehenswürdigkeiten der Stadt.} c) \all{Der Tourismus schafft Arbeitsplätze.}
\item \all{das Hotel}, \all{der Strand}, \all{das Meer}, \all{das Museum}, \all{der Reiseleiter}, \all{die Umweltverschmutzung}.
\item \all{besuchen} $\rightarrow$ \all{der Besuch}, \all{der Besucher} ; \all{sehen} $\rightarrow$ \all{die Sehenswürdigkeit}, \all{sehenswert}.
\end{enumerate}
\end{corrige}

\begin{attention}[Pièges de traduction pour les francophones]
\begin{itemize}
\item « La mer » se dit \all{das Meer} (neutre !) et non \all{die Meer}. Mais « le lac » se dit \all{der See} — et « la mer » peut aussi se dire \all{die See} dans le langage des marins. Retiens : \all{ans Meer fahren} = aller au bord de la mer.
\item « Visiter » ne se traduit pas toujours par \all{visitieren} (ce verbe existe mais signifie « fouiller, perquisitionner » !). On dit \all{besuchen} (personnes) ou \all{besichtigen} (sites).
\item « Faire un voyage » = \all{eine Reise machen} (et non \all{eine Reise tun}).
\item N'oublie pas la majuscule à tous ces noms : \all{der Strand}, \all{das Hotel}, \all{die Reise}.
\end{itemize}
\end{attention}

\begin{savaistu}[D'où vient le mot « Urlaub » ?]
Le mot allemand \all{der Urlaub} (les congés, les vacances) vient du verbe \all{erlauben} (permettre, autoriser). À l'origine, \all{Urlaub} désignait littéralement la « permission » accordée à un soldat ou à un employé de s'absenter de son poste. Aujourd'hui encore, quand un Allemand dit \all{Ich habe Urlaub}, il dit en quelque sorte : « j'ai la permission de ne pas travailler » ! D'où l'expression \all{in den Urlaub fahren} : partir en vacances, littéralement « partir dans la permission ».
\end{savaistu}

Maintenant que le vocabulaire du tourisme est bien installé — acteurs, lieux, activités, problèmes et familles de mots —, nous pouvons passer à la grammaire : comment dire précisément \emph{où} l'on va et \emph{où} l'on est ? C'est le rôle des prépositions de lieu avec l'accusatif et le datif, que nous étudierons dans la section suivante.

\coursec{Grammaire 1 : les prépositions de lieu (accusatif/datif)}

Dans la section précédente, nous avons découvert le vocabulaire du tourisme : \all{der Tourist}, \all{das Hotel}, \all{die Reise}, \all{der Strand}, \all{die Sehenswürdigkeit}. Mais pour employer ce vocabulaire dans des phrases correctes, il faut savoir dire \emph{où l'on va} et \emph{où l'on est}. En allemand, ce choix n'est pas anodin : il détermine le \cle{cas} (accusatif ou datif) du nom qui suit la préposition. C'est l'une des règles les plus importantes — et les plus rentables — de votre apprentissage.

\courssub{Observation : partir d'exemples du texte}

Reprenons deux phrases tirées de notre texte support \emph{Tourismus – Segen oder Problem ?} :

\begin{exemplebox}[Deux phrases à comparer]
\all{Die Touristen reisen ans Meer.} « Les touristes vont à la mer. »\par
\all{Sie liegen am Strand.} « Ils sont allongés sur la plage. »
\end{exemplebox}

Observons attentivement. Dans les deux phrases, on parle du même lieu (la mer, la plage), et la préposition est la même à la base : \all{an}. Pourtant, la forme diffère : \all{an\textbf{s}} (an + das, accusatif) dans la première, \all{a\textbf{m}} (an + dem, datif) dans la seconde. Pourquoi ?

\begin{itemize}
\item Dans la première phrase, les touristes \emph{se déplacent} : ils quittent leur ville pour aller \emph{vers} la mer. Il y a \cle{mouvement}, une destination. La question est : \all{Wohin reisen die Touristen ?} (Vers où voyagent-ils ?)
\item Dans la seconde phrase, les touristes ne bougent pas : ils sont \emph{déjà} sur la plage et y restent. Il y a \cle{position}, un état. La question est : \all{Wo liegen sie ?} (Où sont-ils allongés ?)
\end{itemize}

Toute la règle tient dans cette opposition : \all{wohin ?} (vers où ?) appelle l'\cle{accusatif}, \all{wo ?} (où ?) appelle le \cle{datif}.

\begin{popfigure}
\begin{tikzpicture}[font=\small, node distance=1.2cm and 1.5cm]
% Nœud central : le lieu
\node[draw, thick, popblue, fill=popblueL, rounded corners=4pt, align=center, minimum width=3cm, minimum height=2cm] (lieu) {\textbf{der Strand / das Meer}};
% Flèche mouvement (gauche)
\node[align=center, popdark, left=of lieu, xshift=-1.5cm] (mouv_text) {\textbf{WOHIN ?} (vers où ?)\\ mouvement $\rightarrow$ \textbf{Akkusativ}};
\draw[-{Stealth[length=3mm]}, very thick, poporange] (mouv_text.east) -- (lieu.west);
\node[align=center, popdark, below=of mouv_text] (mouv_ex) {\all{Die Touristen reisen}\\ \all{\textbf{ans} Meer.}};
% Position (droite / bas)
\node[align=center, popdark, below=of lieu, yshift=-0.5cm] (pos_text) {\textbf{WO ?} (où ?)\\ position $\rightarrow$ \textbf{Dativ}};
\fill[popgreen] (lieu.south) circle (3pt);
\node[align=center, popdark, right=of lieu, xshift=1.5cm] (pos_ex) {\all{Sie liegen \textbf{am} Strand.}};
\draw[popgreen, thick] (lieu.south) -- (pos_text.north);
\end{tikzpicture}
\legende{Mouvement vers un lieu (accusatif) ou position dans un lieu (datif) : la question posée décide du cas.}
\end{popfigure}

\courssub{La règle : les prépositions à deux cas}

Neuf prépositions de lieu peuvent être suivies \emph{soit de l'accusatif, soit du datif}. On les appelle \cle{prépositions à deux cas} (en allemand : \all{Wechselpräpositionen}, « prépositions alternantes »).

\begin{propriete}[WOHIN ? ou WO ? — la règle d'or]
Avec les prépositions \all{an}, \all{auf}, \all{in}, \all{über}, \all{unter}, \all{vor}, \all{hinter}, \all{neben}, \all{zwischen} :
\begin{itemize}
\item si le groupe exprime un \textbf{mouvement vers un lieu} (question \all{wohin ?} « vers où ? »), on met l'\textbf{accusatif} ;
\item si le groupe exprime une \textbf{position, un état, un lieu où l'on se trouve} (question \all{wo ?} « où ? »), on met le \textbf{datif}.
\end{itemize}
Le verbe est souvent décisif : \all{gehen}, \all{fahren}, \all{fliegen}, \all{reisen}, \all{kommen}, \all{legen}, \all{stellen} (mouvement) contre \all{sein}, \all{liegen}, \all{sitzen}, \all{stehen}, \all{wohnen}, \all{bleiben}, \all{schlafen} (position).
\end{propriete}

Voici le tableau des neuf prépositions avec leurs sens principaux :

\begin{center}
\begin{tabularx}{\linewidth}{|l|X|X|}
\hline
\rowcolor{popblueL} \textbf{Préposition} & \textbf{Sens (mouvement / position)} & \textbf{Exemple} \\
\hline
\all{an} & à, au bord de (vertical ou rive) & \all{ans Meer fahren / am Meer sein} \\
\hline
\all{auf} & sur (surface) & \all{auf den Berg steigen / auf dem Berg stehen} \\
\hline
\all{in} & dans, en, à & \all{ins Hotel gehen / im Hotel wohnen} \\
\hline
\all{über} & au-dessus de & \all{über die Stadt fliegen / über der Stadt schweben} \\
\hline
\all{unter} & sous, parmi & \all{unter den Baum gehen / unter dem Baum sitzen} \\
\hline
\all{vor} & devant & \all{vor das Museum treten / vor dem Museum warten} \\
\hline
\all{hinter} & derrière & \all{hinter das Hotel gehen / hinter dem Hotel parken} \\
\hline
\all{neben} & à côté de & \all{sich neben den Guide setzen / neben dem Guide sitzen} \\
\hline
\all{zwischen} & entre & \all{zwischen die Stühle treten / zwischen den Stühlen stehen} \\
\hline
\end{tabularx}
\end{center}

Rappel des formes de l'article défini aux deux cas, car c'est l'article qui porte la marque du cas :

\begin{center}
\begin{tabular}{|l|l|l|l|}
\hline
\rowcolor{popblueL} & \textbf{Masculin} & \textbf{Féminin} & \textbf{Neutre} \\
\hline
Nominatif & der & die & das \\
\hline
Akkusativ & \textbf{den} & die & das \\
\hline
Dativ & \textbf{dem} & \textbf{der} & \textbf{dem} \\
\hline
\end{tabular}
\end{center}

\begin{exemplebox}[Paires modèles à mémoriser]
\all{Wir fahren in die Berge.} « Nous allons à la montagne. » (mouvement, accusatif : \all{die Berge}, pluriel identique)\par
\all{Wir sind in den Bergen.} « Nous sommes à la montagne. » (position, datif pluriel en \all{-n} : \all{den Bergen})\par
\all{Er geht ins Hotel.} « Il entre à l'hôtel. » (mouvement, accusatif)\par
\all{Er schläft im Hotel.} « Il dort à l'hôtel. » (position, datif)\par
\all{Sie hängt das Foto an die Wand.} « Elle accroche la photo au mur. » (mouvement)\par
\all{Das Foto hängt an der Wand.} « La photo est accrochée au mur. » (position)
\end{exemplebox}

Notez bien la différence avec le français : en français, « à la montagne », « à l'hôtel », « au mur » ne changent jamais de forme, que l'on y aille ou que l'on y soit. L'allemand, lui, \emph{code grammaticalement} la différence entre le but et le lieu. C'est une richesse qu'il faut apprendre à exploiter.

\courssub{Les contractions obligatoires}

En allemand parlé comme écrit, certaines prépositions fusionnent obligatoirement (ou presque) avec l'article. Ces \cle{contractions} sont à connaître par cœur :

\begin{propriete}[Contractions à connaître]
\begin{itemize}
\item \all{an + das = \textbf{ans}} (accusatif neutre) : \all{Wir fahren ans Meer.} « Nous allons à la mer. »
\item \all{an + dem = \textbf{am}} (datif masculin/neutre) : \all{Wir liegen am Strand.} « Nous sommes allongés sur la plage. »
\item \all{in + das = \textbf{ins}} (accusatif neutre) : \all{Sie geht ins Hotel.} « Elle entre à l'hôtel. »
\item \all{in + dem = \textbf{im}} (datif masculin/neutre) : \all{Sie wohnt im Hotel.} « Elle loge à l'hôtel. »
\end{itemize}
On trouve aussi, moins systématiques : \all{fürs} (für das), \all{übers} (über das), \all{unters} (unter das), \all{vorm} (vor dem), \all{hinterm} (hinter dem), \all{zum} (zu dem), \all{zur} (zu der), \all{beim} (bei dem).
\end{propriete}

\begin{attention}[Ne pas contracter quand l'article est emphatique]
Si l'article a une valeur démonstrative (« \emph{ce} musée-là, pas un autre »), on ne contracte pas : \all{Wir gehen in das Museum, das ich dir gezeigt habe.} « Nous allons dans ce musée que je t'ai montré. » Au niveau Première, retenez simplement que \all{ans}, \all{am}, \all{ins}, \all{im} sont la norme.
\end{attention}

\courssub{Les prépositions toujours suivies du datif}

Bonne nouvelle : certaines prépositions de lieu (et de déplacement) sont \emph{toujours} suivies du datif, qu'il y ait mouvement ou non. Inutile de se poser la question \all{wohin/wo} !

\begin{propriete}[Datif automatique : aus, bei, mit, nach, von, zu]
\begin{itemize}
\item \all{aus} : de, en provenance de — \all{Die Touristen kommen aus Frankreich.} « Les touristes viennent de France. »
\item \all{bei} : chez, auprès de — \all{Ich wohne bei meiner Tante.} « Je loge chez ma tante. » ; \all{beim Reiseleiter} « auprès du guide ».
\item \all{mit} : avec (moyen ou personne) — \all{Wir reisen mit dem Bus.} « Nous voyageons en bus. »
\item \all{nach} : vers (pays sans article, villes), après — \all{Wir fliegen nach Deutschland.} « Nous volons vers l'Allemagne. » ; \all{nach Hause} « à la maison » (mouvement).
\item \all{von} : de, depuis — \all{Er kommt von der Arbeit.} « Il revient du travail. »
\item \all{zu} : à, chez (personne, lieu précis) — \all{Sie geht zum Bahnhof.} « Elle va à la gare. » ; \all{zu Hause} « à la maison » (position).
\end{itemize}
\end{propriete}

\begin{exemplebox}[Application au thème du tourisme]
\all{Die Familie Ngono fliegt nach Deutschland.} « La famille Ngono part en avion pour l'Allemagne. »\par
\all{Sie reist mit dem Zug von München nach Wien.} « Elle voyage en train de Munich à Vienne. »\par
\all{Der Reiseleiter spricht bei der Gruppe über die Sehenswürdigkeiten.} « Le guide parle des curiosités devant le groupe. »\par
\all{Nach dem Ausflug gehen die Touristen zum Hotel zurück.} « Après l'excursion, les touristes retournent à l'hôtel. »
\end{exemplebox}

\courssub{Cas particulier : les noms de pays}

C'est un point classique du programme et des sujets d'examen. Deux constructions coexistent :

\begin{propriete}[nach ou in + pays ?]
\begin{itemize}
\item \all{nach} + pays \textbf{sans article} (la grande majorité : neutres) : \all{nach Kamerun}, \all{nach Deutschland}, \all{nach Frankreich}, \all{nach Österreich}. Et pour la position : \all{in Kamerun}, \all{in Deutschland} (datif sans article visible).
\item \all{in} + pays \textbf{avec article} (féminins, masculins ou pluriels) : \all{in die Schweiz} (mouvement) / \all{in der Schweiz} (position) ; \all{in den Senegal} / \all{im Senegal} ; \all{in die Vereinigten Staaten} / \all{in den Vereinigten Staaten}.
\end{itemize}
Sont avec article : \all{die Schweiz}, \all{die Türkei}, \all{die Elfenbeinküste}, \all{der Senegal}, \all{der Kongo}, \all{der Libanon}, \all{die USA} (pluriel).
\end{propriete}

\begin{exemplebox}[Paires à retenir]
\all{Wir fliegen in die Schweiz.} « Nous partons en Suisse. » / \all{Wir sind in der Schweiz.} « Nous sommes en Suisse. »\par
\all{Die Touristen reisen nach Kamerun.} « Les touristes voyagent au Cameroun. » / \all{Sie besuchen in Kamerun den Waza-Nationalpark.} « Au Cameroun, ils visitent le parc national de Waza. »\par
\all{Mein Onkel fliegt in den Senegal.} « Mon oncle part au Sénégal. » / \all{Er arbeitet im Senegal.} « Il travaille au Sénégal. »
\end{exemplebox}

\begin{methode}[Comment choisir le bon cas en trois étapes]
\begin{enumerate}
\item \textbf{Repérer la préposition.} Si c'est \all{aus}, \all{bei}, \all{mit}, \all{nach}, \all{von}, \all{zu} : datif automatique, problème résolu. Si c'est une préposition à deux cas (\all{an}, \all{auf}, \all{in}, etc.), passer à l'étape 2.
\item \textbf{Poser la question au verbe.} Demandez-vous : « vers où ? » (\all{wohin ?}) ou « où exactement ? » (\all{wo ?}). Si le verbe exprime un déplacement vers un but (\all{gehen}, \all{fahren}, \all{fliegen}, \all{legen}...), c'est un mouvement ; s'il exprime un état stable (\all{sein}, \all{wohnen}, \all{liegen}, \all{bleiben}...), c'est une position.
\item \textbf{Appliquer la règle et vérifier l'article.} Mouvement $\rightarrow$ accusatif (\all{den}, \all{die}, \all{das}) ; position $\rightarrow$ datif (\all{dem}, \all{der}, \all{dem} ; pluriel \all{den + n}). Pensez aux contractions : \all{ins}, \all{im}, \all{ans}, \all{am}.
\end{enumerate}
\end{methode}

\courssub{Texte d'application : Un voyage en Suisse}

Lisons maintenant un court texte qui met en scène toutes ces prépositions. Soulignez mentalement chaque préposition de lieu et justifiez son cas.

\begin{textebox}[Familie Kamdem reist in die Schweiz]
\all{Im Juli fliegt die Familie Kamdem von Douala in die Schweiz. Der Flug dauert lange, aber die Kinder sind voller Freude. Am Flughafen in Zürich wartet ein Reiseleiter bei der Gruppe. „Willkommen in der Schweiz !“, sagt er freundlich.}\par
\all{Die Familie fährt mit dem Bus ins Hotel. Das Hotel liegt an einem See, zwischen den Bergen. Die Kinder laufen sofort ans Wasser, aber das Wasser ist kalt ! Am Abend essen alle im Restaurant. Am nächsten Tag steigen sie auf einen Berg. Oben auf dem Berg machen sie viele Fotos.}\par
\all{Nach einer Woche reist die Familie nach Kamerun zurück. „Die Schweiz ist schön“, sagt der Vater, „aber zu Hause ist es am schönsten.“}
\end{textebox}

\textbf{Glossaire}

\begin{itemize}
\item \all{der Flughafen, die Flughäfen} : l'aéroport
\item \all{der Reiseleiter, die Reiseleiter} : le guide (accompagnateur)
\item \all{der See, die Seen} : le lac
\item \all{der Berg, die Berge} : la montagne
\item \all{steigen (stieg, ist gestiegen)} : monter, gravir
\item \all{oben} : en haut
\item \all{zurückreisen} : retourner (voyager en retour)
\item \all{am schönsten} : le plus beau (superlatif de \all{schön})
\end{itemize}

\textbf{Questions de compréhension}

\begin{enumerate}
\item \all{Wohin fliegt die Familie Kamdem ?} (Où la famille Kamdem part-elle en avion ?)
\item \all{Wo liegt das Hotel ?} (Où se trouve l'hôtel ?)
\item \all{Was machen die Kinder am Wasser ?} (Que font les enfants au bord de l'eau ?)
\item \all{Wohin reist die Familie nach einer Woche ?} (Où la famille repart-elle après une semaine ?)
\end{enumerate}

Relevons ensemble les prépositions du texte : \all{in die Schweiz} (mouvement, accusatif — pays à article), \all{bei der Gruppe} (datif automatique avec \all{bei}), \all{in der Schweiz} (position, datif), \all{mit dem Bus} (datif automatique avec \all{mit}), \all{ins Hotel} (mouvement), \all{an einem See} (position : l'hôtel \emph{se trouve} au bord du lac), \all{zwischen den Bergen} (position, datif pluriel en \all{-n}), \all{ans Wasser} (mouvement : les enfants courent \emph{vers} l'eau), \all{im Restaurant} (position), \all{auf einen Berg} (mouvement : ils gravissent), \all{auf dem Berg} (position : ils sont en haut), \all{nach Kamerun} (pays sans article), \all{zu Hause} (position, datif figé).

\courssub{Exercice résolu et pièges à éviter}

\begin{exoresolu}[Complétez avec la bonne forme]
Complétez avec l'article correct (accusatif ou datif) et justifiez.\par
1. \all{Die Touristen gehen in \dots{} Museum.} (das Museum)\par
2. \all{Die Touristen sitzen in \dots{} Museum und trinken Kaffee.} (das Museum)\par
3. \all{Wir fahren an \dots{} Meer.} (das Meer)\par
4. \all{Wir bleiben zwei Wochen an \dots{} Meer.} (das Meer)\par
5. \all{Sie fliegt in \dots{} Schweiz.} (die Schweiz)\par
6. \all{Sie wohnt in \dots{} Schweiz.} (die Schweiz)
\tcblower
\textbf{Solution détaillée}\par
1. \all{Die Touristen gehen in \textbf{das} Museum} (mieux : \all{\textbf{ins} Museum}). Le verbe \all{gehen} exprime un mouvement vers un lieu (\all{wohin ?}) $\rightarrow$ accusatif ; contraction \all{in + das = ins}.\par
2. \all{Die Touristen sitzen in \textbf{dem} Museum} (mieux : \all{\textbf{im} Museum}). \all{sitzen} exprime une position (\all{wo ?}) $\rightarrow$ datif ; contraction \all{in + dem = im}.\par
3. \all{Wir fahren an \textbf{das} Meer} (\all{\textbf{ans} Meer}). \all{fahren} = mouvement $\rightarrow$ accusatif ; \all{an + das = ans}.\par
4. \all{Wir bleiben zwei Wochen an \textbf{dem} Meer} (\all{\textbf{am} Meer}). \all{bleiben} = position, durée sur place $\rightarrow$ datif ; \all{an + dem = am}.\par
5. \all{Sie fliegt in \textbf{die} Schweiz.} Pays avec article + mouvement $\rightarrow$ accusatif féminin : \all{die} (forme identique au nominatif, attention !).\par
6. \all{Sie wohnt in \textbf{der} Schweiz.} Position $\rightarrow$ datif féminin : \all{der}.
\end{exoresolu}

\begin{attention}[L'erreur n° 1 des francophones]
Ne dites jamais \all{*Wir fliegen in der Schweiz} pour « Nous partons en Suisse » ! En français, « en Suisse » sert pour les deux idées, mais l'allemand distingue : le \textbf{mouvement exige l'accusatif} — \all{Wir fliegen in \textbf{die} Schweiz} — tandis que \all{in \textbf{der} Schweiz} signifie « en Suisse » (on y est déjà). Autres pièges fréquents :
\begin{itemize}
\item oublier le \all{-n} du datif pluriel : \all{in den Berge\textbf{n}} (et non \all{*in den Berge}) ;
\item employer \all{nach} avec un pays à article : on dit \all{in die Schweiz}, jamais \all{*nach der Schweiz} ;
\item oublier que \all{zu} et \all{bei} prennent toujours le datif : \all{zum Bahnhof}, \all{beim Reiseleiter}.
\end{itemize}
\end{attention}

\begin{exercice}[À vous de jouer]
Traduisez en allemand en faisant attention au cas :\par
1. Nous allons à l'hôtel. 2. Nous sommes à l'hôtel. 3. Les touristes voyagent au Cameroun. 4. Ils visitent le marché à Yaoundé. 5. Elle met la valise sur le lit. 6. La valise est sur le lit.
\end{exercice}

\begin{corrige}[Exercice « À vous de jouer »]
1. \all{Wir gehen ins Hotel.} (mouvement, accusatif)\par
2. \all{Wir sind im Hotel.} (position, datif)\par
3. \all{Die Touristen reisen nach Kamerun.} (pays sans article)\par
4. \all{Sie besuchen den Markt in Yaoundé.} (\all{besuchen} + accusatif ; ville sans article)\par
5. \all{Sie stellt den Koffer auf das Bett.} (mouvement : \all{den Koffer} est l'objet direct, \all{auf das Bett} = destination, accusatif)\par
6. \all{Der Koffer ist auf dem Bett.} (position, datif)
\end{corrige}

\begin{savaistu}[Les prépositions de lieu au Bac]
Au baccalauréat camerounais, les prépositions de lieu sont presque toujours testées dans les questions de grammaire qui accompagnent le texte de compréhension — et les textes de voyage ou de tourisme y reviennent régulièrement ! Un exercice typique : « Complétez : \all{Die Familie fährt in \dots{} Berge / Sie wohnt in \dots{} Bergen.} » Maîtriser la distinction \all{wohin/wo}, c'est donc sécuriser des points faciles à l'examen. À retenir aussi : les Allemands disent \all{auf dem Land} (à la campagne, position) et \all{aufs Land} (à la campagne, mouvement) — la même logique, partout.
\end{savaistu}

\begin{aretenir}[L'essentiel de la section]
\begin{itemize}
\item Neuf prépositions à deux cas : \all{an}, \all{auf}, \all{in}, \all{über}, \all{unter}, \all{vor}, \all{hinter}, \all{neben}, \all{zwischen}.
\item Mouvement (\all{wohin ?}) $\rightarrow$ \textbf{accusatif} ; position (\all{wo ?}) $\rightarrow$ \textbf{datif}.
\item Contractions : \all{ans}, \all{am}, \all{ins}, \all{im}.
\item Toujours datif : \all{aus}, \all{bei}, \all{mit}, \all{nach}, \all{von}, \all{zu}.
\item Pays sans article : \all{nach Kamerun} ; pays avec article : \all{in die Schweiz} / \all{in der Schweiz}.
\end{itemize}
\end{aretenir}

\coursec{Grammaire 2 : weil et dass + parfait avec sein}

Dans la section précédente, nous avons appris à dire \emph{où} l'on va et \emph{où} l'on se trouve grâce aux prépositions de lieu : \all{Wir fahren \textbf{ans} Meer} (mouvement, accusatif), \all{Wir sind \textbf{am} Meer} (position, datif). Mais un bon touriste ne dit pas seulement où il va : il explique aussi \emph{pourquoi} il voyage et \emph{ce qu'il pense} du tourisme. Pour cela, l'allemand utilise deux petits mots essentiels : \cle{weil} (parce que) et \cle{dass} (que). Et pour raconter ses voyages passés, il faut maîtriser le parfait avec l'auxiliaire \cle{sein}. C'est l'objet de cette section.

\courssub{Observation : le verbe qui recule}

\begin{textebox}[Deux phrases du texte support]
\all{Viele Touristen reisen ans Meer, weil sie die Sonne lieben.}

\all{Ich denke, dass der Ökotourismus eine gute Lösung ist.}
\end{textebox}

Traduisons : « Beaucoup de touristes vont à la mer parce qu'ils aiment le soleil. » et « Je pense que l'écotourisme est une bonne solution. »

Observons attentivement la seconde partie de chaque phrase :

\begin{itemize}
\item \all{weil sie die Sonne \textbf{lieben}} : le verbe conjugué \all{lieben} se trouve à la \textbf{dernière place}.
\item \all{dass der Ökotourismus eine gute Lösung \textbf{ist}} : le verbe conjugué \all{ist} se trouve à la \textbf{dernière place}.
\end{itemize}

En français, rien de tel : « parce qu'ils \emph{aiment} le soleil » — le verbe reste juste après le sujet. En allemand, après \all{weil} et \all{dass}, le verbe conjugué \emph{recule} jusqu'à la fin de la phrase. C'est la règle d'or de la subordination.

\begin{propriete}[La règle d'or de la subordonnée]
Après \cle{weil} et \cle{dass}, le verbe conjugué occupe la \textbf{DERNIÈRE place} de la subordonnée :

\all{..., weil sie die Sonne \textbf{lieben}.}

\all{..., dass er nach Kribi \textbf{fährt}.}

Structure : [phrase principale : S + V + ...] + [weil/dass + S + ... + \textbf{V en fin}].

Remarque : la subordonnée est toujours séparée de la principale par une \textbf{virgule}.
\end{propriete}

\begin{popfigure}
\begin{tikzpicture}[font=\small]
\node[draw=popblue, fill=popblueL, rounded corners, align=center, minimum width=4.6cm, minimum height=1.1cm] (main) at (0,0) {\textbf{Phrase principale}\\ S + V + ...};
\node[draw=poporange, fill=poporangeL, rounded corners, align=center, minimum width=1.6cm, minimum height=1.1cm] (conj) at (4.1,0) {\textbf{weil}\\ \textbf{dass}};
\node[draw=popgreen, fill=popgreenL, rounded corners, align=center, minimum width=5.4cm, minimum height=1.1cm] (sub) at (8.6,0) {S + ... + \textbf{V en dernière place}};
\node[draw=popink, fill=popink!10, rounded corners, align=center] (verb) at (10.9,-1.7) {\textbf{verbe conjugué}};
\draw[-{Stealth}, thick, popink] (main) -- (conj);
\draw[-{Stealth}, thick, popink] (conj) -- (sub);
\draw[-{Stealth}, very thick, poporange, bend left=25] (8.0,-0.6) to node[above, popdark, align=center]{\footnotesize le verbe recule !} (verb.north);
\node[align=center, popmuted] at (4.5,-2.4) {\all{Viele reisen ans Meer, \textbf{weil} sie die Sonne \textbf{lieben}.}};
\end{tikzpicture}
\legende{Structure de la phrase complexe : après weil/dass, le verbe conjugué recule en dernière position.}
\end{popfigure}

\courssub{La subordonnée de cause avec weil}

\begin{definition}[weil]
\cle{weil} signifie « parce que ». Il introduit la \textbf{cause}, la \textbf{raison} : il répond à la question \all{Warum?} (pourquoi ?).
\end{definition}

\begin{exemplebox}[weil : exprimer la cause]
\all{Viele reisen ans Meer, weil sie die Sonne lieben.} = Beaucoup vont à la mer parce qu'ils aiment le soleil.

\all{Die Preise steigen, weil viele Touristen kommen.} = Les prix montent parce que beaucoup de touristes viennent.

\all{Wir bleiben zu Hause, weil das Wetter schlecht ist.} = Nous restons à la maison parce que le temps est mauvais.

\all{Der Ökotourismus wächst, weil die Menschen die Natur schützen wollen.} = L'écotourisme se développe parce que les gens veulent protéger la nature.
\end{exemplebox}

Remarque importante : la subordonnée en \all{weil} peut aussi se placer \emph{avant} la phrase principale. Dans ce cas, la subordonnée entière occupe la position 1, et le verbe de la phrase principale vient \emph{immédiatement après la virgule} (position 2) :

\all{Weil sie die Sonne lieben, \textbf{reisen} viele Touristen ans Meer.} = Parce qu'ils aiment le soleil, beaucoup de touristes vont à la mer.

\courssub{La subordonnée d'opinion avec dass}

\begin{definition}[dass]
\cle{dass} signifie « que ». Il introduit une \textbf{opinion}, une \textbf{pensée} ou un \textbf{fait rapporté}, après des verbes comme \all{denken} (penser), \all{glauben} (croire), \all{sagen} (dire), \all{hoffen} (espérer), \all{meinen} (penser, estimer), \all{wissen} (savoir).
\end{definition}

\begin{exemplebox}[dass : exprimer une opinion]
\all{Ich denke, dass der Ökotourismus eine gute Lösung ist.} = Je pense que l'écotourisme est une bonne solution.

\all{Er sagt, dass er im Sommer nach Österreich gefahren ist.} = Il dit qu'il est allé en Autriche en été.

\all{Wir hoffen, dass der Massentourismus die Natur nicht zerstört.} = Nous espérons que le tourisme de masse ne détruit pas la nature.

\all{Sie glaubt, dass Kamerun viele schöne Sehenswürdigkeiten hat.} = Elle croit que le Cameroun a beaucoup de belles curiosités.
\end{exemplebox}

\begin{attention}[Orthographe : dass ou das ?]
\all{dass} (conjonction « que ») s'écrit avec \textbf{ss} ; \all{das} (article ou pronom « le/ce ») avec un seul \textbf{s}. Test : si l'on peut remplacer le mot par \all{dieses} (« ce »), on écrit \all{das} ; sinon, \all{dass}.

\all{Ich weiß, dass das Hotel teuer ist.} = Je sais que l'hôtel est cher. (premier \all{dass} = conjonction ; second \all{das} = article devant \all{Hotel})
\end{attention}

\begin{attention}[Piège classique : weil ou denn ?]
Le français « car / parce que » a deux traductions en allemand :

\begin{itemize}
\item \cle{denn} : coordination, le verbe reste en \textbf{position 2} : \all{Viele reisen ans Meer, denn sie \textbf{lieben} die Sonne.}
\item \cle{weil} : subordination, le verbe va \textbf{en fin} : \all{Viele reisen ans Meer, weil sie die Sonne \textbf{lieben}.}
\end{itemize}

Au Bac, écrire \all{weil sie lieben die Sonne} est une faute grave de syntaxe. Après \all{weil}, le verbe recule, toujours !
\end{attention}

\begin{center}
\begin{tabularx}{\linewidth}{|X|X|X|}
\hline
\rowcolor{popblueL} Français & Deutsch & Remarque \\
\hline
parce qu'ils aiment le soleil & weil sie die Sonne \textbf{lieben} & verbe en fin (subordination) \\
\hline
car ils aiment le soleil & denn sie \textbf{lieben} die Sonne & verbe en position 2 (coordination) \\
\hline
je pense qu'il vient & ich denke, dass er \textbf{kommt} & verbe en fin après dass \\
\hline
il dit qu'il est allé & er sagt, dass er gefahren \textbf{ist} & auxiliaire en fin au parfait \\
\hline
\end{tabularx}
\end{center}

\courssub{Le parfait avec sein : raconter ses voyages}

Pour raconter un voyage passé, l'allemand utilise le parfait (\all{Perfekt}) : auxiliaire au présent en position 2 + participe passé en fin de phrase. La plupart des verbes prennent l'auxiliaire \all{haben} : \all{Wir \textbf{haben} die Stadt \textbf{besucht}.} Mais les verbes de \textbf{déplacement} et de \textbf{changement d'état} prennent l'auxiliaire \cle{sein}.

\begin{propriete}[Le parfait avec sein]
On conjugue le parfait avec \all{sein} (et non \all{haben}) pour :

\begin{enumerate}
\item les verbes de \textbf{mouvement / déplacement} : \all{reisen} (voyager), \all{fahren} (aller en véhicule), \all{fliegen} (voler, prendre l'avion), \all{gehen} (aller à pied), \all{kommen} (venir), \all{aufsteigen} (monter) ;
\item les verbes de \textbf{changement d'état} : \all{sterben} (mourir), \all{einschlafen} (s'endormir), \all{aufstehen} (se lever) ;
\item deux verbes à retenir par cœur : \all{sein} lui-même (\all{ich bin gewesen}) et \all{bleiben} (rester : \all{ich bin geblieben}).
\end{enumerate}

Forme : \textbf{sein au présent (position 2) + participe passé en fin de phrase}.
\end{propriete}

\begin{center}
\begin{tabular}{|l|l|l|}
\hline
\rowcolor{popgreenL} Personne & sein (Präsens) & Partizip II (exemples) \\
\hline
ich & bin & gereist \\
\hline
du & bist & gefahren \\
\hline
er / sie / es & ist & geflogen \\
\hline
wir & sind & gegangen \\
\hline
ihr & seid & gekommen \\
\hline
sie / Sie & sind & geblieben \\
\hline
\end{tabular}
\end{center}

\begin{aretenir}[Participes passés irréguliers à connaître]
\begin{itemize}
\item \all{reisen} $\rightarrow$ \all{gereist} (régulier) ; \all{fahren} $\rightarrow$ \all{gefahren} ; \all{fliegen} $\rightarrow$ \all{geflogen}
\item \all{gehen} $\rightarrow$ \all{gegangen} ; \all{kommen} $\rightarrow$ \all{gekommen} ; \all{bleiben} $\rightarrow$ \all{geblieben} ; \all{sein} $\rightarrow$ \all{gewesen}
\end{itemize}
Astuce : les verbes forts forment leur participe avec \all{ge- ... -en} et changent souvent de voyelle ; les verbes faibles avec \all{ge- ... -t} (\all{reisen} $\rightarrow$ \all{gereist}, \all{besuchen} $\rightarrow$ \all{besucht}).
\end{aretenir}

\begin{exemplebox}[Le parfait avec sein en contexte touristique]
\all{Ich bin nach Kribi gereist.} = J'ai voyagé à Kribi / Je suis allé à Kribi.

\all{Wir sind ans Meer gefahren.} = Nous sommes allés à la mer (en voiture).

\all{Sie sind nach Deutschland geflogen.} = Ils ont pris l'avion pour l'Allemagne.

\all{Der Bus ist um sechs Uhr gekommen.} = Le bus est arrivé à six heures.

\all{Ich bin eine Woche am Strand geblieben.} = Je suis resté une semaine à la plage.
\end{exemplebox}

\begin{attention}[sein ou haben ? Le critère du mouvement]
Le parfait avec \all{sein} concerne les verbes de \textbf{MOUVEMENT} (et de changement d'état). Les verbes d'action ordinaires se conjuguent avec \all{haben} :

\all{Wir \textbf{haben} die Stadt \textbf{besucht}.} = Nous avons visité la ville.

\all{Ich \textbf{habe} viele Fotos \textbf{gemacht}.} = J'ai pris beaucoup de photos.

\all{Sie \textbf{haben} ein Hotel \textbf{gebucht}.} = Ils ont réservé un hôtel.

Erreur fréquente des francophones : le français « être » dans « je suis allé » aide à retenir \all{ich bin gegangen}, mais attention : « j'ai visité » = \all{ich habe besucht}, jamais \all{ich bin besucht} !
\end{attention}

\courssub{Combiner les trois points : une phrase complète}

Réunissons maintenant les prépositions de lieu, la subordonnée en \all{weil} et le parfait avec \all{sein} dans une seule phrase :

\all{Wir \textbf{sind} nach Kribi \textbf{gefahren}, weil wir den Strand sehen \textbf{wollten}.} = Nous sommes allés à Kribi parce que nous voulions voir la plage.

Analysons : \all{sind ... gefahren} = parfait avec \all{sein} (verbe de mouvement) ; \all{nach Kribi} = préposition de lieu (villes sans article) ; \all{weil wir ... wollten} = subordonnée de cause, verbe conjugué \all{wollten} en dernière place. Trois règles, une seule phrase correcte !

\begin{methode}[Rédiger une phrase complexe en quatre étapes]
\begin{enumerate}
\item Écrire d'abord la \textbf{phrase principale} : \all{Wir sind nach Kribi gefahren.}
\item Choisir la conjonction : \cle{weil} pour une \textbf{cause} (pourquoi ?), \cle{dass} pour une \textbf{opinion} (je pense que...).
\item Dans la subordonnée, placer le \textbf{verbe conjugué à la toute fin} : \all{..., weil wir den Strand sehen wollten.}
\item Vérifier l'\textbf{auxiliaire du parfait} : \all{sein} si le verbe exprime un déplacement (\all{ich bin gefahren}), \all{haben} sinon (\all{ich habe besucht}).
\end{enumerate}
\end{methode}

\begin{textebox}[Meine Reise nach Kribi]
\all{Letzten Sommer bin ich mit meiner Familie nach Kribi gefahren, weil wir das Meer und den Strand lieben. Wir sind mit dem Bus gereist, und die Fahrt hat vier Stunden gedauert. Ich denke, dass Kribi einer der schönsten Orte Kameruns ist. Wir haben in einem kleinen Hotel am Strand gewohnt und viele Fotos gemacht. Mein Bruder ist jeden Morgen früh aufgestanden, weil er den Sonnenaufgang sehen wollte. Er sagt, dass der Ökotourismus für Kamerun sehr wichtig ist, denn viele Touristen wollen die Natur schützen. Ich glaube, dass er Recht hat: der Massentourismus kann die Strände verschmutzen, weil zu viele Menschen gleichzeitig kommen. Am Ende sind wir nach Yaoundé zurückgefahren, aber wir sind noch eine Woche in der Region geblieben. Ich hoffe, dass wir nächstes Jahr wieder ans Meer fahren!}

\emph{Glossaire} : die Fahrt, -en (le trajet) ; dauern (durer) ; der Ort, -e (le lieu) ; wohnen (loger, habiter) ; der Sonnenaufgang, die Sonnenaufgänge (le lever du soleil) ; Recht haben (avoir raison) ; verschmutzen (polluer) ; zurückfahren (rentrer) ; gleichzeitig (en même temps).
\end{textebox}

Questions de compréhension : 1) \all{Warum ist die Familie nach Kribi gefahren?} 2) \all{Wie lange hat die Fahrt gedauert?} 3) \all{Was sagt der Bruder über den Ökotourismus?} 4) \all{Welches Problem nennt der Text beim Massentourismus?}

Réponses attendues : 1) \all{Sie ist nach Kribi gefahren, weil sie das Meer und den Strand liebt.} 2) \all{Die Fahrt hat vier Stunden gedauert.} 3) \all{Er sagt, dass der Ökotourismus für Kamerun sehr wichtig ist.} 4) \all{Der Massentourismus kann die Strände verschmutzen.}

\begin{exoresolu}[Construisez des phrases complexes]
Énoncé : reliez les phrases avec \all{weil} ou \all{dass}, puis mettez la phrase principale au parfait quand il s'agit d'un voyage passé.

a) Viele Touristen kommen nach Kamerun. Sie wollen die Natur sehen. (weil)

b) Ich denke ... Der Ökotourismus ist eine gute Lösung. (dass)

c) Wir / nach Österreich / fliegen (Perfekt) — wir wollen die Berge sehen. (weil)

\tcblower
\textbf{Solution détaillée.}

a) Cause $\rightarrow$ \all{weil} ; le verbe \all{wollen} recule en fin : \all{Viele Touristen kommen nach Kamerun, weil sie die Natur sehen \textbf{wollen}.} = Beaucoup de touristes viennent au Cameroun parce qu'ils veulent voir la nature.

b) Opinion après \all{denken} $\rightarrow$ \all{dass} ; le verbe \all{ist} recule en fin : \all{Ich denke, dass der Ökotourismus eine gute Lösung \textbf{ist}.}

c) \all{fliegen} est un verbe de déplacement $\rightarrow$ auxiliaire \all{sein} : \all{Wir \textbf{sind} nach Österreich \textbf{geflogen}, weil wir die Berge sehen \textbf{wollten}.} = Nous avons pris l'avion pour l'Autriche parce que nous voulions voir les montagnes. Vérification : participe irrégulier \all{geflogen}, verbe de la subordonnée \all{wollten} bien en dernière place.
\end{exoresolu}

\begin{exercice}[À vous de jouer]
1) Traduisez : « Beaucoup de touristes sont allés à la mer parce qu'ils aiment le soleil. »

2) Complétez avec \all{weil}, \all{dass} ou \all{denn} : a) \all{Ich glaube, ... der Massentourismus ein Problem ist.} b) \all{Wir bleiben im Hotel, ... es regnet.} c) \all{Sie reist gern, ... sie liebt neue Länder.}

3) Mettez au parfait : a) \all{ich / reisen / nach Douala} ; b) \all{wir / besuchen / die Stadt} ; c) \all{sie / fliegen / in die Schweiz}.
\end{exercice}

\begin{corrige}[Exercice « À vous de jouer »]
1) \all{Viele Touristen sind ans Meer gefahren, weil sie die Sonne lieben.} (mouvement $\rightarrow$ \all{sein} ; verbe \all{lieben} en fin de subordonnée).

2) a) \all{dass} (opinion après \all{glauben}) ; b) \all{weil} (cause, verbe \all{regnet} en fin) ; c) \all{denn} (verbe \all{liebt} en position 2, donc coordination).

3) a) \all{Ich bin nach Douala gereist.} b) \all{Wir haben die Stadt besucht.} (pas de mouvement $\rightarrow$ \all{haben}) ; c) \all{Sie sind in die Schweiz geflogen.} (pays avec article $\rightarrow$ \all{in die} + accusatif, mouvement).
\end{corrige}

\begin{aretenir}[L'essentiel de la section]
\begin{itemize}
\item Après \cle{weil} (cause) et \cle{dass} (opinion), le verbe conjugué va à la \textbf{dernière place}.
\item \cle{denn} = coordination : le verbe reste en position 2.
\item Parfait avec \cle{sein} : verbes de déplacement et de changement d'état (\all{ich bin gereist / gefahren / geflogen / gegangen / gekommen / geblieben}).
\item Tous les autres verbes prennent \all{haben} : \all{Wir haben die Stadt besucht.}
\end{itemize}
\end{aretenir}

\coursec{Méthode du commentaire et commentaire modèle}

Dans la section précédente, nous avons appris à construire des subordonnées avec \all{weil} et \all{dass} et à raconter un voyage au parfait avec l'auxiliaire \all{sein}. Ces outils grammaticaux ne sont pas une fin en soi : ils vont maintenant nous servir à \emph{parler d'un texte}. En effet, à l'examen comme en classe, on ne vous demande pas seulement de comprendre le texte \all{Tourismus – Segen oder Problem?}, mais aussi de le \emph{commenter} : dire de quoi il parle, quelle est l'idée de l'auteur, et ce que vous, vous en pensez. C'est exactement là que \all{weil} (pour justifier) et \all{dass} (pour rapporter une opinion) deviennent indispensables. Cette section vous donne la méthode complète, pas à pas, puis un commentaire modèle rédigé que vous pourrez imiter.

\courssub{Qu'est-ce qu'un commentaire de texte ?}

\begin{definition}[Le commentaire de texte]
Un \cle{commentaire de texte} (\all{der Kommentar}, \all{die Stellungnahme}) est un court texte dans lequel on : 1) \emph{présente} le document (thème, type, source), 2) dégage son \emph{idée principale} (\all{die Hauptidee}), 3) relève les \emph{arguments} pour et contre (\all{die Argumente}), et 4) donne son \emph{avis personnel} justifié (\all{die eigene Meinung}). En allemand, on dit \all{einen Text kommentieren} ou \all{zu einem Text Stellung nehmen} (prendre position sur un texte).
\end{definition}

Observez d'abord ces deux phrases, l'une maladroite, l'autre correcte :

\all{Der Text ist gut. Ich mag Tourismus.} — « Le texte est bien. J'aime le tourisme. » (trop pauvre : pas de présentation, pas d'argument, pas de justification)

\all{Der Text handelt vom Tourismus und seinen Problemen. Meiner Meinung nach ist der Ökotourismus wichtig, weil er die Natur schützt.} — « Le texte traite du tourisme et de ses problèmes. À mon avis, l'écotourisme est important, parce qu'il protège la nature. » (présentation + avis + justification avec \all{weil})

La différence est claire : un bon commentaire est \emph{structuré} et \emph{justifié}. Voyons la méthode en quatre étapes.

\courssub{Les quatre étapes du commentaire}

\begin{methode}[Commenter un texte en quatre étapes]
\begin{enumerate}
\item \textbf{Présenter le texte} : indiquez le titre, le thème et le type de document. Utilisez la formule \all{Der Text handelt von...} (le texte traite de...). Attention : après \all{von}, le nom est au \emph{datif} : \all{von dem Tourismus} $\rightarrow$ \all{vom Tourismus}.
\item \textbf{Dégager l'idée principale} (\all{die Hauptidee}) : demandez-vous « Que veut dire l'auteur ? » et rapportez sa pensée avec \all{dass} : \all{Der Autor meint, dass...} / \all{Der Autor zeigt, dass...} (verbe conjugué en fin de subordonnée !).
\item \textbf{Relever les arguments} pour et contre : utilisez le couple de connecteurs \all{einerseits... andererseits...} (d'un côté... de l'autre...). Citez le texte entre guillemets allemands : \all{Der Autor sagt: „...“}.
\item \textbf{Donner son avis justifié} : annoncez votre opinion avec \all{Meiner Meinung nach...} ou \all{Ich bin der Meinung, dass...}, puis justifiez \emph{toujours} avec une subordonnée en \all{weil}.
\end{enumerate}
\end{methode}

\begin{popfigure}
\begin{center}
\begin{tikzpicture}[
  node distance=0.55cm,
  etape/.style={rectangle, rounded corners=4pt, draw=popblue, fill=popblueL, text width=4.6cm, align=center, minimum height=1.1cm, font=\small},
  fleche/.style={-{Stealth[length=3mm]}, thick, poporange}
]
\node[etape, align=center] (e1){\textbf{1. Présentation}\\ Der Text handelt von...};
\node[etape, below=of e1, align=center] (e2){\textbf{2. Idée principale}\\ Der Autor meint, dass...};
\node[etape, below=of e2, align=center] (e3){\textbf{3. Arguments + / $-$}\\ Einerseits... andererseits...};
\node[etape, below=of e3, draw=poppurple, fill=poppurpleL, align=center] (e4){\textbf{4. Avis personnel}\\ Meiner Meinung nach... weil...};
\draw[fleche] (e1) -- (e2);
\draw[fleche] (e2) -- (e3);
\draw[fleche] (e3) -- (e4);
\end{tikzpicture}
\end{center}
\legende{Organigramme du commentaire de texte : de la présentation à l'avis personnel.}
\end{popfigure}

\courssub{Les expressions-outils du commentaire}

Voici la « boîte à outils » du commentateur. Apprenez ces formules par cœur : elles sont réutilisables avec n'importe quel texte, à l'écrit comme à l'oral.

\begin{center}
\begin{tabularx}{\linewidth}{|X|X|X|}
\hline
\rowcolor{popblueL} Étape & Deutsch & Français \\
\hline
Présentation & Der Text handelt vom Tourismus. & Le texte traite du tourisme. \\
\hline
Présentation & Es geht um die Probleme des Tourismus. & Il s'agit des problèmes du tourisme. \\
\hline
Idée principale & Der Autor meint, dass der Tourismus wichtig ist. & L'auteur pense que le tourisme est important. \\
\hline
Idée principale & Der Autor zeigt, dass Hotels die Strände verschmutzen. & L'auteur montre que les hôtels polluent les plages. \\
\hline
Arguments & Einerseits bringt der Tourismus Geld, andererseits verschmutzt er die Umwelt. & D'un côté le tourisme rapporte de l'argent, de l'autre il pollue l'environnement. \\
\hline
Avis & Meiner Meinung nach ist der Ökotourismus die beste Lösung. & À mon avis, l'écotourisme est la meilleure solution. \\
\hline
Avis & Ich bin der Meinung, dass wir die Natur schützen müssen. & Je suis d'avis que nous devons protéger la nature. \\
\hline
Justification & ..., weil er die Natur schützt und Arbeit bringt. & ..., parce qu'il protège la nature et crée du travail. \\
\hline
\end{tabularx}
\end{center}

\begin{exemplebox}[Exemples traduits]
\all{Der Text handelt vom Tourismus und seinen Problemen.} — « Le texte traite du tourisme et de ses problèmes. »

\all{Einerseits bringt der Tourismus Geld, andererseits verschmutzt er die Strände.} — « D'un côté le tourisme rapporte de l'argent, de l'autre il pollue les plages. »

\all{Der Autor sagt: „Viele Touristen kommen ans Meer.“ Das zeigt, dass der Strandtourismus sehr beliebt ist.} — « L'auteur dit : „Beaucoup de touristes viennent à la mer.“ Cela montre que le tourisme balnéaire est très apprécié. »

\all{Ich bin der Meinung, dass der Kulturtourismus interessanter ist, weil man viel über das Land lernt.} — « Je suis d'avis que le tourisme culturel est plus intéressant, parce qu'on apprend beaucoup sur le pays. »
\end{exemplebox}

\begin{attention}[Piège : la place du verbe après « Meiner Meinung nach »]
En français, « à mon avis, le tourisme est important » : l'ordre des mots ne change pas. En allemand, \all{Meiner Meinung nach} occupe la \emph{première position} de la phrase : le verbe conjugué doit donc venir \emph{immédiatement après}, en position 2, et le sujet passe derrière le verbe (inversion).

Correct : \all{Meiner Meinung nach \textbf{ist} der Ökotourismus wichtig.}

Faux : \all{Meiner Meinung nach der Ökotourismus ist wichtig.}

Même règle avec \all{einerseits} et \all{andererseits} : \all{Einerseits \textbf{bringt} der Tourismus Geld.} — le verbe reste en position 2. Et n'oubliez pas : dans la subordonnée en \all{weil} ou \all{dass}, le verbe conjugué va \emph{à la fin} : \all{..., weil er die Natur \textbf{schützt}.} Autre piège : ne confondez pas \all{der Autor} (l'auteur) avec \all{das Auto} (la voiture) — faux ami classique !
\end{attention}

\courssub{Commentaire modèle rédigé}

Lisez maintenant le commentaire modèle complet (8 à 10 lignes) sur notre texte support. Repérez, en lisant, les quatre étapes de la méthode.

\begin{textebox}[Kommentar — Modell]
Der Text „Tourismus – Segen oder Problem?“ handelt von den verschiedenen Arten des Tourismus und ihren Folgen. Der Autor zeigt, dass der Tourismus für viele Länder sehr wichtig ist: er bringt Geld und schafft Arbeitsplätze, zum Beispiel in Hotels und Restaurants. Einerseits ist der Tourismus also ein Segen für die Wirtschaft. Andererseits erklärt der Autor, dass große Hotels die Strände verschmutzen und dass die Preise für die Einheimischen steigen. Der Autor sagt: „Die Natur ist in Gefahr.“ Das ist ein großes Problem. Meiner Meinung nach ist der Ökotourismus die beste Lösung, weil er die Natur schützt und den Einheimischen hilft. Ich bin der Meinung, dass auch Kamerun vom Ökotourismus profitieren kann, weil unser Land eine wunderbare Natur hat.
\end{textebox}

\textbf{Glossaire} : die Art, -en (le type, la sorte) ; die Folge, -n (la conséquence) ; der Segen, - (la bénédiction, le bienfait ; presque toujours au singulier) ; die Wirtschaft, -en (l'économie) ; verschmutzen (polluer) ; der Einheimische, -n (l'habitant du pays — nom décliné comme un adjectif : \all{die Einheimischen} au pluriel) ; steigen, ist gestiegen (augmenter, monter) ; in Gefahr (en danger) ; die Lösung, -en (la solution) ; schützen (protéger) ; profitieren von + datif (profiter de).

\courssub{Analyse guidée du modèle}

Décortiquons ensemble ce modèle, étape par étape :

\begin{enumerate}
\item \textbf{Présentation} : \all{Der Text „...“ handelt von den verschiedenen Arten des Tourismus und ihren Folgen.} — titre entre guillemets allemands, formule \all{handelt von} + datif pluriel (\all{von den... Arten... und ihren Folgen}).
\item \textbf{Idée principale} : \all{Der Autor zeigt, dass der Tourismus für viele Länder sehr wichtig ist.} — subordonnée en \all{dass}, verbe \all{ist} à la fin.
\item \textbf{Arguments} : le couple \all{Einerseits... Andererseits...} oppose le pour (argent, emplois : \all{Geld}, \all{Arbeitsplätze}) et le contre (environnement, prix : \all{die Strände verschmutzen}, \all{die Preise steigen}). Notez la \emph{citation} du texte : \all{Der Autor sagt: „Die Natur ist in Gefahr.“} — citer avant de commenter est la marque d'un bon travail.
\item \textbf{Avis personnel} : \all{Meiner Meinung nach ist der Ökotourismus die beste Lösung, weil...} — inversion correcte après \all{Meiner Meinung nach}, justification en \all{weil} avec les verbes \all{schützt} et \all{hilft} en fin de phrase. Le commentaire s'ouvre enfin sur le Cameroun : excellent réflexe pour personnaliser sa copie et montrer que l'on sait réinvestir le texte.
\end{enumerate}

\begin{methode}[Deux règles d'or du commentaire]
\begin{enumerate}
\item \textbf{Citez toujours le texte avant de le commenter} : \all{Der Autor sagt: „...“} ou \all{Im Text steht: „...“}. Une affirmation sans citation n'est pas un commentaire, c'est une opinion en l'air.
\item \textbf{Justifiez toujours votre avis avec \all{weil}} : un avis sans justification ne vaut rien. La structure minimale est : \all{Meiner Meinung nach} + verbe en position 2 + ..., \all{weil} + verbe à la fin.
\end{enumerate}
\end{methode}

\begin{exoresolu}[Reconstituer un commentaire]
Énoncé — Remettez les phrases suivantes dans l'ordre logique d'un commentaire (1 à 4) :
a) \all{Meiner Meinung nach ist der Strandtourismus schlecht für die Umwelt, weil die Hotels das Meer verschmutzen.}
b) \all{Der Text handelt vom Tourismus in Kamerun.}
c) \all{Einerseits bringt der Tourismus Geld, andererseits steigen die Preise.}
d) \all{Der Autor meint, dass der Tourismus Chancen und Probleme hat.}
\tcblower
Solution — L'ordre correct est : \textbf{b) $\rightarrow$ d) $\rightarrow$ c) $\rightarrow$ a)}.
\begin{enumerate}
\item b) Présentation du texte avec \all{handelt von} : on annonce d'abord le thème.
\item d) Idée principale rapportée avec \all{dass} : \all{Der Autor meint, dass...} (verbe \all{hat} à la fin).
\item c) Arguments pour et contre avec \all{einerseits... andererseits...} (verbe en position 2 dans chaque membre de phrase).
\item a) Avis personnel : \all{Meiner Meinung nach ist...} (inversion correcte) + justification \all{weil die Hotels das Meer verschmutzen} (verbe \all{verschmutzen} à la fin).
\end{enumerate}
Ce squelette — présentation, idée, arguments, avis justifié — est le plan universel de tout commentaire de texte en Première.
\end{exoresolu}

\begin{exercice}[À vous de commenter]
Rédigez un mini-commentaire de 5 phrases sur le texte \all{Tourismus – Segen oder Problem?} en suivant le plan : 1) \all{Der Text handelt von...} ; 2) \all{Der Autor zeigt, dass...} ; 3) \all{Einerseits...} ; 4) \all{Andererseits...} ; 5) \all{Meiner Meinung nach..., weil...}. Soulignez dans votre production les connecteurs et les verbes en fin de subordonnée.
\end{exercice}

\begin{corrige}[Exercice — éléments de correction]
Exemple de production attendue : \all{Der Text handelt vom Tourismus und seinen Problemen. Der Autor zeigt, dass der Tourismus Geld und Arbeitsplätze bringt. Einerseits ist er gut für die Wirtschaft. Andererseits verschmutzen die Hotels die Strände und die Preise steigen. Meiner Meinung nach muss der Tourismus die Natur respektieren, weil die Umwelt sehr wichtig ist.} — Vérifiez : verbe en position 2 après \all{Meiner Meinung nach} et \all{Einerseits}/\all{Andererseits} ; verbe à la fin après \all{dass} et \all{weil} ; noms avec majuscule et article correct (\all{die Wirtschaft}, \all{die Strände}, \all{die Umwelt}).
\end{corrige}

\begin{experience}[Entraînement oral : mon mini-commentaire sur le tourisme au Cameroun]
Par groupes de quatre, puis devant la classe, chaque élève présente à l'oral un mini-commentaire en \textbf{5 phrases} sur le tourisme au Cameroun (Kribi, le mont Cameroun, les chutes de la Lobé, le parc de Waza...). Plan imposé :
\begin{enumerate}
\item \all{In Kamerun gibt es viele Sehenswürdigkeiten, zum Beispiel...}
\item \all{Der Tourismus bringt..., weil...}
\item \all{Aber es gibt auch Probleme: ...}
\item \all{Meiner Meinung nach...}
\item \all{Ich bin der Meinung, dass..., weil...}
\end{enumerate}
Les camarades écoutent et cochent sur une grille : présentation ? citation ou exemple ? connecteur \all{einerseits/andererseits} ? avis justifié avec \all{weil} ? Le professeur corrige uniquement la place des verbes et les articles.
\end{experience}

\begin{savaistu}[Le tourisme au Cameroun : « l'Afrique en miniature »]
Le Cameroun est souvent surnommé « l'Afrique en miniature » parce qu'il réunit presque tous les paysages du continent : plages de l'Atlantique (Kribi, Limbé), montagnes (le mont Cameroun, plus de 4 000 mètres), savanes et parcs animaliers (Waza, Bénoué), forêts équatoriales et chutes spectaculaires (les chutes de la Lobé, qui se jettent directement dans l'océan). En allemand, on dirait : \all{Kamerun ist „Afrika im Kleinen“: es gibt Strände, Berge, Savannen und Regenwälder.} — « Le Cameroun est „l'Afrique en miniature“ : il y a des plages, des montagnes, des savanes et des forêts tropicales. » Un atout formidable pour l'écotourisme... à condition de protéger la nature !
\end{savaistu}

\begin{aretenir}[L'essentiel de la section]
Un commentaire de texte suit toujours quatre étapes : \cle{présentation} (\all{Der Text handelt von...}), \cle{idée principale} (\all{Der Autor meint, dass...}), \cle{arguments} (\all{einerseits... andererseits...} + citation \all{Der Autor sagt: „...“}), \cle{avis justifié} (\all{Meiner Meinung nach..., weil...}). Deux règles d'or : citer le texte avant de commenter, et justifier son avis avec \all{weil}. Surveillez la place du verbe : position 2 après \all{Meiner Meinung nach}, position finale après \all{weil} et \all{dass}.
\end{aretenir}

\coursec{Production guidée : Mein Traumurlaub}

Dans la section précédente, tu as appris à commenter un texte sur le tourisme : lire, repérer les idées, résumer, donner ton avis. Il est temps maintenant de passer de la lecture à l'écriture : tu vas rédiger toi-même un court texte sur tes vacances de rêve. Bonne nouvelle : tu possèdes déjà tous les outils nécessaires — le lexique du tourisme, les prépositions de lieu, les subordonnées en \all{weil} et \all{dass}, et le parfait avec \all{sein}. Cette section te montre comment les assembler, étape par étape, pour produire un texte correct et bien structuré, comme on te le demandera à l'examen de fin d'année et plus tard au Bac.

\courssub{La tâche : comprendre la consigne}

\begin{definition}[Die Aufgabe — la tâche]
Tu dois rédiger un paragraphe de \cle{8 à 10 lignes} intitulé \all{„Mein Traumurlaub“} (« Mes vacances de rêve »). Le texte doit présenter : ta destination, le type de tourisme choisi, les activités prévues ou déjà vécues, et ton avis sur les problèmes du tourisme.
\end{definition}

Le plan est \textbf{imposé} : tu ne choisis pas l'ordre des idées. C'est une chance : il te suffit de suivre les quatre étapes, et chaque étape mobilise un outil grammatical de la leçon.

\begin{center}
\begin{tabularx}{\linewidth}{|l|X|X|}
\hline
\rowcolor{popblueL} \textbf{Étape} & \textbf{Contenu} & \textbf{Outil grammatical obligatoire} \\
\hline
1. \all{Ziel} & Où je veux aller & préposition de lieu + accusatif : \all{Ich möchte nach Kamerun reisen.} \\
\hline
2. \all{Grund} & Pourquoi & subordonnée en \all{weil} : \all{..., weil ich das Meer liebe.} \\
\hline
3. \all{Erlebnisse} & Ce que j'ai déjà visité & parfait avec \all{sein}/\all{haben} : \all{Ich bin nach Kribi gefahren.} \\
\hline
4. \all{Meinung} & Mon avis & subordonnée en \all{dass} : \all{Ich denke, dass ...} \\
\hline
\end{tabularx}
\end{center}

\begin{popfigure}
\begin{tikzpicture}[
  bubble/.style={rectangle, rounded corners=6pt, draw=popblue, fill=popblueL, align=center, minimum width=2.6cm, minimum height=1.3cm, font=\small},
  arr/.style={-{Stealth[length=3mm]}, thick, poporange}
]
\node[bubble, align=center] (a) at (0,0){\textbf{1. Ziel}\\Wohin?\\nach / in + Akk.};
\node[bubble, draw=popgreen, fill=popgreenL, align=center] (b) at (3.9,0){\textbf{2. Grund}\\Warum?\\weil-Satz};
\node[bubble, draw=poppurple, fill=poppurpleL, align=center] (c) at (7.8,0){\textbf{3. Erlebnisse}\\Was schon?\\Perfekt mit sein};
\node[bubble, draw=poppink, fill=poppinkL, align=center] (d) at (11.7,0){\textbf{4. Meinung}\\Was denke ich?\\dass-Satz};
\draw[arr] (a) -- (b);
\draw[arr] (b) -- (c);
\draw[arr] (c) -- (d);
\end{tikzpicture}
\legende{Le plan de rédaction en quatre bulles : Ziel $\rightarrow$ Grund $\rightarrow$ Erlebnisse $\rightarrow$ Meinung.}
\end{popfigure}

\courssub{La banque de phrases de départ}

Pour démarrer chaque étape sans stress, utilise ces amorces. Elles garantissent la bonne structure ; tu n'as qu'à les compléter.

\begin{center}
\begin{tabularx}{\linewidth}{|X|X|}
\hline
\rowcolor{popblueL} \textbf{Deutsch} & \textbf{Français} \\
\hline
\all{Ich möchte nach ... reisen, weil ...} & Je voudrais voyager en/à ..., parce que ... \\
\hline
\all{Mein Traumziel ist ..., weil es dort ... gibt.} & Ma destination de rêve est ..., parce qu'il y a ... là-bas. \\
\hline
\all{Letztes Jahr bin ich mit meiner Familie nach ... gefahren.} & L'année dernière, je suis allé(e) à ... avec ma famille. \\
\hline
\all{Wir haben ... besucht / gesehen / fotografiert.} & Nous avons visité / vu / photographié ... \\
\hline
\all{Ich denke, dass ...} & Je pense que ... \\
\hline
\all{Meiner Meinung nach ist der Tourismus ..., weil ...} & À mon avis, le tourisme est ..., parce que ... \\
\hline
\end{tabularx}
\end{center}

\begin{exemplebox}[Deux phrases modèles traduites]
\all{Letztes Jahr bin ich mit meiner Familie nach Kribi gefahren, weil wir das Meer lieben. Ich denke, dass Kribi die schönste Stadt Kameruns ist.}

« L'année dernière, je suis allé à Kribi avec ma famille parce que nous aimons la mer. Je pense que Kribi est la plus belle ville du Cameroun. »

Observe : après \all{weil} et \all{dass}, le verbe conjugué va \textbf{à la fin} de la subordonnée (\all{lieben}, \all{ist}). Au parfait, le verbe de mouvement \all{fahren} prend l'auxiliaire \all{sein} : \all{ich bin gefahren}. Note aussi le génitif \all{Kameruns} : « du Cameroun ».
\end{exemplebox}

\courssub{Un texte modèle à observer}

Avant d'écrire, lis ce texte modèle. Repère les quatre étapes du plan et souligne mentalement les outils grammaticaux : prépositions de lieu, verbes en fin de subordonnée, auxiliaires du parfait.

\begin{textebox}[Mein Traumurlaub — ein Modelltext]
Mein Traumurlaub führt mich nach Deutschland. Ich möchte nach Berlin reisen, weil die Hauptstadt viele Sehenswürdigkeiten hat. Dort möchte ich Kulturtourismus machen: Ich will Museen besuchen und alte Kirchen fotografieren.

Letztes Jahr bin ich mit meiner Familie nach Kribi gefahren, weil wir das Meer lieben. Wir sind an den Strand gegangen und haben viele Fotos gemacht. Das war mein erster Urlaub am Meer, und ich bin sehr glücklich gewesen.

Aber der Tourismus hat auch Probleme. Ich denke, dass zu viele Touristen die Natur zerstören können. Meiner Meinung nach müssen die Hotels die Strände sauber halten, weil die Umwelt wichtig für alle ist. Ich glaube, dass Kamerun ein schönes Land für den Tourismus ist. Mein Traum ist es, eines Tages nach Österreich in die Alpen zu fahren, weil ich die Berge noch nie gesehen habe.
\end{textebox}

\textbf{Glossaire :} \all{die Sehenswürdigkeit, -en} (le site touristique) ; \all{der Kulturtourismus} (le tourisme culturel) ; \all{das Museum, die Museen} (le musée) ; \all{die Kirche, -n} (l'église) ; \all{zerstören} (détruire) ; \all{die Umwelt} (l'environnement) ; \all{sauber halten} (garder propre) ; \all{die Alpen} (les Alpes) ; \all{der Berg, -e} (la montagne).

\textbf{Questions de vérification :}
\begin{enumerate}
\item \all{Wohin möchte der Autor reisen?} (Où l'auteur veut-il voyager ?)
\item \all{Was hat er letztes Jahr gemacht?} (Qu'a-t-il fait l'année dernière ?)
\item \all{Welches Problem des Tourismus nennt er?} (Quel problème du tourisme mentionne-t-il ?)
\end{enumerate}

\begin{corrige}[Questions sur le texte modèle]
\begin{enumerate}
\item \all{Er möchte nach Deutschland, nach Berlin reisen.} — Il veut voyager en Allemagne, à Berlin.
\item \all{Er ist mit seiner Familie nach Kribi gefahren und hat viele Fotos gemacht.} — Il est allé à Kribi avec sa famille et a fait beaucoup de photos.
\item \all{Zu viele Touristen können die Natur zerstören.} — Trop de touristes peuvent détruire la nature.
\end{enumerate}
\end{corrige}

\courssub{La méthode de rédaction en cinq étapes}

\begin{methode}[Comment rédiger « Mein Traumurlaub »]
\begin{enumerate}
\item \textbf{Prépare ton plan.} Note quatre mots-clés : \all{Ziel} (destination), \all{Grund} (raison), \all{Erlebnisse} (souvenirs), \all{Meinung} (avis). Une idée par étape suffit.
\item \textbf{Choisis ton vocabulaire.} Sélectionne au moins \cle{6 mots du lexique} de la leçon : par exemple \all{der Tourismus, der Tourist, das Hotel, die Reise, der Strand, die Sehenswürdigkeit, der Urlaub, das Meer}.
\item \textbf{Rédige en suivant les amorces.} Étape 1 : une phrase avec \all{nach} ou \all{in} + accusatif. Étape 2 : une phrase avec \all{weil} (verbe à la fin). Étape 3 : deux phrases au parfait, dont au moins une avec \all{sein}. Étape 4 : une phrase avec \all{dass}.
\item \textbf{Vérifie le contrat grammatical.} Tu dois avoir : \cle{2 prépositions de lieu}, \cle{1 phrase avec weil}, \cle{1 phrase avec dass}, \cle{2 verbes au parfait}.
\item \textbf{Relis avec la grille d'auto-correction} ci-dessous, puis recopie au propre.
\end{enumerate}
\end{methode}

\begin{aretenir}[Grille d'auto-correction]
Relis ton texte et coche chaque point :
\begin{itemize}
\item \textbf{Articles et majuscules} : chaque nom a-t-il son article correct (\all{der/die/das}) et sa majuscule ?
\item \textbf{Cas après prépositions} : mouvement $\rightarrow$ accusatif (\all{Ich fahre an den Strand}) ; position $\rightarrow$ datif (\all{Ich bin am Strand}).
\item \textbf{Verbe en fin de subordonnée} : après \all{weil} et \all{dass}, le verbe conjugué est-il bien à la dernière place ?
\item \textbf{Auxiliaire sein} : les verbes de mouvement (\all{fahren, gehen, reisen, fliegen, kommen}) prennent-ils \all{sein} au parfait ?
\item \textbf{Place du verbe} : en phrase principale, le verbe conjugué est-il en deuxième position, même si la phrase commence par un complément (\all{Letztes Jahr bin ich ...}) ?
\end{itemize}
\end{aretenir}

\begin{attention}[Le piège de « nach » et des pays avec article]
Après \all{nach}, on utilise les pays \textbf{sans article} et \textbf{sans déclinaison} : \all{nach Kamerun}, \all{nach Deutschland}, \all{nach Gabun}. Mais certains pays prennent un article en allemand, et alors on dit \all{in} + accusatif : \all{in die Schweiz} (la Suisse), \all{in die Türkei} (la Turquie), \all{in den Senegal} (le Sénégal). Pour les villes et les îles, toujours \all{nach} : \all{nach Berlin}, \all{nach Kribi}. Ne dis jamais \emph{« nach die Schweiz »} !
\end{attention}

\courssub{Exercice résolu : corriger un texte d'élève}

\begin{exoresolu}[Finde die fünf Fehler]
Voici le brouillon d'un élève de Première. Il contient \textbf{cinq erreurs}. Trouve-les et corrige-les.

\all{Ich möchte nach die Schweiz reisen, weil ich liebe die Berge. Letztes Jahr ich bin nach Kribi gefahren. Wir sind am Strand gespielt und ich habe im Hotel geschlafen. Ich denke, dass der Tourismus ist wichtig für Kamerun.}
\tcblower
\textbf{Les cinq erreurs et leurs corrections :}
\begin{enumerate}
\item \all{nach die Schweiz} $\rightarrow$ \all{in die Schweiz} : la Suisse prend l'article, donc \all{in} + accusatif, jamais \all{nach}.
\item \all{weil ich liebe die Berge} $\rightarrow$ \all{weil ich die Berge liebe} : après \all{weil}, le verbe conjugué va à la fin de la subordonnée.
\item \all{Letztes Jahr ich bin ... gefahren} $\rightarrow$ \all{Letztes Jahr bin ich ... gefahren} : quand la phrase commence par un complément, le verbe reste en deuxième position et le sujet suit (inversion).
\item \all{Wir sind am Strand gespielt} $\rightarrow$ \all{Wir haben am Strand gespielt} : \all{spielen} n'est pas un verbe de mouvement ni de changement d'état ; il prend l'auxiliaire \all{haben} au parfait.
\item \all{dass der Tourismus ist wichtig} $\rightarrow$ \all{dass der Tourismus wichtig für Kamerun ist} : après \all{dass}, le verbe \all{ist} va à la fin.
\end{enumerate}
\textbf{Texte corrigé :} \all{Ich möchte in die Schweiz reisen, weil ich die Berge liebe. Letztes Jahr bin ich nach Kribi gefahren. Wir haben am Strand gespielt und ich habe im Hotel geschlafen. Ich denke, dass der Tourismus wichtig für Kamerun ist.}

« Je voudrais voyager en Suisse parce que j'aime les montagnes. L'année dernière, je suis allé à Kribi. Nous avons joué sur la plage et j'ai dormi à l'hôtel. Je pense que le tourisme est important pour le Cameroun. »
\end{exoresolu}

\courssub{Production orale : l'interview à Kribi}

\begin{experience}[Ein Interview am Strand von Kribi]
Travaille en binôme. L'un est \textbf{journaliste} pour une radio camerounaise, l'autre est \textbf{touriste} en vacances à Kribi. Le journaliste pose au moins quatre questions avec \all{wohin} (mouvement) ou \all{wo} (position) ; le touriste répond avec des phrases complètes, dont au moins une avec \all{weil} et une avec \all{dass}.

\textbf{Questions possibles du journaliste :}
\begin{itemize}
\item \all{Wohin reisen Sie dieses Jahr?} (Où voyagez-vous cette année ?)
\item \all{Wo wohnen Sie hier in Kribi?} (Où logez-vous ici à Kribi ?)
\item \all{Warum sind Sie nach Kribi gekommen?} (Pourquoi êtes-vous venu à Kribi ?)
\item \all{Was denken Sie über den Tourismus in Kamerun?} (Que pensez-vous du tourisme au Cameroun ?)
\end{itemize}

\textbf{Exemple de réponse :} \all{Ich bin nach Kribi gekommen, weil die Strände hier wunderschön sind. Ich wohne in einem Hotel am Meer. Ich denke, dass Kamerun viele schöne Sehenswürdigkeiten hat.}

« Je suis venu à Kribi parce que les plages sont magnifiques ici. Je loge dans un hôtel au bord de la mer. Je pense que le Cameroun a beaucoup de beaux sites touristiques. »

Ensuite, échangez les rôles. Le professeur écoute et note les verbes bien placés en fin de subordonnée !
\end{experience}

\begin{exercice}[À toi d'écrire]
Rédige maintenant ton propre texte \all{„Mein Traumurlaub“} de 8 à 10 lignes en suivant le plan imposé (Ziel, Grund, Erlebnisse, Meinung). Respecte le contrat : 6 mots du lexique, 2 prépositions de lieu, 1 phrase avec \all{weil}, 1 avec \all{dass}, 2 verbes au parfait. Relis-toi avec la grille d'auto-correction avant de rendre ta copie.
\end{exercice}

\begin{savaistu}[La production écrite au Bac]
Au baccalauréat, la production écrite sur un thème de société comme le tourisme vaut souvent 4 à 5 points. Les correcteurs attribuent environ la moitié des points à la \textbf{structure} : un plan clair, des connecteurs corrects (\all{weil}, \all{dass}, \all{meiner Meinung nach}), et le verbe à la bonne place. Un texte simple mais bien structuré gagne toujours contre un texte ambitieux plein de fautes d'ordre des mots. Entraîne-toi donc à écrire court, clair et correct !
\end{savaistu}

\newpage

\coursec{Activité d'intégration}

\courssub{Objectifs de l'activité}

À l'issue de cette activité d'intégration, tu seras capable de :
\begin{itemize}
\item mobiliser dans un échange oral réel le \cle{vocabulaire du tourisme} étudié dans la leçon (\all{der Tourismus, der Tourist, das Hotel, die Reise, der Strand, die Sehenswürdigkeit, der Arbeitsplatz, die Umwelt}...) ;
\item employer correctement les \cle{prépositions de lieu} avec le bon cas (accusatif pour le mouvement, datif pour la position) ;
\item construire des \cle{subordonnées} avec \all{weil} et \all{dass} en plaçant le verbe conjugué à la fin ;
\item raconter une expérience au \cle{parfait avec sein} (verbes de déplacement et de changement d'état) ;
\item argumenter, réagir à l'avis d'un camarade, puis rédiger une courte synthèse écrite personnelle.
\end{itemize}

\courssub{Situation}

Le lycée bilingue de Yaoundé organise une « Semaine du tourisme ». À cette occasion, le club d'allemand reçoit une journaliste de la radio locale qui prépare un reportage intitulé \all{„Der Tourismus in Kamerun – Chance oder Gefahr?“}. Elle souhaite enregistrer un débat entre les élèves de la classe de Première A. Pour préparer le débat, le club a reçu le texte d'information suivant :

\begin{textebox}[„Kamerun – Afrika im Kleinformat“]
Kamerun wird oft „Afrika im Kleinformat“ genannt, weil man hier fast alles findet: Strände in Kribi und Limbe, Berge wie den Mount Kamerun, den Regenwald im Süden und die Savanne im Norden. Viele Touristen sind schon nach Kamerun gereist und haben die Sehenswürdigkeiten besucht. Der Tourismus schafft Arbeitsplätze: Hotels, Restaurants und Reisebüros brauchen Personal.\par
Aber es gibt auch Probleme. Manche Touristen lassen Müll am Strand, und die Natur leidet. Die Straßen zu vielen schönen Orten sind schlecht, und die Hotels sind oft zu teuer. Die Frage bleibt: Ist der Tourismus für Kamerun eine Chance oder eine Gefahr?
\end{textebox}

\begin{definition}[Glossaire du texte]
\all{der Regenwald, -wälder} : la forêt tropicale ; \all{die Savanne, -n} : la savane ; \all{der Arbeitsplatz, -plätze} : l'emploi ; \all{das Personal} (singulier) : le personnel ; \all{der Müll} (singulier) : les déchets ; \all{leiden (litt, hat gelitten)} : souffrir ; \all{die Gefahr, -en} : le danger ; \all{das Reisebüro, -s} : l'agence de voyages ; \all{schaffen} : créer.
\end{definition}

La classe est divisée en deux groupes : les \all{Befürworter} (les partisans du développement du tourisme au Cameroun) et les \all{Kritiker} (ceux qui dénoncent ses problèmes). Un élève joue le rôle du \all{Moderator} (l'animateur) et la journaliste « enregistre » le débat. À la fin, chaque élève rédige sa conclusion personnelle.

\courssub{Consignes}

\begin{enumerate}
\item \textbf{Lecture et compréhension.} Lis le texte „Kamerun – Afrika im Kleinformat“ à voix haute avec ton groupe, puis réponds par écrit en allemand, par des phrases complètes, à ces trois questions :
\begin{enumerate}
\item \all{Warum wird Kamerun „Afrika im Kleinformat“ genannt?}
\item \all{Nenne zwei Vorteile des Tourismus in Kamerun.}
\item \all{Nenne zwei Probleme des Tourismus in Kamerun.}
\end{enumerate}
\item \textbf{Préparation des arguments.} Dans ton camp (\all{Befürworter} ou \all{Kritiker}), rédige \textbf{trois arguments} en phrases complètes. Chaque argument doit contenir une subordonnée avec \all{weil} ou \all{dass} et au moins un mot du lexique du tourisme. Exemple : \all{Der Tourismus schafft Arbeitsplätze, weil viele Hotels gebaut werden.}
\item \textbf{Souvenirs de voyage.} Ajoute une phrase au parfait avec \all{sein} pour raconter une expérience (réelle ou imaginée) : \all{Ich bin letztes Jahr nach Kribi gefahren.} Attention au choix de l'auxiliaire !
\item \textbf{Le débat oral (15 minutes).} Le \all{Moderator} donne la parole à tour de rôle à chaque camp. Présente tes arguments, écoute ceux des autres et réagis avec les formules de la boîte méthode ci-dessous. Utilise au moins une fois une préposition de lieu : \all{Die Touristen fahren an den Strand. / Der Müll liegt am Strand.}
\item \textbf{Vote et conclusion.} Après le débat, la classe vote : le tourisme au Cameroun est-il plutôt une chance ou une menace ? Le \all{Moderator} annonce le résultat en allemand : \all{Die Mehrheit denkt, dass der Tourismus eine Chance ist.}
\item \textbf{Synthèse écrite individuelle.} Rédige un court texte de six lignes intitulé \all{„Meine Meinung zum Tourismus in Kamerun“}. Il doit obligatoirement contenir : une préposition de lieu (avec le bon cas), une subordonnée avec \all{weil}, une subordonnée avec \all{dass} et un verbe au parfait avec \all{sein}.
\end{enumerate}

\begin{methode}[Réussir un débat en allemand]
\begin{enumerate}
\item \textbf{Donner son avis} : \all{Ich denke, dass ... / Meiner Meinung nach ... / Ich glaube, dass ...}
\item \textbf{Être d'accord} : \all{Da hast du recht. / Ich bin auch dieser Meinung.}
\item \textbf{N'être pas d'accord} : \all{Das stimmt nicht. / Ich sehe das anders, weil ... / Aber hast du daran gedacht, dass ...?}
\item \textbf{Conclure} : \all{Zusammenfassend kann man sagen, dass ...}
\item Parler lentement, en phrases complètes, et vérifier mentalement la place du verbe avant de parler : en deuxième position dans la principale, à la fin dans la subordonnée.
\end{enumerate}
\end{methode}

\courssub{Ressources à mobiliser}

\begin{aretenir}[Lexique utile pour le débat]
\begin{itemize}
\item \all{der Tourismus} (singulier, le tourisme), \all{der Tourist, -en} (le touriste), \all{die Touristin, -nen} (la touriste), \all{die Reise, -n} (le voyage), \all{das Hotel, -s} (l'hôtel), \all{der Strand, -strände} (la plage), \all{die Sehenswürdigkeit, -en} (le site touristique) ;
\item \all{der Arbeitsplatz, -plätze} (l'emploi), \all{die Natur, -en} (la nature), \all{die Umwelt} (l'environnement), \all{der Müll} (les déchets), \all{die Straße, -n} (la route), \all{das Geld, -er} (l'argent), \all{die Chance, -n} (la chance, l'opportunité) ;
\item verbes : \all{reisen} (voyager), \all{besuchen} (visiter), \all{schaffen} (créer), \all{verschmutzen} (polluer), \all{verdienen} (gagner), \all{schützen} (protéger), \all{verbringen (verbrachte, hat verbracht)} (passer du temps).
\end{itemize}
\end{aretenir}

\begin{propriete}[Rappels grammaticaux indispensables]
\begin{enumerate}
\item \textbf{Prépositions de lieu à deux cas} (\all{in, an, auf, vor, hinter, neben, über, unter, zwischen}) : \textbf{accusatif} pour le mouvement (question \all{wohin?} « vers où ? ») : \all{Wir fahren ans Meer.} « Nous allons à la mer. » ; \textbf{datif} pour la position (question \all{wo?} « où ? ») : \all{Wir sind am Meer.} « Nous sommes à la mer. » Contractions fréquentes : \all{an dem = am}, \all{an das = ans}, \all{in dem = im}, \all{in das = ins}.
\item \textbf{Subordonnée avec \all{weil} ou \all{dass}} : le verbe conjugué va à la \textbf{fin} de la subordonnée : \all{Ich glaube, dass der Tourismus wichtig ist.} « Je crois que le tourisme est important. » / \all{Die Natur leidet, weil die Touristen Müll lassen.} « La nature souffre parce que les touristes laissent des déchets. »
\item \textbf{Parfait avec \all{sein}} : pour les verbes de déplacement ou de changement d'état : \all{ich bin gefahren / gereist / gekommen / geblieben / gestiegen}. Tous les autres verbes prennent \all{haben} : \all{wir haben die Sehenswürdigkeiten besucht}.
\end{enumerate}
\end{propriete}

\begin{attention}[Pièges à éviter pendant le débat]
\begin{itemize}
\item Ne dis jamais \all{*Ich habe nach Kribi gefahren} : le verbe \all{fahren} (déplacement) exige l'auxiliaire \all{sein} : \all{Ich \textbf{bin} nach Kribi gefahren.}
\item Après \all{weil} et \all{dass}, le verbe conjugué ne reste pas en deuxième position : on dit \all{..., weil viele Hotels gebaut \textbf{werden}} et non \all{*weil viele Hotels werden gebaut}.
\item \all{zu Hause} et \all{nach Hause} ne prennent jamais d'article ; mais \all{ans Meer, in die Berge, am Strand} exigent l'article contracté.
\item Faux ami : \all{die Chance} signifie « la chance, l'opportunité » (et non « la chance » au sens de la fortune : \all{das Glück}).
\item N'oublie pas la majuscule à tous les noms : \all{der Strand, das Hotel, die Reise}.
\end{itemize}
\end{attention}

\courssub{Proposition de production et corrigé commenté}

\begin{exoresolu}[„Meine Meinung zum Tourismus in Kamerun“ — production modèle]
Rédige une synthèse de six lignes respectant les quatre contraintes : une préposition de lieu, une subordonnée avec \all{weil}, une subordonnée avec \all{dass}, un verbe au parfait avec \all{sein}.
\tcblower
\textbf{Production modèle :}\par
\all{Meiner Meinung nach ist der Tourismus in Kamerun eine große Chance. Ich denke, dass viele Touristen gern nach Kamerun kommen, weil das Land so schön ist. Letztes Jahr bin ich mit meiner Familie nach Kribi gefahren, und wir haben drei Tage am Strand verbracht. Der Tourismus ist wichtig, weil er viele Arbeitsplätze in den Hotels schafft. Aber ich glaube auch, dass wir die Natur und die Strände schützen müssen. Zusammenfassend kann man sagen: Der Tourismus ist eine Chance, wenn wir die Umwelt respektieren.}\par
« À mon avis, le tourisme au Cameroun est une grande chance. Je pense que beaucoup de touristes viennent volontiers au Cameroun parce que le pays est si beau. L'année dernière, je suis allé(e) avec ma famille à Kribi et nous avons passé trois jours à la plage. Le tourisme est important parce qu'il crée beaucoup d'emplois dans les hôtels. Mais je crois aussi que nous devons protéger la nature et les plages. En résumé, on peut dire : le tourisme est une chance si nous respectons l'environnement. »\par
\textbf{Commentaire des choix de langue :}
\begin{itemize}
\item \all{nach Kamerun kommen} : préposition de lieu \all{nach} (pays sans article) pour le mouvement ; \all{am Strand} : \all{an} + datif (\all{an dem} = \all{am}) car il s'agit d'une position ; \all{in den Hotels} : \all{in} + datif pluriel.
\item Subordonnée avec \all{dass} : \all{..., dass viele Touristen gern nach Kamerun kommen} — verbe \all{kommen} à la fin. Subordonnée avec \all{weil} : \all{..., weil das Land so schön ist} — verbe \all{ist} à la fin.
\item Parfait avec \all{sein} : \all{ich bin ... nach Kribi gefahren} (verbe de déplacement). En revanche, \all{wir haben drei Tage verbracht} prend \all{haben}, car \all{verbringen} n'exprime ni déplacement ni changement d'état.
\item Le texte s'ouvre avec \all{Meiner Meinung nach} et se conclut avec \all{Zusammenfassend kann man sagen} : structure claire d'une prise de position argumentée, comme on te le demandera à l'examen.
\end{itemize}
\end{exoresolu}

\courssub{Bilan de l'activité}

Cette activité t'a permis de réinvestir l'ensemble des acquis de la leçon dans une situation de communication authentique :
\begin{itemize}
\item \textbf{Compréhension} : tu as lu et exploité un texte d'information sur le tourisme au Cameroun ;
\item \textbf{Lexique} : tu as employé le vocabulaire du tourisme (\all{der Strand, das Hotel, der Arbeitsplatz, die Sehenswürdigkeit}...) dans tes propres arguments ;
\item \textbf{Grammaire} : tu as choisi le bon cas après les prépositions de lieu et placé correctement le verbe dans les subordonnées avec \all{weil} et \all{dass} ;
\item \textbf{Conjugaison} : tu as raconté une expérience au parfait avec l'auxiliaire \all{sein} ;
\item \textbf{Expression orale et écrite} : tu as défendu un point de vue, réagi à celui des autres, puis rédigé une synthèse personnelle structurée.
\end{itemize}

\begin{experience}[Pour aller plus loin]
En équipe, préparez une petite affiche publicitaire en allemand pour promouvoir un site touristique du Cameroun (Kribi, Limbe, Waza, Foumban...) : un slogan, trois arguments avec \all{weil} ou \all{dass}, et une invitation : \all{Kommen Sie nach Kamerun! Viele Touristen sind schon hier gewesen!} Présentez ensuite votre affiche à la classe en deux minutes.
\end{experience}

\newpage

\coursec{Méthodes et exercices résolus}

\courssub{Savoir-faire 1 : répondre à des questions de compréhension}

\begin{methode}[Comprendre et répondre à des questions sur un texte]
\begin{enumerate}
\item \textbf{Lire deux fois le texte} : une première lecture globale (qui ? quoi ? où ?), une seconde lecture en soulignant les mots-clés du thème (ici : \all{der Tourismus, der Tourist, das Hotel, der Strand, die Sehenswürdigkeit, das Problem}).
\item \textbf{Repérer la ligne} où se trouve la réponse : la question reprend presque toujours un mot du texte (ou son synonyme).
\item \textbf{Répondre par une phrase complète} en allemand : reprendre les mots de la question, jamais un simple « ja » ou « nein ».
\item \textbf{Respecter la place du verbe} : verbe conjugué en 2\textsuperscript{e} position dans la phrase principale.
\item \textbf{Justifier avec une citation courte} entre guillemets allemands „…“ si la consigne le demande (\all{Begründen Sie Ihre Antwort}).
\item \textbf{Se relire} : majuscule aux noms, Umlaute, accord de l'article.
\end{enumerate}
\end{methode}

\begin{exoresolu}[Questions de compréhension sur \all{Tourismus – Segen oder Problem ?}]
Énoncé — Texte : \all{Viele Touristen reisen nach Kamerun. Sie besuchen die Strände von Kribi und den Nationalpark Waza. Der Tourismus bringt Geld und Arbeit, aber er hat auch Probleme : die Hotels sind teuer und die Natur leidet.} Questions : 1) \all{Wohin reisen viele Touristen ?} 2) \all{Was besuchen sie ?} 3) \all{Welche Probleme hat der Tourismus ? Nennen Sie zwei.}
\tcblower
\textbf{Raisonnement} : on repère dans le texte les mots de chaque question. Question 1 : \all{wohin} renvoie à \all{nach Kamerun}. Question 2 : \all{besuchen} est repris mot pour mot. Question 3 : le signal \all{Probleme} introduit la réponse ; il faut choisir deux éléments et les formuler en phrase complète.\par
\textbf{Réponses rédigées} :\par
1) \all{Viele Touristen reisen nach Kamerun.} « Beaucoup de touristes voyagent au Cameroun. »\par
2) \all{Sie besuchen die Strände von Kribi und den Nationalpark Waza.} « Ils visitent les plages de Kribi et le parc national de Waza. »\par
3) \all{Die Hotels sind teuer, und die Natur leidet.} « Les hôtels sont chers et la nature souffre. »\par
\textbf{Conclusion} : chaque réponse est une phrase complète, avec le verbe en 2\textsuperscript{e} position et la majuscule aux noms (\all{Touristen, Strände, Nationalpark, Hotels, Natur}).
\end{exoresolu}

\courssub{Savoir-faire 2 : employer les prépositions de lieu (accusatif/datif)}

\begin{methode}[Choisir le cas après une préposition de lieu]
\begin{enumerate}
\item \textbf{Identifier la préposition} : \all{in, an, auf, nach, zu, vor, hinter, neben, über, unter, zwischen}.
\item \textbf{Poser la bonne question} : \all{wo ?} (où ? = position, repos) $\rightarrow$ \textbf{datif} ; \all{wohin ?} (vers où ? = mouvement) $\rightarrow$ \textbf{accusatif}.
\item \textbf{Cas particuliers à mémoriser} : \all{nach} + pays sans article et villes (\all{nach Kamerun, nach Deutschland, nach Yaoundé}) ; \all{zu} + personne ou lieu précis (\all{zum Markt, zur Schule}) ; \all{in} + pays avec article (\all{in die Schweiz, in der Schweiz}).
\item \textbf{Accorder l'article} : datif masculin \all{dem}, accusatif masculin \all{den} ; contractions utiles : \all{in dem = im, an dem = am, zu dem = zum, zu der = zur, in das = ins, an das = ans}.
\item \textbf{Vérifier le verbe} : verbes de mouvement (\all{gehen, fahren, reisen, fliegen}) appellent \all{wohin} ; verbes de position (\all{sein, liegen, wohnen, bleiben}) appellent \all{wo}.
\end{enumerate}
\end{methode}

\begin{center}
\begin{tabularx}{\linewidth}{|X|X|X|}
\hline
\rowcolor{popblueL} Français & Deutsch & Remarque \\
\hline
Nous allons à la plage. & \all{Wir gehen an den Strand.} & mouvement $\rightarrow$ accusatif \\
\hline
Nous sommes à la plage. & \all{Wir sind am Strand.} & position $\rightarrow$ datif (\all{an dem = am}) \\
\hline
Elle voyage en Suisse. & \all{Sie reist in die Schweiz.} & pays avec article $\rightarrow$ \all{in} + accusatif \\
\hline
Elle habite en Suisse. & \all{Sie wohnt in der Schweiz.} & position $\rightarrow$ datif \\
\hline
\end{tabularx}
\end{center}

\begin{exoresolu}[Prépositions de lieu : thème grammatical]
Énoncé — Traduisez en allemand : 1) « Nous voyageons au Cameroun. » 2) « Les touristes vont à la plage. » 3) « L'hôtel se trouve près du marché. » 4) « Elle pose la valise sur le lit. »
\tcblower
\textbf{Raisonnement phrase par phrase} :\par
1) « au Cameroun » : pays sans article $\rightarrow$ \all{nach} + mouvement (\all{reisen}). \all{Wir reisen nach Kamerun.}\par
2) « à la plage » : mouvement (\all{gehen}) $\rightarrow$ \all{wohin ?} $\rightarrow$ accusatif : \all{an den Strand}. \all{Die Touristen gehen an den Strand.}\par
3) « près du marché » : position (\all{liegen}) $\rightarrow$ \all{wo ?} $\rightarrow$ datif : \all{neben dem Markt}. \all{Das Hotel liegt neben dem Markt.}\par
4) « sur le lit » : mouvement (\all{legen}) $\rightarrow$ accusatif : \all{auf das Bett}. \all{Sie legt den Koffer auf das Bett.}\par
\textbf{Conclusion} : la question \all{wo ? / wohin ?} décide du cas ; c'est le verbe qui donne l'indice.
\end{exoresolu}

\courssub{Savoir-faire 3 : construire des subordonnées avec \all{weil} et \all{dass}}

\begin{methode}[La subordonnée : le verbe à la fin]
\begin{enumerate}
\item \textbf{Identifier la conjonction} : \all{weil} (parce que, cause) ou \all{dass} (que, après \all{ich denke, ich glaube, er sagt}).
\item \textbf{Appliquer la règle d'or} : dans la subordonnée, le \textbf{verbe conjugué va à la dernière place}. \all{Ich bleibe zu Hause, weil ich krank bin.}
\item \textbf{Mettre la virgule} avant \all{weil} et \all{dass}, toujours.
\item \textbf{Si la subordonnée commence la phrase}, elle compte comme la 1\textsuperscript{re} position : le verbe de la principale vient juste après, puis le sujet. \all{Weil das Wetter schön ist, fahren wir ans Meer.}
\item \textbf{Ne pas confondre} \all{weil} (cause) et \all{denn} (car) : après \all{denn}, le verbe reste en 2\textsuperscript{e} position. À l'examen, privilégier \all{weil}.
\end{enumerate}
\end{methode}

\begin{exoresolu}[Transformation et traduction avec \all{weil} et \all{dass}]
Énoncé — 1) Reliez avec \all{weil} : \all{Die Touristen kommen nach Kribi. Die Strände sind schön.} 2) Traduisez : « Je pense que le tourisme est important pour le Cameroun. » 3) Traduisez : « Nous restons à l'hôtel parce qu'il pleut. »
\tcblower
\textbf{Raisonnement} : 1) La deuxième phrase donne la cause $\rightarrow$ \all{weil} + verbe \all{sind} à la fin. 2) « je pense que » $\rightarrow$ \all{ich denke, dass...} avec \all{ist} à la fin ; « pour le Cameroun » = \all{für Kamerun}. 3) « à l'hôtel » : position $\rightarrow$ datif \all{im Hotel} ; « parce qu'il pleut » $\rightarrow$ \all{weil es regnet}, verbe final.\par
\textbf{Réponses rédigées} :\par
1) \all{Die Touristen kommen nach Kribi, weil die Strände schön sind.} « Les touristes viennent à Kribi parce que les plages sont belles. »\par
2) \all{Ich denke, dass der Tourismus für Kamerun wichtig ist.} « Je pense que le tourisme est important pour le Cameroun. »\par
3) \all{Wir bleiben im Hotel, weil es regnet.} « Nous restons à l'hôtel parce qu'il pleut. »\par
\textbf{Conclusion} : virgule + conjonction + verbe conjugué en dernière position : c'est la structure qui rapporte des points.
\end{exoresolu}

\courssub{Savoir-faire 4 : conjuguer au parfait avec \all{sein}}

\begin{methode}[Le Perfekt avec \all{sein}]
\begin{enumerate}
\item \textbf{Choisir l'auxiliaire} : \all{sein} pour les verbes de \textbf{mouvement} (\all{gehen, fahren, fliegen, reisen, kommen, aufstehen}) et de \textbf{changement d'état} (\all{sterben, einschlafen}), plus \all{sein} et \all{bleiben}. Tous les autres prennent \all{haben}.
\item \textbf{Former la structure} : sujet + \all{sein} conjugué au présent + participe II \textbf{à la fin}. \all{Wir sind nach Kamerun geflogen.}
\item \textbf{Former le participe II} : verbes faibles \all{ge...t} (\all{machen $\rightarrow$ gemacht}) ; verbes forts \all{ge...en} (\all{gehen $\rightarrow$ gegangen, fahren $\rightarrow$ gefahren, fliegen $\rightarrow$ geflogen}) ; verbes en \all{-ieren} \textbf{sans} \all{ge-} (\all{studieren $\rightarrow$ studiert, fotografieren $\rightarrow$ fotografiert}).
\item \textbf{Apprendre par cœur} les participes des verbes de voyage : \all{gehen $\rightarrow$ gegangen, fahren $\rightarrow$ gefahren, fliegen $\rightarrow$ geflogen, reisen $\rightarrow$ gereist, kommen $\rightarrow$ gekommen, bleiben $\rightarrow$ geblieben}.
\item \textbf{Conjuguer \all{sein}} : \all{ich bin, du bist, er/sie/es ist, wir sind, ihr seid, sie/Sie sind}.
\end{enumerate}
\end{methode}

\begin{center}
\begin{tabular}{|l|l|l|}
\hline
\rowcolor{popblueL} Infinitif & Auxiliaire & Perfekt (exemple) \\
\hline
\all{fahren} (aller en véhicule) & \all{sein} & \all{Wir sind nach Kribi gefahren.} \\
\hline
\all{fliegen} (voler) & \all{sein} & \all{Sie sind nach Deutschland geflogen.} \\
\hline
\all{bleiben} (rester) & \all{sein} & \all{Er ist eine Woche geblieben.} \\
\hline
\all{besuchen} (visiter) & \all{haben} & \all{Wir haben den Markt besucht.} \\
\hline
\all{fotografieren} (photographier) & \all{haben} & \all{Ich habe die Natur fotografiert.} \\
\hline
\end{tabular}
\end{center}

\begin{exoresolu}[Version et thème au Perfekt]
Énoncé — 1) Mettez au Perfekt : \all{Die Familie fährt ans Meer. Sie bleibt zwei Wochen dort.} 2) Traduisez : « Les touristes sont venus à Yaoundé et ils ont visité le marché. »
\tcblower
\textbf{Raisonnement} : 1) \all{fahren} = mouvement $\rightarrow$ \all{sein} + \all{gefahren} ; \all{bleiben} prend toujours \all{sein} $\rightarrow$ \all{ist geblieben}. Attention : \all{die Familie} est singulier $\rightarrow$ \all{ist}. 2) « sont venus » : mouvement $\rightarrow$ \all{sind gekommen} ; « ont visité » : pas de mouvement $\rightarrow$ \all{haben besucht} ; « le marché » = \all{den Markt} (accusatif après \all{besuchen}).\par
\textbf{Réponses rédigées} :\par
1) \all{Die Familie ist ans Meer gefahren. Sie ist zwei Wochen dort geblieben.} « La famille est allée à la mer. Elle y est restée deux semaines. »\par
2) \all{Die Touristen sind nach Yaoundé gekommen, und sie haben den Markt besucht.} « Les touristes sont venus à Yaoundé et ils ont visité le marché. »\par
\textbf{Conclusion} : mouvement ou changement $\rightarrow$ \all{sein} ; le reste $\rightarrow$ \all{haben}. Le participe II se place toujours à la fin.
\end{exoresolu}

\courssub{Savoir-faire 5 : produire un court texte guidé}

\begin{methode}[Rédiger un court paragraphe : \all{Mein Traumurlaub}]
\begin{enumerate}
\item \textbf{Faire un plan en 3 parties} : destination (\all{Wohin ?}), activités (\all{Was mache ich ?}), justification avec \all{weil}.
\item \textbf{Utiliser le vocabulaire du thème} : \all{die Reise, das Hotel, der Strand, die Sehenswürdigkeit, besuchen, fotografieren, sich erholen}.
\item \textbf{Varier les structures} : au moins une phrase avec \all{weil}, une avec \all{dass}, une au Perfekt si on raconte un voyage passé.
\item \textbf{Viser 6 à 8 phrases simples} plutôt que 3 phrases compliquées et fautives.
\item \textbf{Relire avec la grille} : verbe en 2\textsuperscript{e} position, verbe final dans les subordonnées, majuscules aux noms, articles corrects.
\end{enumerate}
\end{methode}

\begin{exoresolu}[Production écrite guidée]
Énoncé — \all{Schreiben Sie einen kurzen Text (6–8 Sätze) : „Mein Traumurlaub“. Benutzen Sie : weil, dass, eine Präposition mit Dativ, eine mit Akkusativ.}
\tcblower
\textbf{Raisonnement} : plan en 3 temps : destination (Kribi), activités (plage, poisson, photos), justification (\all{weil}). On place volontairement : un datif de position (\all{im Hotel}), un accusatif de mouvement (\all{an den Strand}), une subordonnée \all{weil} et une \all{dass}.\par
\textbf{Texte modèle} :\par
\all{Mein Traumurlaub ist eine Reise nach Kribi. Ich fahre mit meiner Familie ans Meer. Wir wohnen in einem kleinen Hotel am Strand. Jeden Morgen gehen wir an den Strand und schwimmen im Meer. Am Nachmittag besuchen wir die Wasserfälle von Lobé, weil sie sehr schön sind. Ich denke, dass Kribi der schönste Ort in Kamerun ist. Dort esse ich frischen Fisch und fotografiere die Natur. Dieser Urlaub ist ein Traum!}\par
« Mes vacances de rêve, c'est un voyage à Kribi. Je pars avec ma famille à la mer. Nous logeons dans un petit hôtel au bord de la plage. Chaque matin, nous allons à la plage et nous nageons dans la mer. L'après-midi, nous visitons les chutes de Lobé parce qu'elles sont très belles. Je pense que Kribi est le plus bel endroit du Cameroun. Là-bas, je mange du poisson frais et je photographie la nature. Ces vacances sont un rêve ! »\par
\textbf{Conclusion} : le texte contient les 4 structures exigées (\all{weil}, \all{dass}, datif \all{im Hotel / am Strand}, accusatif \all{an den Strand / ans Meer}) et reste simple et correct.
\end{exoresolu}

\begin{attention}[Les 5 erreurs qui coûtent des points à l'examen]
\begin{enumerate}
\item \textbf{Oublier la majuscule aux noms} : \all{das hotel, der strand} sont faux ; écrire \all{das Hotel, der Strand, der Tourismus}. C'est la faute la plus sanctionnée.
\item \textbf{Verbe mal placé dans la subordonnée} : jamais \all{weil die Strände sind schön} ; toujours \all{weil die Strände schön \textbf{sind}}. Le verbe conjugué va à la fin après \all{weil} et \all{dass}.
\item \textbf{Confondre datif et accusatif après les prépositions de lieu} : \all{wir gehen im Strand} est faux ; mouvement $\rightarrow$ \all{wir gehen \textbf{an den} Strand} ; position $\rightarrow$ \all{wir sind \textbf{am} Strand}.
\item \textbf{Mauvais auxiliaire au Perfekt} : \all{wir haben nach Kribi gefahren} est faux ; les verbes de mouvement prennent \all{sein} : \all{wir \textbf{sind} nach Kribi gefahren}.
\item \textbf{Calquer le français} : « visiter » ne se traduit pas par \all{visitieren} (faux ami, registre médical) ; on dit \all{besuchen} ou \all{besichtigen} : \all{Wir besuchen die Sehenswürdigkeiten.} De même, « passer les vacances » = \all{die Ferien verbringen}, jamais \all{passieren} (qui signifie « se passer, arriver »).
\end{enumerate}
\end{attention}

\newpage

\coursec{Exercices}

\courssub{Je vérifie mes connaissances}

\begin{exercice}[QCM : le tourisme et ses formes]
Entourez la bonne réponse.
\begin{enumerate}
\item \all{Der Ökotourismus} est un tourisme qui :
\begin{enumerate}
\item détruit la nature
\item respecte la nature et la culture locale
\item coûte toujours très cher
\end{enumerate}
\item \all{Der Massentourismus} signifie :
\begin{enumerate}
\item le tourisme de très nombreux visiteurs
\item le tourisme des sportifs
\item le tourisme d'affaires
\end{enumerate}
\item \all{Die Sehenswürdigkeit} est :
\begin{enumerate}
\item un hôtel de luxe
\item un monument ou un lieu à visiter
\item un moyen de transport
\end{enumerate}
\item Après \all{weil}, le verbe conjugué se place :
\begin{enumerate}
\item en deuxième position
\item à la fin de la subordonnée
\item juste après \all{weil}
\end{enumerate}
\end{enumerate}
\end{exercice}

\begin{exercice}[Vrai ou faux ? Justifiez]
Répondez par \all{richtig} ou \all{falsch} et justifiez chaque réponse par une courte phrase en allemand.
\begin{enumerate}
\item \all{Der Tourismus bringt nur Vorteile für ein Land.}
\item \all{Nach} s'emploie devant les noms de pays sans article, comme \all{Kamerun}.
\item \all{Ich bin nach Yaoundé gereist} est une phrase correcte au parfait.
\item \all{Dass} introduit une subordonnée où le verbe est en deuxième position.
\end{enumerate}
\end{exercice}

\begin{exercice}[Vocabulaire : trouvez le mot]
Complétez chaque phrase avec le mot qui convient : \all{der Strand}, \all{die Reise}, \all{das Hotel}, \all{die Sehenswürdigkeit}, \all{der Tourist}, \all{der Tourismus}.
\begin{enumerate}
\item Wir liegen den ganzen Tag am \ldots{} und schwimmen im Meer.
\item Das \ldots{} in Douala ist sehr modern, es hat 200 Zimmer.
\item Die \ldots{} nach Österreich dauert zwei Wochen.
\item Der Eiffelturm ist eine berühmte \ldots{} in Paris.
\item Viele \ldots{} besuchen jedes Jahr den Berg Kamerun.
\item Der \ldots{} ist für die Wirtschaft des Landes sehr wichtig.
\end{enumerate}
\end{exercice}

\begin{exercice}[Conjugaison : le parfait avec \all{sein}]
Mettez les verbes entre parenthèses au parfait. Attention : auxiliaire \all{sein} ou \all{haben} ?
\begin{enumerate}
\item Wir \ldots{} nach Deutschland \ldots{} (reisen).
\item Die Touristen \ldots{} ans Meer \ldots{} (fahren).
\item Ich \ldots{} lange in Kribi \ldots{} (bleiben).
\item Er \ldots{} ein Hotel \ldots{} (besuchen).
\item Ihr \ldots{} nach Yaoundé \ldots{} (fliegen).
\item Sie (elle) \ldots{} viele Fotos \ldots{} (machen).
\end{enumerate}
\end{exercice}

\courssub{Je m'entraîne}

\begin{exercice}[Prépositions de lieu : le bon cas]
Complétez avec la bonne préposition et le bon article (accusatif = mouvement, datif = position).
\begin{enumerate}
\item Wir fliegen \ldots{} \ldots{} Schweiz. (in die / in der)
\item Die Familie wohnt jetzt \ldots{} \ldots{} Schweiz. (in die / in der)
\item Die Touristen gehen \ldots{} \ldots{} Meer. (ans / am)
\item Sie liegen schon \ldots{} \ldots{} Meer. (ans / am)
\item Er fährt \ldots{} \ldots{} Hotel. (ins / im)
\item Er schläft \ldots{} \ldots{} Hotel. (ins / im)
\item Nächstes Jahr reisen wir \ldots{} Kamerun. (nach / in)
\end{enumerate}
\end{exercice}

\begin{exercice}[Transformation : \all{weil} et \all{dass}]
Reliez les deux phrases : d'abord avec \all{weil}, puis avec \all{dass} (commencez par \all{Ich denke, \ldots}).
\begin{enumerate}
\item Die Touristen kommen. Die Preise steigen.
\item Viele Leute fliegen nach Mallorca. Die Flüge sind billig.
\item Wir bleiben zu Hause. Wir haben kein Geld.
\item Der Ökotourismus ist wichtig. Er schützt die Natur.
\end{enumerate}
Exemple : \all{Die Touristen kommen. Die Preise steigen.} $\rightarrow$ \all{Die Preise steigen, weil die Touristen kommen.} / \all{Ich denke, dass die Preise steigen.}
\end{exercice}

\begin{textebox}[Meine Reise nach Kamerun]
Letzten Sommer bin ich mit meiner Familie nach Kamerun \ldots{}. Wir haben zuerst Douala besucht, \ldots{} dort der Hafen sehr interessant ist. Dann sind wir nach Kribi gefahren. Wir sind eine Woche in einem kleinen \ldots{} am Meer \ldots{}. Jeden Tag sind wir an den \ldots{} gegangen. Kamerun hat viele \ldots{} : den Berg Kamerun, die Wasserfälle von Ekom und die Lobé-Wasserfälle. Es war ein wunderbarer Urlaub !
\end{textebox}

\begin{exercice}[Texte à trous : un voyage au Cameroun]
Complétez le texte ci-dessus avec les mots suivants : \all{gereist}, \all{Sehenswürdigkeiten}, \all{Strand}, \all{weil}, \all{Hotel}, \all{geblieben}.
\end{exercice}

\begin{exercice}[Appariements : mots et définitions]
Reliez chaque mot allemand à sa définition française.
\begin{center}
\begin{tabularx}{\linewidth}{|X|X|}
\hline
\rowcolor{popblueL} \textbf{Deutsch} & \textbf{Français} \\
\hline
1. \all{der Massentourismus} & a. voyage organisé pour les affaires \\
\hline
2. \all{der Ökotourismus} & b. affluence de très nombreux visiteurs \\
\hline
3. \all{die Geschäftsreise} & c. tourisme respectueux de la nature \\
\hline
4. \all{der Badeurlaub} & d. avantage \\
\hline
5. \all{der Vorteil} & e. vacances à la plage \\
\hline
6. \all{der Nachteil} & f. inconvénient \\
\hline
\end{tabularx}
\end{center}
\end{exercice}

\courssub{Je me prépare au Bac}

\begin{textebox}[Ökotourismus in Österreich]
Immer mehr Touristen kommen nach Österreich, aber sie wollen nicht mehr nur Ski fahren oder einkaufen. Viele Besucher suchen heute einen Urlaub in der Natur : sie wandern in den Alpen, schlafen auf einem Bauernhof und essen Produkte aus der Region. Das nennt man Ökotourismus. Letztes Jahr sind schon viele Familien in die Alpen gefahren, und einige Jugendliche haben auf einem Bauernhof gearbeitet.

Diese Art von Tourismus hat viele Vorteile. Die Touristen respektieren die Umwelt, weil sie wenig Müll produzieren und die Natur nicht zerstören. Außerdem verdienen die Bauern und die kleinen Dörfer Geld, und die Jugendlichen bleiben in ihrer Region, weil sie dort Arbeit finden.

Aber es gibt auch Probleme. Manche Hotels sind weit von den Städten, und die Touristen müssen mit dem Auto fahren. Das ist schlecht für die Umwelt. Außerdem ist ein Urlaub auf dem Bauernhof oft teurer als eine Reise ans Meer. Trotzdem glauben viele Experten, dass der Ökotourismus die Zukunft des Reisens ist.

(Environ 160 mots)
\end{textebox}

\begin{exercice}[Compréhension de l'écrit : \all{Ökotourismus in Österreich}]
Lisez le texte ci-dessus, puis répondez aux questions.

\textbf{A. Vrai ou faux ?} (Justifiez chaque réponse par une citation du texte.)
\begin{enumerate}
\item Die Touristen in Österreich wollen heute nur Ski fahren.
\item Die Bauern verdienen Geld mit dem Ökotourismus.
\item Der Ökotourismus hat keine Nachteile.
\end{enumerate}

\textbf{B. Questions de compréhension.} (Répondez en allemand, par des phrases complètes.)
\begin{enumerate}
\item Was machen die Ökotouristen in Österreich ? (2 activités)
\item Warum bleiben die Jugendlichen in ihrer Region ?
\item Welches Problem hat der Urlaub auf dem Bauernhof ?
\item Was glauben viele Experten ?
\end{enumerate}

\textbf{C. Question personnelle.}
\begin{enumerate}
\item Möchten Sie auch einen Urlaub auf dem Bauernhof machen ? Warum (nicht) ? (2--3 Sätze)
\end{enumerate}

\textbf{D. Questions de langue.}
\begin{enumerate}
\item Trouvez dans le texte deux verbes conjugués au parfait et donnez leur infinitif.
\item Relevez une subordonnée avec \all{weil} et soulignez le verbe conjugué.
\item Mettez au parfait : \all{Die Touristen respektieren die Umwelt.}
\end{enumerate}
\end{exercice}

\begin{exercice}[Thème et version]
\textbf{A. Thème.} Traduisez en allemand les phrases suivantes.
\begin{enumerate}
\item Les touristes vont à la plage parce qu'ils aiment le soleil.
\item Je pense que l'écotourisme protège la nature.
\item Nous sommes allés en Suisse l'année dernière.
\item Beaucoup de touristes sont venus au Cameroun parce que le pays est très beau.
\item Elle est restée une semaine à l'hôtel, puis elle est rentrée à Yaoundé.
\end{enumerate}

\textbf{B. Version.} Traduisez en français le texte suivant.

\begin{textebox}[Massentourismus auf Mallorca]
Jeden Sommer kommen Millionen Touristen nach Mallorca. Die Strände sind voll, die Hotels sind komplett reserviert, und die Preise steigen. Viele Einwohner sind nicht zufrieden, weil das Leben auf der Insel immer teurer wird. Der Tourismus bringt viel Geld, aber er bringt auch Probleme : zu viel Müll, zu viele Autos und zu wenig Wasser. Die Regierung sucht jetzt neue Ideen für einen besseren Tourismus.
\end{textebox}
\end{exercice}

\begin{exercice}[Production écrite : \all{Mein Traumurlaub}]
Rédigez un court texte en allemand de \textbf{60 à 80 mots} sur le sujet suivant :

\textbf{Sujet :} \all{Welche Art von Tourismus mögen Sie ? Warum ? Beschreiben Sie Ihren Traumurlaub.}

Consignes :
\begin{itemize}
\item Présentez le type de tourisme que vous préférez (plage, nature, ville, écotourisme\ldots).
\item Donnez au moins deux raisons avec \all{weil} ou \all{dass}.
\item Utilisez au moins une phrase au parfait (un voyage que vous avez déjà fait).
\item Employez le vocabulaire de la leçon : \all{die Reise}, \all{der Strand}, \all{das Hotel}, \all{die Sehenswürdigkeit}, \all{besuchen}, \all{reisen}\ldots
\end{itemize}
\end{exercice}

\begin{methode}[Réussir la production écrite : \all{Mein Traumurlaub}]
\begin{enumerate}
\item \textbf{Introduction} (1 phrase) : annoncez votre type de tourisme préféré. \all{Ich mag den Ökotourismus sehr.}
\item \textbf{Raisons} (2--3 phrases) : utilisez \all{weil} et \all{dass} avec le verbe à la fin. \all{Ich mag ihn, weil er die Natur schützt. Ich denke, dass dieser Tourismus wichtig für die Zukunft ist.}
\item \textbf{Expérience au parfait} (1--2 phrases) : \all{Letztes Jahr bin ich mit meiner Familie nach Kribi gefahren. Wir haben ein kleines Hotel am Strand besucht.}
\item \textbf{Conclusion} (1 phrase) : votre rêve. \all{Mein Traumurlaub ist eine Reise in die Alpen.}
\item \textbf{Vérification} : relisez en contrôlant la place du verbe, les majuscules des noms et le choix de l'auxiliaire au parfait.
\end{enumerate}
\end{methode}

\begin{exemplebox}[Modèle de production : \all{Mein Traumurlaub}]
\all{Ich mag den Ökotourismus sehr, weil er die Natur und die Kultur respektiert. Ich denke, dass dieser Tourismus die Zukunft ist. Letztes Jahr bin ich mit meiner Familie nach Kribi gefahren. Wir sind eine Woche in einem kleinen Hotel geblieben, und jeden Tag sind wir an den Strand gegangen. Es war wunderbar ! Mein Traumurlaub ist eine Reise nach Österreich : ich möchte in den Alpen wandern und die Sehenswürdigkeiten von Wien besuchen.}

« J'aime beaucoup l'écotourisme, parce qu'il respecte la nature et la culture. Je pense que ce tourisme est l'avenir. L'année dernière, je suis allé avec ma famille à Kribi. Nous sommes restés une semaine dans un petit hôtel, et chaque jour nous sommes allés à la plage. C'était merveilleux ! Mes vacances de rêve sont un voyage en Autriche : je voudrais faire de la randonnée dans les Alpes et visiter les monuments de Vienne. » (75 mots)
\end{exemplebox}

\newpage

\coursec{Corrigés des exercices}

\begin{corrige}[Exercice 1 — QCM : le tourisme et ses formes]
\begin{enumerate}
\item \textbf{b} : \all{Der Ökotourismus respektiert die Natur und die Kultur.} („L'écotourisme respecte la nature et la culture.“) C'est justement sa définition : un tourisme responsable (\cle{der Ökotourismus}).
\item \textbf{a} : \all{Massentourismus} = tourisme de masse. Famille de mots : \all{die Masse} (la masse) + \all{der Tourismus} (\cle{der Massentourismus}).
\item \textbf{b} : \all{die Sehenswürdigkeit} = le monument, le site à visiter. Étymologie : \all{sehen} (voir) + \all{würdig} (digne) : littéralement „digne d'être vu“ (\cle{die Sehenswürdigkeit}, \cle{die Sehenswürdigkeiten}).
\item \textbf{b} : dans la subordonnée introduite par \all{weil}, le verbe conjugué va \textbf{à la fin} : \all{Ich bleibe zu Hause, weil ich kein Geld habe.} („Je reste à la maison parce que je n'ai pas d'argent.“)
\end{enumerate}
\end{corrige}

\begin{corrige}[Exercice 2 — Vrai ou faux ? Justifiez]
\begin{enumerate}
\item \all{Falsch. Der Tourismus bringt Geld, aber auch Probleme wie Müll und hohe Preise.} („Faux. Le tourisme rapporte de l'argent, mais aussi des problèmes comme les déchets et les prix élevés.“)
\item \all{Richtig.} \all{Nach} + pays sans article : \all{Wir reisen nach Kamerun.} („Nous voyageons au Cameroun.“) Mais attention : \all{in die Schweiz}, \all{in den Senegal}, car ces pays ont un article (\cle{Präpositionen mit Länder}).
\item \all{Richtig.} \all{Reisen} est un verbe de déplacement : parfait avec \all{sein} : \all{ich bin gereist.} („Je suis allé / j'ai voyagé.“)
\item \all{Falsch. Nach „dass“ steht das Verb am Ende : Ich denke, dass der Tourismus wichtig ist.} („Faux. Après dass, le verbe est à la fin : Je pense que le tourisme est important.“)
\end{enumerate}
\textbf{Points de vigilance :}
\begin{itemize}
\item Ne pas écrire \all{ich habe gereist} (auxiliaire \cle{haben} incorrect pour \cle{reisen}).
\item Ne jamais placer le verbe en deuxième position après \all{weil} ou \all{dass} (verbe en \cle{position finale}).
\end{itemize}
\end{corrige}

\begin{corrige}[Exercice 3 — Vocabulaire : trouvez le mot]
\begin{enumerate}
\item \all{Wir liegen den ganzen Tag am Strand und schwimmen im Meer.} (\all{am Strand} = \all{an dem Strand}, \cle{Dativ} : position ; „Nous restons toute la journée à la plage et nous nageons dans la mer.“)
\item \all{Das Hotel in Douala ist sehr modern, es hat 200 Zimmer.}
\item \all{Die Reise nach Österreich dauert zwei Wochen.} („Le voyage en Autriche dure deux semaines.“)
\item \all{Der Eiffelturm ist eine berühmte Sehenswürdigkeit in Paris.}
\item \all{Viele Touristen besuchen jedes Jahr den Berg Kamerun.} (Pluriel de \all{der Tourist} : \cle{die Touristen} — nom masculin \textbf{régulier}, pas de n-déclinaison : \all{der Tourist, des Touristen, die Touristen}).
\item \all{Der Tourismus ist für die Wirtschaft des Landes sehr wichtig.} („Le tourisme est très important pour l'économie du pays.“)
\end{enumerate}
\end{corrige}

\begin{corrige}[Exercice 4 — Conjugaison : le parfait avec \all{sein}]
\begin{enumerate}
\item \all{Wir sind nach Deutschland gereist.} (déplacement $\rightarrow$ \cle{sein})
\item \all{Die Touristen sind ans Meer gefahren.} (déplacement : \all{ans} = \all{an das}, \cle{Akkusativ})
\item \all{Ich bin lange in Kribi geblieben.} (\all{bleiben} prend toujours \cle{sein} — changement d'état / permanence)
\item \all{Er hat ein Hotel besucht.} (verbe transitif, pas de déplacement $\rightarrow$ \cle{haben})
\item \all{Ihr seid nach Yaoundé geflogen.} (déplacement $\rightarrow$ \cle{sein} ; participe irrégulier : \all{fliegen -- flog -- ist geflogen})
\item \all{Sie hat viele Fotos gemacht.} (\cle{haben} ; „Elle a fait beaucoup de photos.“)
\end{enumerate}
\textbf{Rappel :} \cle{sein} + \cle{Partizip II} pour les verbes de déplacement (\all{reisen, fahren, fliegen, gehen, kommen, laufen, schwimmen}) ou de changement d'état (\all{bleiben, werden, sterben, aufwachen, einschlafen}) ; \cle{haben} pour les autres. Le participe passé se place \textbf{à la fin} de la phrase.
\end{corrige}

\begin{corrige}[Exercice 5 — Prépositions de lieu : le bon cas]
\begin{enumerate}
\item \all{Wir fliegen in die Schweiz.} (mouvement $\rightarrow$ \cle{Akkusativ} : \all{in die})
\item \all{Die Familie wohnt jetzt in der Schweiz.} (position $\rightarrow$ \cle{Dativ} : \all{in der})
\item \all{Die Touristen gehen ans Meer.} (mouvement : \all{ans} = \all{an das}, \cle{Akkusativ})
\item \all{Sie liegen schon am Meer.} (position : \all{am} = \all{an dem}, \cle{Dativ})
\item \all{Er fährt ins Hotel.} (mouvement : \all{ins} = \all{in das}, \cle{Akkusativ})
\item \all{Er schläft im Hotel.} (position : \all{im} = \all{in dem}, \cle{Dativ})
\item \all{Nächstes Jahr reisen wir nach Kamerun.} (pays sans article $\rightarrow$ \cle{nach} + \cle{Dativ})
\end{enumerate}
\textbf{Question-test :} \all{wohin ?} (où aller ?) $\rightarrow$ \cle{Akkusativ} ; \all{wo ?} (où être ?) $\rightarrow$ \cle{Dativ}. 
\textbf{Point de vigilance :} le français utilise „en/à“ dans les deux cas ; l'allemand distingue \textbf{toujours} mouvement (\cle{wohin}) et position (\cle{wo}).
\end{corrige}

\begin{corrige}[Exercice 6 — Transformation : \all{weil} et \all{dass}]
\begin{enumerate}
\item \all{Die Preise steigen, weil die Touristen kommen.} / \all{Ich denke, dass die Preise steigen.}
\item \all{Viele Leute fliegen nach Mallorca, weil die Flüge billig sind.} / \all{Ich denke, dass viele Leute nach Mallorca fliegen.}
\item \all{Wir bleiben zu Hause, weil wir kein Geld haben.} / \all{Ich denke, dass wir zu Hause bleiben.}
\item \all{Der Ökotourismus ist wichtig, weil er die Natur schützt.} / \all{Ich denke, dass der Ökotourismus die Natur schützt.}
\end{enumerate}
\textbf{Point de vigilance :} après \cle{weil} et \cle{dass}, le verbe conjugué va \textbf{à la fin} de la subordonnée ; dans la principale qui suit une subordonnée en tête, le verbe se place en \textbf{première position} (\all{Weil ich müde bin, bleibe ich zu Hause.}).
\end{corrige}

\begin{corrige}[Exercice 7 — Texte à trous : un voyage au Cameroun]
\all{Letzten Sommer bin ich mit meiner Familie nach Kamerun \textbf{gereist}. Wir haben zuerst Douala besucht, \textbf{weil} dort der Hafen sehr interessant ist. Dann sind wir nach Kribi \textbf{gefahren}. Wir sind eine Woche in einem kleinen \textbf{Hotel} am Meer \textbf{geblieben}. Jeden Tag sind wir an den \textbf{Strand} gegangen. Kamerun hat viele \textbf{Sehenswürdigkeiten} : den Berg Kamerun, die \textbf{Ekom-Wasserfälle} und die \textbf{Lobé-Wasserfälle}. Es war ein wunderbarer Urlaub !}

Traduction : « L'été dernier, je suis allé avec ma famille au Cameroun. Nous avons d'abord visité Douala, parce que le port y est très intéressant. Ensuite, nous sommes allés à Kribi. Nous sommes restés une semaine dans un petit hôtel au bord de la mer. Chaque jour, nous sommes allés à la plage. Le Cameroun a beaucoup de sites touristiques : le mont Cameroun, les chutes d'Ekom et les chutes de la Lobé. C'était des vacances merveilleuses ! »

\textbf{Justifications :} \all{gereist}, \all{gefahren}, \all{geblieben}, \all{gegangen} avec l'auxiliaire \cle{sein} (verbes de déplacement / \all{bleiben}) ; \all{weil} introduit la cause avec le verbe en finale (\all{ist}) ; \all{Sehenswürdigkeiten} au pluriel après \all{viele} ; noms composés \all{Ekom-Wasserfälle}, \all{Lobé-Wasserfälle} (majuscule, trait d'union).
\end{corrige}

\begin{corrige}[Exercice 8 — Appariements : mots et définitions]
1 -- b ; 2 -- c ; 3 -- a ; 4 -- e ; 5 -- d ; 6 -- f.

\textbf{Remarques :} \all{der Massentourismus} = \all{die Masse} (la masse) + \all{der Tourismus} (mot composé) ; \all{die Geschäftsreise} = \all{das Geschäft} (l'affaire, le commerce) + \all{die Reise} (le voyage) ; \all{der Vorteil} / \all{der Nachteil} forment un couple d'antonymes (\all{vor-} = avant, \all{nach-} = après/arrière).
\end{corrige}

\begin{corrige}[Exercice 9 — Compréhension de l'écrit : \all{Ökotourismus in Österreich}]
\textbf{Barème indicatif (sur 20) :} A : 3 pts (1 pt par réponse justifiée) ; B : 8 pts (2 pts par question) ; C : 3 pts ; D : 6 pts (2 pts par question).

\textbf{A. Vrai ou faux ?}
\begin{enumerate}
\item \all{Falsch : „Sie wollen nicht mehr nur Ski fahren oder einkaufen.“}
\item \all{Richtig : „Außerdem verdienen die Bauern und die kleinen Dörfer Geld.“}
\item \all{Falsch : „Aber es gibt auch Probleme.“}
\end{enumerate}

\textbf{B. Questions de compréhension.}
\begin{enumerate}
\item \all{Sie wandern in den Alpen, schlafen auf einem Bauernhof und essen Produkte aus der Region.} („Ils font de la randonnée dans les Alpes, dorment dans une ferme et mangent des produits de la région.“)
\item \all{Sie bleiben in ihrer Region, weil sie dort Arbeit finden.} („Ils restent dans leur région parce qu'ils y trouvent du travail.“)
\item \all{Er ist oft teurer als eine Reise ans Meer, und die Touristen müssen mit dem Auto fahren.} („Il est souvent plus cher qu'un voyage à la mer, et les touristes doivent prendre la voiture.“)
\item \all{Viele Experten glauben, dass der Ökotourismus die Zukunft des Reisens ist.} („Beaucoup d'experts pensent que l'écotourisme est l'avenir du voyage.“)
\end{enumerate}

\textbf{C. Question personnelle (exemple de réponse).}\\
\all{Ja, ich möchte auch einen Urlaub auf dem Bauernhof machen, weil ich die Natur liebe. Ich denke, dass dieser Urlaub sehr ruhig und interessant ist.} („Oui, je voudrais aussi passer des vacances à la ferme, parce que j'aime la nature. Je pense que ces vacances sont très calmes et intéressantes.“)

\textbf{D. Questions de langue.}
\begin{enumerate}
\item \all{sind gefahren} (infinitif : \all{fahren}) ; \all{haben gearbeitet} (infinitif : \all{arbeiten}).
\item \all{Die Touristen respektieren die Umwelt, weil sie wenig Müll \underline{produzieren}.} (verbe conjugué en finale : \all{produzieren})
\item \all{Die Touristen haben die Umwelt respektiert.} (verbe transitif $\rightarrow$ auxiliaire \cle{haben})
\end{enumerate}

\textbf{Points de vigilance :}
\begin{itemize}
\item Répondre par des phrases complètes (verbe en deuxième position dans la principale).
\item Citer le texte exactement (guillemets) pour justifier Vrai/Faux.
\item Ne pas confondre \all{das Dorf} (le village) et \all{die Stadt} (la ville).
\end{itemize}
\end{corrige}

\begin{corrige}[Exercice 10 — Thème et version]
\textbf{Barème indicatif (sur 20) :} thème : 10 pts (2 pts par phrase) ; version : 10 pts (sens global 4 pts, exactitude lexicale et grammaticale 6 pts).

\textbf{A. Thème.}
\begin{enumerate}
\item \all{Die Touristen gehen an den Strand, weil sie die Sonne lieben.} (mouvement : \all{an den Strand}, \cle{Akkusativ})
\item \all{Ich denke, dass der Ökotourismus die Natur schützt.} (\cle{dass} + verbe à la fin)
\item \all{Letztes Jahr sind wir in die Schweiz gefahren (gereist).} (\all{die Schweiz} a un article $\rightarrow$ \all{in die}, \cle{Akkusativ}, pas \all{nach} !)
\item \all{Viele Touristen sind nach Kamerun gekommen, weil das Land sehr schön ist.} (\all{kommen} $\rightarrow$ \cle{sein} ; \all{nach} + pays sans article)
\item \all{Sie ist eine Woche im Hotel geblieben, dann ist sie nach Yaoundé zurückgefahren.} (\all{bleiben} et \all{zurückfahren} avec \cle{sein} ; \all{im} = \all{in dem}, \cle{Dativ})
\end{enumerate}

\textbf{B. Version.}\\
Chaque été, des millions de touristes viennent à Majorque. Les plages sont pleines, les hôtels sont complets et les prix augmentent. Beaucoup d'habitants ne sont pas satisfaits, parce que la vie sur l'île devient de plus en plus chère. Le tourisme rapporte beaucoup d'argent, mais il apporte aussi des problèmes : trop de déchets, trop de voitures et trop peu d'eau. Le gouvernement cherche maintenant de nouvelles idées pour un meilleur tourisme.

\textbf{Proposition de traduction :}
\all{Jeden Sommer kommen Millionen Touristen nach Mallorca. Die Strände sind voll, die Hotels sind ausgebucht und die Preise steigen. Viele Einwohner sind nicht zufrieden, weil das Leben auf der Insel immer teurer wird. Der Tourismus bringt viel Geld, aber auch Probleme: zu viel Müll, zu viele Autos und zu wenig Wasser. Die Regierung sucht jetzt neue Ideen für einen besseren Tourismus.}

\textbf{Points de vigilance :}
\begin{itemize}
\item \all{die Insel} = l'île (féminin, faux ami : pas *\all{der Insel}) ; \cle{ausgebucht} ou \cle{voll} = complet (pas *\all{reserviert} qui signifie „réservé“).
\item \all{die Regierung} = le gouvernement ; \all{zu wenig Wasser} = trop peu d'eau.
\item \all{nach Mallorca} (île = pays sans article) ; \all{auf der Insel} (position, \cle{Dativ}).
\end{itemize}
\end{corrige}

\begin{corrige}[Exercice 11 — Production écrite : \all{Mein Traumurlaub}]
\textbf{Barème indicatif (sur 20) :} respect du sujet et de la longueur (60--80 mots) : 4 pts ; richesse du vocabulaire de la leçon : 4 pts ; correction grammaticale (place du verbe, \cle{weil}/\cle{dass}, parfait \cle{sein}/\cle{haben}) : 8 pts ; orthographe (majuscules des noms, Umlaute, ß) : 4 pts.

\textbf{Proposition de rédaction modèle :}

\all{Ich mag den Ökotourismus sehr, weil er die Natur und die Kultur respektiert. Ich denke, dass dieser Tourismus die Zukunft ist. Letztes Jahr bin ich mit meiner Familie nach Kribi gefahren. Wir sind eine Woche in einem kleinen Hotel geblieben, und jeden Tag sind wir an den Strand gegangen. Es war wunderbar! Mein Traumurlaub ist eine Reise nach Österreich: ich möchte in den Alpen wandern und die Sehenswürdigkeiten von Wien besuchen.}

« J'aime beaucoup l'écotourisme, parce qu'il respecte la nature et la culture. Je pense que ce tourisme est l'avenir. L'année dernière, je suis allé avec ma famille à Kribi. Nous sommes restés une semaine dans un petit hôtel, et chaque jour nous sommes allés à la plage. C'était merveilleux ! Mes vacances de rêve sont un voyage en Autriche : je voudrais faire de la randonnée dans les Alpes et visiter les monuments de Vienne. » (75 mots)

\textbf{Points de vigilance :}
\begin{itemize}
\item Verbe conjugué en \textbf{deuxième position} dans la principale, à la \textbf{fin} après \cle{weil} et \cle{dass}.
\item Parfait : \cle{sein} pour \all{fahren, gehen, bleiben, reisen, kommen, fliegen} ; \cle{haben} pour \all{besuchen, machen, essen, sehen}.
\item Majuscule à \textbf{tous} les noms : \all{die Reise}, \all{der Strand}, \all{das Hotel}, \all{die Alpen}, \all{die Sehenswürdigkeiten}.
\item Prépositions de lieu : \all{nach Kribi} (pays/ville), \all{an den Strand} (mouvement, \cle{Akkusativ}), \all{am Strand} (position, \cle{Dativ}), \all{in den Alpen} (position/région, \cle{Dativ} pluriel en \textit{-n}).
\end{itemize}
\end{corrige}

\newpage

\coursec{Fiche bilan}

\begin{popfigure}
\begin{tikzpicture}[mindmap, grow cyclic, every node/.style={concept}, text width=2.6cm, align=center,
root concept/.append style={concept color=popblue, font=\bfseries, text=white, text width=3cm},
level 1/.append style={level distance=3.6cm, sibling angle=51.4, font=\small},
level 1 concept/.append style={text width=2.4cm}]
\node[root concept]{Tourismus : Segen oder Problem?}
child[concept color=poporange]{node[align=center]{Lexique\\der Strand, die Reise, das Hotel, die Sehenswürdigkeit}}
child[concept color=popgreen]{node[align=center]{Prépositions de lieu\\wo ? + Dativ\\wohin ? + Akkusativ}}
child[concept color=popgold]{node[align=center]{Subordonnées\\weil + verbe final\\dass + verbe final}}
child[concept color=popgreenL]{node[align=center]{Perfekt mit sein\\bin gefahren, ist geblieben}}
child[concept color=poppurple]{node[align=center]{Types de tourisme\\der Massentourismus, der Ökotourismus}}
child[concept color=poppink]{node[align=center]{Problèmes\\die Umwelt, die Preise, der Müll}}
child[concept color=popblueL]{node[align=center]{Production\\Mein Traumurlaub}};
\end{tikzpicture}
\legende{Carte mentale de la leçon : le tourisme, ses formes, ses problèmes et les outils de langue associés.}
\end{popfigure}

\begin{aretenir}[L'essentiel de la leçon]
\textbf{1. Lexique clé (à connaître avec article et pluriel)}
\begin{itemize}
\item \all{der Tourismus} (sans pluriel) : le tourisme ; \all{der Tourist, -en} : le touriste ; \all{die Touristin, -nen} : la touriste.
\item \all{die Reise, -n} : le voyage ; \all{reisen} (verbe régulier, avec \all{sein} au Perfekt : \all{ich bin gereist}) : voyager ; \all{der Urlaub, -e} : les vacances.
\item \all{das Hotel, -s} : l'hôtel ; \all{die Jugendherberge, -n} : l'auberge de jeunesse ; \all{die Übernachtung, -en} : la nuitée.
\item \all{der Strand, die Strände} : la plage ; \all{das Meer, -e} : la mer ; \all{die Küste, -n} : la côte.
\item \all{die Sehenswürdigkeit, -en} : le site touristique (la curiosité) ; \all{besichtigen} : visiter (un lieu) ; \all{besuchen} : visiter (une personne ou un lieu).
\item \all{der Massentourismus} : le tourisme de masse ; \all{der Ökotourismus} : l'écotourisme ; \all{der Kulturtourismus} : le tourisme culturel.
\item \all{die Umwelt} (sans pluriel) : l'environnement ; \all{der Müll} (sans pluriel) : les déchets ; \all{verschmutzen} : polluer ; \all{schützen} : protéger.
\end{itemize}

\textbf{2. Grammaire à connaître par cœur}
\begin{center}
\begin{tabularx}{\linewidth}{|X|X|X|}
\hline
\rowcolor{popblueL} Règle & Forme & Exemple \\
\hline
Prépositions de lieu à deux cas (Wechselpräpositionen) & \cle{wo ?} (position) + datif ; \cle{wohin ?} (mouvement) + accusatif & \all{Wir sind \textbf{am} Strand. / Wir fahren \textbf{an den} Strand.} « Nous sommes à la plage. / Nous allons à la plage. » \\
\hline
Prépositions concernées & an, auf, in, über, unter, vor, hinter, neben, zwischen & \all{Das Hotel liegt \textbf{in der} Stadt. / Wir gehen \textbf{in die} Stadt.} « L'hôtel se trouve dans la ville. / Nous allons en ville. » \\
\hline
Subordonnée avec \all{weil} & exprime la cause ; verbe conjugué en dernière position & \all{Wir bleiben hier, \textbf{weil} das Wetter schön \textbf{ist}.} « Nous restons ici parce que le temps est beau. » \\
\hline
Subordonnée avec \all{dass} & exprime une opinion ou un fait rapporté ; verbe conjugué en dernière position & \all{Ich denke, \textbf{dass} der Tourismus wichtig \textbf{ist}.} « Je pense que le tourisme est important. » \\
\hline
Perfekt avec \all{sein} & verbes de déplacement et de changement d'état : gehen, fahren, fliegen, kommen, reisen, bleiben, sein & \all{Wir \textbf{sind} nach Kribi \textbf{gefahren}.} « Nous sommes allés à Kribi. » \\
\hline
\end{tabularx}
\end{center}

\textbf{3. Méthodes express}
\begin{itemize}
\item \textbf{Comprendre le texte} : lire le titre, repérer les mots du champ lexical du tourisme, puis répondre aux questions en recopiant la structure de la question.
\item \textbf{Choisir le cas après une préposition de lieu} : se demander « mouvement ou position ? » ; mouvement (\cle{wohin ?}) = accusatif, position (\cle{wo ?}) = datif.
\item \textbf{Construire une subordonnée} : conjonction (\all{weil}/\all{dass}) + sujet + compléments + verbe conjugué en dernière position, sans oublier la virgule avant la conjonction.
\item \textbf{Rédiger « Mein Traumurlaub »} : destination (wohin ?), activités, moyen de transport au Perfekt, une phrase avec \all{weil} et une avec \all{dass}.
\end{itemize}
\end{aretenir}

\begin{attention}[Pièges fréquents des francophones]
\begin{itemize}
\item Ne pas confondre \all{das Meer} (la mer) et \all{das Mehr} (le surplus) ; et attention au faux ami : \all{die See} signifie aussi « la mer » dans le Nord de l'Allemagne.
\item \all{der Urlaub} désigne les vacances (congés), pas « l'auberge » : l'auberge se dit \all{die Herberge}.
\item Après \all{weil} et \all{dass}, le verbe conjugué va \textbf{à la fin} : on écrit \all{weil das Wetter schön ist} et jamais \all{weil das Wetter ist schön}.
\item Au Perfekt, les verbes de déplacement prennent \all{sein} : \all{ich bin geflogen} (j'ai pris l'avion), et non \all{ich habe geflogen} dans le sens « je me suis déplacé en avion ».
\item Contractions obligatoires : \all{an + dem = am}, \all{an + das = ans}, \all{in + dem = im}, \all{in + das = ins}.
\end{itemize}
\end{attention}

\begin{aretenir}[Checklist avant le Bac]
\begin{itemize}
\item[$\square$] Je connais l'article et le pluriel du vocabulaire du tourisme (\all{der Strand, die Strände} ; \all{die Sehenswürdigkeit, -en}).
\item[$\square$] Je distingue les types de tourisme : \all{der Massentourismus}, \all{der Ökotourismus}, \all{der Kulturtourismus}.
\item[$\square$] Je sais citer deux avantages et deux problèmes du tourisme (emplois, devises ; pollution, hausse des prix).
\item[$\square$] Je choisis correctement entre accusatif et datif après \all{in, an, auf, neben, vor, hinter, über, unter, zwischen}.
\item[$\square$] Je sais que \all{wo ?} appelle le datif et \all{wohin ?} l'accusatif.
\item[$\square$] Je place le verbe conjugué à la fin après \all{weil} et \all{dass}.
\item[$\square$] Je n'oublie pas la virgule avant la subordonnée.
\item[$\square$] Je conjugue le Perfekt avec \all{sein} pour les verbes de déplacement : \all{ich bin gefahren, wir sind geblieben}.
\item[$\square$] Je mets la majuscule à tous les noms allemands (\all{das Hotel, die Reise}).
\item[$\square$] Je sais rédiger un court paragraphe « Mein Traumurlaub » avec au moins une subordonnée.
\item[$\square$] Je relis ma production : place du verbe, cas, orthographe des Umlaute (ä, ö, ü).
\end{itemize}
\end{aretenir}

\findecours

\end{document}
